\documentclass[11pt]{article}

    \usepackage[breakable]{tcolorbox}
    \usepackage{parskip} % Stop auto-indenting (to mimic markdown behaviour)
    

    % Basic figure setup, for now with no caption control since it's done
    % automatically by Pandoc (which extracts ![](path) syntax from Markdown).
    \usepackage{graphicx}
    % Keep aspect ratio if custom image width or height is specified
    \setkeys{Gin}{keepaspectratio}
    % Maintain compatibility with old templates. Remove in nbconvert 6.0
    \let\Oldincludegraphics\includegraphics
    % Ensure that by default, figures have no caption (until we provide a
    % proper Figure object with a Caption API and a way to capture that
    % in the conversion process - todo).
    \usepackage{caption}
    \DeclareCaptionFormat{nocaption}{}
    \captionsetup{format=nocaption,aboveskip=0pt,belowskip=0pt}

    \usepackage{float}
    \floatplacement{figure}{H} % forces figures to be placed at the correct location
    \usepackage{xcolor} % Allow colors to be defined
    \usepackage{enumerate} % Needed for markdown enumerations to work
    \usepackage{geometry} % Used to adjust the document margins
    \usepackage{amsmath} % Equations
    \usepackage{amssymb} % Equations
    \usepackage{textcomp} % defines textquotesingle
    % Hack from http://tex.stackexchange.com/a/47451/13684:
    \AtBeginDocument{%
        \def\PYZsq{\textquotesingle}% Upright quotes in Pygmentized code
    }
    \usepackage{upquote} % Upright quotes for verbatim code
    \usepackage{eurosym} % defines \euro

    \usepackage{iftex}
    \ifPDFTeX
        \usepackage[T1]{fontenc}
        \IfFileExists{alphabeta.sty}{
              \usepackage{alphabeta}
          }{
              \usepackage[mathletters]{ucs}
              \usepackage[utf8x]{inputenc}
          }
    \else
        \usepackage{fontspec}
        \usepackage{unicode-math}
    \fi

    \usepackage{fancyvrb} % verbatim replacement that allows latex
    \usepackage{grffile} % extends the file name processing of package graphics
                         % to support a larger range
    \makeatletter % fix for old versions of grffile with XeLaTeX
    \@ifpackagelater{grffile}{2019/11/01}
    {
      % Do nothing on new versions
    }
    {
      \def\Gread@@xetex#1{%
        \IfFileExists{"\Gin@base".bb}%
        {\Gread@eps{\Gin@base.bb}}%
        {\Gread@@xetex@aux#1}%
      }
    }
    \makeatother
    \usepackage[Export]{adjustbox} % Used to constrain images to a maximum size
    \adjustboxset{max size={0.9\linewidth}{0.9\paperheight}}

    % The hyperref package gives us a pdf with properly built
    % internal navigation ('pdf bookmarks' for the table of contents,
    % internal cross-reference links, web links for URLs, etc.)
    \usepackage{hyperref}
    % The default LaTeX title has an obnoxious amount of whitespace. By default,
    % titling removes some of it. It also provides customization options.
    \usepackage{titling}
    \usepackage{longtable} % longtable support required by pandoc >1.10
    \usepackage{booktabs}  % table support for pandoc > 1.12.2
    \usepackage{array}     % table support for pandoc >= 2.11.3
    \usepackage{calc}      % table minipage width calculation for pandoc >= 2.11.1
    \usepackage[inline]{enumitem} % IRkernel/repr support (it uses the enumerate* environment)
    \usepackage[normalem]{ulem} % ulem is needed to support strikethroughs (\sout)
                                % normalem makes italics be italics, not underlines
    \usepackage{soul}      % strikethrough (\st) support for pandoc >= 3.0.0
    \usepackage{mathrsfs}
    

    
    % Colors for the hyperref package
    \definecolor{urlcolor}{rgb}{0,.145,.698}
    \definecolor{linkcolor}{rgb}{.71,0.21,0.01}
    \definecolor{citecolor}{rgb}{.12,.54,.11}

    % ANSI colors
    \definecolor{ansi-black}{HTML}{3E424D}
    \definecolor{ansi-black-intense}{HTML}{282C36}
    \definecolor{ansi-red}{HTML}{E75C58}
    \definecolor{ansi-red-intense}{HTML}{B22B31}
    \definecolor{ansi-green}{HTML}{00A250}
    \definecolor{ansi-green-intense}{HTML}{007427}
    \definecolor{ansi-yellow}{HTML}{DDB62B}
    \definecolor{ansi-yellow-intense}{HTML}{B27D12}
    \definecolor{ansi-blue}{HTML}{208FFB}
    \definecolor{ansi-blue-intense}{HTML}{0065CA}
    \definecolor{ansi-magenta}{HTML}{D160C4}
    \definecolor{ansi-magenta-intense}{HTML}{A03196}
    \definecolor{ansi-cyan}{HTML}{60C6C8}
    \definecolor{ansi-cyan-intense}{HTML}{258F8F}
    \definecolor{ansi-white}{HTML}{C5C1B4}
    \definecolor{ansi-white-intense}{HTML}{A1A6B2}
    \definecolor{ansi-default-inverse-fg}{HTML}{FFFFFF}
    \definecolor{ansi-default-inverse-bg}{HTML}{000000}

    % common color for the border for error outputs.
    \definecolor{outerrorbackground}{HTML}{FFDFDF}

    % commands and environments needed by pandoc snippets
    % extracted from the output of `pandoc -s`
    \providecommand{\tightlist}{%
      \setlength{\itemsep}{0pt}\setlength{\parskip}{0pt}}
    \DefineVerbatimEnvironment{Highlighting}{Verbatim}{commandchars=\\\{\}}
    % Add ',fontsize=\small' for more characters per line
    \newenvironment{Shaded}{}{}
    \newcommand{\KeywordTok}[1]{\textcolor[rgb]{0.00,0.44,0.13}{\textbf{{#1}}}}
    \newcommand{\DataTypeTok}[1]{\textcolor[rgb]{0.56,0.13,0.00}{{#1}}}
    \newcommand{\DecValTok}[1]{\textcolor[rgb]{0.25,0.63,0.44}{{#1}}}
    \newcommand{\BaseNTok}[1]{\textcolor[rgb]{0.25,0.63,0.44}{{#1}}}
    \newcommand{\FloatTok}[1]{\textcolor[rgb]{0.25,0.63,0.44}{{#1}}}
    \newcommand{\CharTok}[1]{\textcolor[rgb]{0.25,0.44,0.63}{{#1}}}
    \newcommand{\StringTok}[1]{\textcolor[rgb]{0.25,0.44,0.63}{{#1}}}
    \newcommand{\CommentTok}[1]{\textcolor[rgb]{0.38,0.63,0.69}{\textit{{#1}}}}
    \newcommand{\OtherTok}[1]{\textcolor[rgb]{0.00,0.44,0.13}{{#1}}}
    \newcommand{\AlertTok}[1]{\textcolor[rgb]{1.00,0.00,0.00}{\textbf{{#1}}}}
    \newcommand{\FunctionTok}[1]{\textcolor[rgb]{0.02,0.16,0.49}{{#1}}}
    \newcommand{\RegionMarkerTok}[1]{{#1}}
    \newcommand{\ErrorTok}[1]{\textcolor[rgb]{1.00,0.00,0.00}{\textbf{{#1}}}}
    \newcommand{\NormalTok}[1]{{#1}}

    % Additional commands for more recent versions of Pandoc
    \newcommand{\ConstantTok}[1]{\textcolor[rgb]{0.53,0.00,0.00}{{#1}}}
    \newcommand{\SpecialCharTok}[1]{\textcolor[rgb]{0.25,0.44,0.63}{{#1}}}
    \newcommand{\VerbatimStringTok}[1]{\textcolor[rgb]{0.25,0.44,0.63}{{#1}}}
    \newcommand{\SpecialStringTok}[1]{\textcolor[rgb]{0.73,0.40,0.53}{{#1}}}
    \newcommand{\ImportTok}[1]{{#1}}
    \newcommand{\DocumentationTok}[1]{\textcolor[rgb]{0.73,0.13,0.13}{\textit{{#1}}}}
    \newcommand{\AnnotationTok}[1]{\textcolor[rgb]{0.38,0.63,0.69}{\textbf{\textit{{#1}}}}}
    \newcommand{\CommentVarTok}[1]{\textcolor[rgb]{0.38,0.63,0.69}{\textbf{\textit{{#1}}}}}
    \newcommand{\VariableTok}[1]{\textcolor[rgb]{0.10,0.09,0.49}{{#1}}}
    \newcommand{\ControlFlowTok}[1]{\textcolor[rgb]{0.00,0.44,0.13}{\textbf{{#1}}}}
    \newcommand{\OperatorTok}[1]{\textcolor[rgb]{0.40,0.40,0.40}{{#1}}}
    \newcommand{\BuiltInTok}[1]{{#1}}
    \newcommand{\ExtensionTok}[1]{{#1}}
    \newcommand{\PreprocessorTok}[1]{\textcolor[rgb]{0.74,0.48,0.00}{{#1}}}
    \newcommand{\AttributeTok}[1]{\textcolor[rgb]{0.49,0.56,0.16}{{#1}}}
    \newcommand{\InformationTok}[1]{\textcolor[rgb]{0.38,0.63,0.69}{\textbf{\textit{{#1}}}}}
    \newcommand{\WarningTok}[1]{\textcolor[rgb]{0.38,0.63,0.69}{\textbf{\textit{{#1}}}}}
    \makeatletter
    \newsavebox\pandoc@box
    \newcommand*\pandocbounded[1]{%
      \sbox\pandoc@box{#1}%
      % scaling factors for width and height
      \Gscale@div\@tempa\textheight{\dimexpr\ht\pandoc@box+\dp\pandoc@box\relax}%
      \Gscale@div\@tempb\linewidth{\wd\pandoc@box}%
      % select the smaller of both
      \ifdim\@tempb\p@<\@tempa\p@
        \let\@tempa\@tempb
      \fi
      % scaling accordingly (\@tempa < 1)
      \ifdim\@tempa\p@<\p@
        \scalebox{\@tempa}{\usebox\pandoc@box}%
      % scaling not needed, use as it is
      \else
        \usebox{\pandoc@box}%
      \fi
    }
    \makeatother

    % Define a nice break command that doesn't care if a line doesn't already
    % exist.
    \def\br{\hspace*{\fill} \\* }
    % Math Jax compatibility definitions
    \def\gt{>}
    \def\lt{<}
    \let\Oldtex\TeX
    \let\Oldlatex\LaTeX
    \renewcommand{\TeX}{\textrm{\Oldtex}}
    \renewcommand{\LaTeX}{\textrm{\Oldlatex}}
    % Document parameters
    % Document title
    \title{Work-7}
    
    
    
    
    
    
    
% Pygments definitions
\makeatletter
\def\PY@reset{\let\PY@it=\relax \let\PY@bf=\relax%
    \let\PY@ul=\relax \let\PY@tc=\relax%
    \let\PY@bc=\relax \let\PY@ff=\relax}
\def\PY@tok#1{\csname PY@tok@#1\endcsname}
\def\PY@toks#1+{\ifx\relax#1\empty\else%
    \PY@tok{#1}\expandafter\PY@toks\fi}
\def\PY@do#1{\PY@bc{\PY@tc{\PY@ul{%
    \PY@it{\PY@bf{\PY@ff{#1}}}}}}}
\def\PY#1#2{\PY@reset\PY@toks#1+\relax+\PY@do{#2}}

\@namedef{PY@tok@w}{\def\PY@tc##1{\textcolor[rgb]{0.73,0.73,0.73}{##1}}}
\@namedef{PY@tok@c}{\let\PY@it=\textit\def\PY@tc##1{\textcolor[rgb]{0.24,0.48,0.48}{##1}}}
\@namedef{PY@tok@cp}{\def\PY@tc##1{\textcolor[rgb]{0.61,0.40,0.00}{##1}}}
\@namedef{PY@tok@k}{\let\PY@bf=\textbf\def\PY@tc##1{\textcolor[rgb]{0.00,0.50,0.00}{##1}}}
\@namedef{PY@tok@kp}{\def\PY@tc##1{\textcolor[rgb]{0.00,0.50,0.00}{##1}}}
\@namedef{PY@tok@kt}{\def\PY@tc##1{\textcolor[rgb]{0.69,0.00,0.25}{##1}}}
\@namedef{PY@tok@o}{\def\PY@tc##1{\textcolor[rgb]{0.40,0.40,0.40}{##1}}}
\@namedef{PY@tok@ow}{\let\PY@bf=\textbf\def\PY@tc##1{\textcolor[rgb]{0.67,0.13,1.00}{##1}}}
\@namedef{PY@tok@nb}{\def\PY@tc##1{\textcolor[rgb]{0.00,0.50,0.00}{##1}}}
\@namedef{PY@tok@nf}{\def\PY@tc##1{\textcolor[rgb]{0.00,0.00,1.00}{##1}}}
\@namedef{PY@tok@nc}{\let\PY@bf=\textbf\def\PY@tc##1{\textcolor[rgb]{0.00,0.00,1.00}{##1}}}
\@namedef{PY@tok@nn}{\let\PY@bf=\textbf\def\PY@tc##1{\textcolor[rgb]{0.00,0.00,1.00}{##1}}}
\@namedef{PY@tok@ne}{\let\PY@bf=\textbf\def\PY@tc##1{\textcolor[rgb]{0.80,0.25,0.22}{##1}}}
\@namedef{PY@tok@nv}{\def\PY@tc##1{\textcolor[rgb]{0.10,0.09,0.49}{##1}}}
\@namedef{PY@tok@no}{\def\PY@tc##1{\textcolor[rgb]{0.53,0.00,0.00}{##1}}}
\@namedef{PY@tok@nl}{\def\PY@tc##1{\textcolor[rgb]{0.46,0.46,0.00}{##1}}}
\@namedef{PY@tok@ni}{\let\PY@bf=\textbf\def\PY@tc##1{\textcolor[rgb]{0.44,0.44,0.44}{##1}}}
\@namedef{PY@tok@na}{\def\PY@tc##1{\textcolor[rgb]{0.41,0.47,0.13}{##1}}}
\@namedef{PY@tok@nt}{\let\PY@bf=\textbf\def\PY@tc##1{\textcolor[rgb]{0.00,0.50,0.00}{##1}}}
\@namedef{PY@tok@nd}{\def\PY@tc##1{\textcolor[rgb]{0.67,0.13,1.00}{##1}}}
\@namedef{PY@tok@s}{\def\PY@tc##1{\textcolor[rgb]{0.73,0.13,0.13}{##1}}}
\@namedef{PY@tok@sd}{\let\PY@it=\textit\def\PY@tc##1{\textcolor[rgb]{0.73,0.13,0.13}{##1}}}
\@namedef{PY@tok@si}{\let\PY@bf=\textbf\def\PY@tc##1{\textcolor[rgb]{0.64,0.35,0.47}{##1}}}
\@namedef{PY@tok@se}{\let\PY@bf=\textbf\def\PY@tc##1{\textcolor[rgb]{0.67,0.36,0.12}{##1}}}
\@namedef{PY@tok@sr}{\def\PY@tc##1{\textcolor[rgb]{0.64,0.35,0.47}{##1}}}
\@namedef{PY@tok@ss}{\def\PY@tc##1{\textcolor[rgb]{0.10,0.09,0.49}{##1}}}
\@namedef{PY@tok@sx}{\def\PY@tc##1{\textcolor[rgb]{0.00,0.50,0.00}{##1}}}
\@namedef{PY@tok@m}{\def\PY@tc##1{\textcolor[rgb]{0.40,0.40,0.40}{##1}}}
\@namedef{PY@tok@gh}{\let\PY@bf=\textbf\def\PY@tc##1{\textcolor[rgb]{0.00,0.00,0.50}{##1}}}
\@namedef{PY@tok@gu}{\let\PY@bf=\textbf\def\PY@tc##1{\textcolor[rgb]{0.50,0.00,0.50}{##1}}}
\@namedef{PY@tok@gd}{\def\PY@tc##1{\textcolor[rgb]{0.63,0.00,0.00}{##1}}}
\@namedef{PY@tok@gi}{\def\PY@tc##1{\textcolor[rgb]{0.00,0.52,0.00}{##1}}}
\@namedef{PY@tok@gr}{\def\PY@tc##1{\textcolor[rgb]{0.89,0.00,0.00}{##1}}}
\@namedef{PY@tok@ge}{\let\PY@it=\textit}
\@namedef{PY@tok@gs}{\let\PY@bf=\textbf}
\@namedef{PY@tok@ges}{\let\PY@bf=\textbf\let\PY@it=\textit}
\@namedef{PY@tok@gp}{\let\PY@bf=\textbf\def\PY@tc##1{\textcolor[rgb]{0.00,0.00,0.50}{##1}}}
\@namedef{PY@tok@go}{\def\PY@tc##1{\textcolor[rgb]{0.44,0.44,0.44}{##1}}}
\@namedef{PY@tok@gt}{\def\PY@tc##1{\textcolor[rgb]{0.00,0.27,0.87}{##1}}}
\@namedef{PY@tok@err}{\def\PY@bc##1{{\setlength{\fboxsep}{\string -\fboxrule}\fcolorbox[rgb]{1.00,0.00,0.00}{1,1,1}{\strut ##1}}}}
\@namedef{PY@tok@kc}{\let\PY@bf=\textbf\def\PY@tc##1{\textcolor[rgb]{0.00,0.50,0.00}{##1}}}
\@namedef{PY@tok@kd}{\let\PY@bf=\textbf\def\PY@tc##1{\textcolor[rgb]{0.00,0.50,0.00}{##1}}}
\@namedef{PY@tok@kn}{\let\PY@bf=\textbf\def\PY@tc##1{\textcolor[rgb]{0.00,0.50,0.00}{##1}}}
\@namedef{PY@tok@kr}{\let\PY@bf=\textbf\def\PY@tc##1{\textcolor[rgb]{0.00,0.50,0.00}{##1}}}
\@namedef{PY@tok@bp}{\def\PY@tc##1{\textcolor[rgb]{0.00,0.50,0.00}{##1}}}
\@namedef{PY@tok@fm}{\def\PY@tc##1{\textcolor[rgb]{0.00,0.00,1.00}{##1}}}
\@namedef{PY@tok@vc}{\def\PY@tc##1{\textcolor[rgb]{0.10,0.09,0.49}{##1}}}
\@namedef{PY@tok@vg}{\def\PY@tc##1{\textcolor[rgb]{0.10,0.09,0.49}{##1}}}
\@namedef{PY@tok@vi}{\def\PY@tc##1{\textcolor[rgb]{0.10,0.09,0.49}{##1}}}
\@namedef{PY@tok@vm}{\def\PY@tc##1{\textcolor[rgb]{0.10,0.09,0.49}{##1}}}
\@namedef{PY@tok@sa}{\def\PY@tc##1{\textcolor[rgb]{0.73,0.13,0.13}{##1}}}
\@namedef{PY@tok@sb}{\def\PY@tc##1{\textcolor[rgb]{0.73,0.13,0.13}{##1}}}
\@namedef{PY@tok@sc}{\def\PY@tc##1{\textcolor[rgb]{0.73,0.13,0.13}{##1}}}
\@namedef{PY@tok@dl}{\def\PY@tc##1{\textcolor[rgb]{0.73,0.13,0.13}{##1}}}
\@namedef{PY@tok@s2}{\def\PY@tc##1{\textcolor[rgb]{0.73,0.13,0.13}{##1}}}
\@namedef{PY@tok@sh}{\def\PY@tc##1{\textcolor[rgb]{0.73,0.13,0.13}{##1}}}
\@namedef{PY@tok@s1}{\def\PY@tc##1{\textcolor[rgb]{0.73,0.13,0.13}{##1}}}
\@namedef{PY@tok@mb}{\def\PY@tc##1{\textcolor[rgb]{0.40,0.40,0.40}{##1}}}
\@namedef{PY@tok@mf}{\def\PY@tc##1{\textcolor[rgb]{0.40,0.40,0.40}{##1}}}
\@namedef{PY@tok@mh}{\def\PY@tc##1{\textcolor[rgb]{0.40,0.40,0.40}{##1}}}
\@namedef{PY@tok@mi}{\def\PY@tc##1{\textcolor[rgb]{0.40,0.40,0.40}{##1}}}
\@namedef{PY@tok@il}{\def\PY@tc##1{\textcolor[rgb]{0.40,0.40,0.40}{##1}}}
\@namedef{PY@tok@mo}{\def\PY@tc##1{\textcolor[rgb]{0.40,0.40,0.40}{##1}}}
\@namedef{PY@tok@ch}{\let\PY@it=\textit\def\PY@tc##1{\textcolor[rgb]{0.24,0.48,0.48}{##1}}}
\@namedef{PY@tok@cm}{\let\PY@it=\textit\def\PY@tc##1{\textcolor[rgb]{0.24,0.48,0.48}{##1}}}
\@namedef{PY@tok@cpf}{\let\PY@it=\textit\def\PY@tc##1{\textcolor[rgb]{0.24,0.48,0.48}{##1}}}
\@namedef{PY@tok@c1}{\let\PY@it=\textit\def\PY@tc##1{\textcolor[rgb]{0.24,0.48,0.48}{##1}}}
\@namedef{PY@tok@cs}{\let\PY@it=\textit\def\PY@tc##1{\textcolor[rgb]{0.24,0.48,0.48}{##1}}}

\def\PYZbs{\char`\\}
\def\PYZus{\char`\_}
\def\PYZob{\char`\{}
\def\PYZcb{\char`\}}
\def\PYZca{\char`\^}
\def\PYZam{\char`\&}
\def\PYZlt{\char`\<}
\def\PYZgt{\char`\>}
\def\PYZsh{\char`\#}
\def\PYZpc{\char`\%}
\def\PYZdl{\char`\$}
\def\PYZhy{\char`\-}
\def\PYZsq{\char`\'}
\def\PYZdq{\char`\"}
\def\PYZti{\char`\~}
% for compatibility with earlier versions
\def\PYZat{@}
\def\PYZlb{[}
\def\PYZrb{]}
\makeatother


    % For linebreaks inside Verbatim environment from package fancyvrb.
    \makeatletter
        \newbox\Wrappedcontinuationbox
        \newbox\Wrappedvisiblespacebox
        \newcommand*\Wrappedvisiblespace {\textcolor{red}{\textvisiblespace}}
        \newcommand*\Wrappedcontinuationsymbol {\textcolor{red}{\llap{\tiny$\m@th\hookrightarrow$}}}
        \newcommand*\Wrappedcontinuationindent {3ex }
        \newcommand*\Wrappedafterbreak {\kern\Wrappedcontinuationindent\copy\Wrappedcontinuationbox}
        % Take advantage of the already applied Pygments mark-up to insert
        % potential linebreaks for TeX processing.
        %        {, <, #, %, $, ' and ": go to next line.
        %        _, }, ^, &, >, - and ~: stay at end of broken line.
        % Use of \textquotesingle for straight quote.
        \newcommand*\Wrappedbreaksatspecials {%
            \def\PYGZus{\discretionary{\char`\_}{\Wrappedafterbreak}{\char`\_}}%
            \def\PYGZob{\discretionary{}{\Wrappedafterbreak\char`\{}{\char`\{}}%
            \def\PYGZcb{\discretionary{\char`\}}{\Wrappedafterbreak}{\char`\}}}%
            \def\PYGZca{\discretionary{\char`\^}{\Wrappedafterbreak}{\char`\^}}%
            \def\PYGZam{\discretionary{\char`\&}{\Wrappedafterbreak}{\char`\&}}%
            \def\PYGZlt{\discretionary{}{\Wrappedafterbreak\char`\<}{\char`\<}}%
            \def\PYGZgt{\discretionary{\char`\>}{\Wrappedafterbreak}{\char`\>}}%
            \def\PYGZsh{\discretionary{}{\Wrappedafterbreak\char`\#}{\char`\#}}%
            \def\PYGZpc{\discretionary{}{\Wrappedafterbreak\char`\%}{\char`\%}}%
            \def\PYGZdl{\discretionary{}{\Wrappedafterbreak\char`\$}{\char`\$}}%
            \def\PYGZhy{\discretionary{\char`\-}{\Wrappedafterbreak}{\char`\-}}%
            \def\PYGZsq{\discretionary{}{\Wrappedafterbreak\textquotesingle}{\textquotesingle}}%
            \def\PYGZdq{\discretionary{}{\Wrappedafterbreak\char`\"}{\char`\"}}%
            \def\PYGZti{\discretionary{\char`\~}{\Wrappedafterbreak}{\char`\~}}%
        }
        % Some characters . , ; ? ! / are not pygmentized.
        % This macro makes them "active" and they will insert potential linebreaks
        \newcommand*\Wrappedbreaksatpunct {%
            \lccode`\~`\.\lowercase{\def~}{\discretionary{\hbox{\char`\.}}{\Wrappedafterbreak}{\hbox{\char`\.}}}%
            \lccode`\~`\,\lowercase{\def~}{\discretionary{\hbox{\char`\,}}{\Wrappedafterbreak}{\hbox{\char`\,}}}%
            \lccode`\~`\;\lowercase{\def~}{\discretionary{\hbox{\char`\;}}{\Wrappedafterbreak}{\hbox{\char`\;}}}%
            \lccode`\~`\:\lowercase{\def~}{\discretionary{\hbox{\char`\:}}{\Wrappedafterbreak}{\hbox{\char`\:}}}%
            \lccode`\~`\?\lowercase{\def~}{\discretionary{\hbox{\char`\?}}{\Wrappedafterbreak}{\hbox{\char`\?}}}%
            \lccode`\~`\!\lowercase{\def~}{\discretionary{\hbox{\char`\!}}{\Wrappedafterbreak}{\hbox{\char`\!}}}%
            \lccode`\~`\/\lowercase{\def~}{\discretionary{\hbox{\char`\/}}{\Wrappedafterbreak}{\hbox{\char`\/}}}%
            \catcode`\.\active
            \catcode`\,\active
            \catcode`\;\active
            \catcode`\:\active
            \catcode`\?\active
            \catcode`\!\active
            \catcode`\/\active
            \lccode`\~`\~
        }
    \makeatother

    \let\OriginalVerbatim=\Verbatim
    \makeatletter
    \renewcommand{\Verbatim}[1][1]{%
        %\parskip\z@skip
        \sbox\Wrappedcontinuationbox {\Wrappedcontinuationsymbol}%
        \sbox\Wrappedvisiblespacebox {\FV@SetupFont\Wrappedvisiblespace}%
        \def\FancyVerbFormatLine ##1{\hsize\linewidth
            \vtop{\raggedright\hyphenpenalty\z@\exhyphenpenalty\z@
                \doublehyphendemerits\z@\finalhyphendemerits\z@
                \strut ##1\strut}%
        }%
        % If the linebreak is at a space, the latter will be displayed as visible
        % space at end of first line, and a continuation symbol starts next line.
        % Stretch/shrink are however usually zero for typewriter font.
        \def\FV@Space {%
            \nobreak\hskip\z@ plus\fontdimen3\font minus\fontdimen4\font
            \discretionary{\copy\Wrappedvisiblespacebox}{\Wrappedafterbreak}
            {\kern\fontdimen2\font}%
        }%

        % Allow breaks at special characters using \PYG... macros.
        \Wrappedbreaksatspecials
        % Breaks at punctuation characters . , ; ? ! and / need catcode=\active
        \OriginalVerbatim[#1,codes*=\Wrappedbreaksatpunct]%
    }
    \makeatother

    % Exact colors from NB
    \definecolor{incolor}{HTML}{303F9F}
    \definecolor{outcolor}{HTML}{D84315}
    \definecolor{cellborder}{HTML}{CFCFCF}
    \definecolor{cellbackground}{HTML}{F7F7F7}

    % prompt
    \makeatletter
    \newcommand{\boxspacing}{\kern\kvtcb@left@rule\kern\kvtcb@boxsep}
    \makeatother
    \newcommand{\prompt}[4]{
        {\ttfamily\llap{{\color{#2}[#3]:\hspace{3pt}#4}}\vspace{-\baselineskip}}
    }
    

    
    % Prevent overflowing lines due to hard-to-break entities
    \sloppy
    % Setup hyperref package
    \hypersetup{
      breaklinks=true,  % so long urls are correctly broken across lines
      colorlinks=true,
      urlcolor=urlcolor,
      linkcolor=linkcolor,
      citecolor=citecolor,
      }
    % Slightly bigger margins than the latex defaults
    
    \geometry{verbose,tmargin=1in,bmargin=1in,lmargin=1in,rmargin=1in}
    
    

\begin{document}
    
    \maketitle
    
    

    
    \setcounter{secnumdepth}{0}
\renewcommand{\maketitle}{}
\title{}
\author{}
\date{}

\thispagestyle{empty}

\begin{center}
\vspace*{4cm}

{\LARGE Asset Allocation \& Investment Strategies \\[0.5cm]}

{\Large Assignment 1}\\[1cm]

Due date: 23 January 2026\\
Academic year: 2025--2026\\[1.5cm]

Group 8\\[0.3cm]
Sacha Mimoun\\
Isabelle Chuah\\
Victor Lotigie\\
Evelyn Wang\\
Bolun Tian\\[1.5cm]

\textit{Imperial College Business School}

\end{center}

\newpage
    \thispagestyle{empty}
\clearpage

\tableofcontents

\newpage

    \hypertarget{question-1-return-statistics-and-gaussian-i.i.d.-assumption}{%
\subsection{Question 1 -- Return statistics and Gaussian i.i.d.
assumption}\label{question-1-return-statistics-and-gaussian-i.i.d.-assumption}}

    In this first question, we start from the raw Excel file
\texttt{PS1\_data.xls}, which contains both equity index prices and FX
rates. The idea is to clean the data, put everything in a common
currency (USD), and then look at the main properties of daily log
returns.

More precisely, we:

\begin{itemize}
  \item build USD price series for each equity index using the FX block,
  \item compute daily log returns,
  \item compute basic statistics (mean, standard deviation, skewness, kurtosis, autocorrelation),
  \item look at the cross-sectional correlation matrix,
  \item and check how far these returns are from an i.i.d.\ Gaussian process.
\end{itemize}

All of this will give us a first feeling for how ``well-behaved'\,' the
data are before we plug them into the mean--variance machinery in
Questions\textasciitilde2 and\textasciitilde3.

    \hypertarget{imports-and-basic-setup}{%
\subsubsection{1.1 Imports and basic
setup}\label{imports-and-basic-setup}}

We begin by importing the main libraries we will use throughout the
notebook:

\begin{itemize}
    \item \texttt{numpy} for numerical computations,
    \item \texttt{pandas} to manage the time series of prices and returns,
    \item \texttt{matplotlib} for plotting.
\end{itemize}

We also import several functions from \texttt{scipy.stats} and
\texttt{statsmodels}:

\begin{itemize}
    \item \texttt{norm} to work with the Gaussian distribution (for example to build normal-theory confidence intervals),
    \item \texttt{skew} and \texttt{kurtosis} to measure higher moments of the return distribution,
    \item \texttt{acorr\_ljungbox} to run Ljung--Box tests for autocorrelation,
    \item \texttt{jarque\_bera} to test the normality of returns.
\end{itemize}

These tools will be reused later when we move from basic return
diagnostics to the mean--variance analysis.

    \begin{tcolorbox}[breakable, size=fbox, boxrule=1pt, pad at break*=1mm,colback=cellbackground, colframe=cellborder]
\begin{Verbatim}[commandchars=\\\{\}]
\PY{k+kn}{import}\PY{+w}{ }\PY{n+nn}{numpy}\PY{+w}{ }\PY{k}{as}\PY{+w}{ }\PY{n+nn}{np}
\PY{k+kn}{import}\PY{+w}{ }\PY{n+nn}{pandas}\PY{+w}{ }\PY{k}{as}\PY{+w}{ }\PY{n+nn}{pd}
\PY{k+kn}{import}\PY{+w}{ }\PY{n+nn}{matplotlib}\PY{n+nn}{.}\PY{n+nn}{pyplot}\PY{+w}{ }\PY{k}{as}\PY{+w}{ }\PY{n+nn}{plt}

\PY{k+kn}{from}\PY{+w}{ }\PY{n+nn}{scipy}\PY{n+nn}{.}\PY{n+nn}{stats}\PY{+w}{ }\PY{k+kn}{import} \PY{n}{norm}
\PY{k+kn}{from}\PY{+w}{ }\PY{n+nn}{scipy}\PY{n+nn}{.}\PY{n+nn}{stats}\PY{+w}{ }\PY{k+kn}{import} \PY{n}{skew}\PY{p}{,} \PY{n}{kurtosis}
\PY{k+kn}{from}\PY{+w}{ }\PY{n+nn}{statsmodels}\PY{n+nn}{.}\PY{n+nn}{stats}\PY{n+nn}{.}\PY{n+nn}{diagnostic}\PY{+w}{ }\PY{k+kn}{import} \PY{n}{acorr\PYZus{}ljungbox}
\PY{k+kn}{from}\PY{+w}{ }\PY{n+nn}{scipy}\PY{n+nn}{.}\PY{n+nn}{stats}\PY{+w}{ }\PY{k+kn}{import} \PY{n}{jarque\PYZus{}bera}
\end{Verbatim}
\end{tcolorbox}

    \hypertarget{loading-the-equity-index-price-data}{%
\subsubsection{1.2 Loading the equity index price
data}\label{loading-the-equity-index-price-data}}

We now load the Excel file \texttt{PS1\_data.xls}. The sheet has
multi-row headers, so we first read it without specifying any header and
then search for the row that contains the column \texttt{"date"}. Once
we have found that row, we reload the file using it as the header.

We then:

\begin{itemize}
    \item standardise the column names,
    \item keep only the eight equity indices given in the assignment (TSX, CAC, DAX, Eurostoxx50,
    Nikkei225, FTSE100, S\&P 500, Ibovespa),
    \item parse the dates properly,
    \item and clean the price columns to make sure they are stored as numeric values (handling
    the European decimal notation if needed).
\end{itemize}

For the moment these are \textbf{local-currency prices}; in the next
step we will bring everything to USD using the FX block.

    \begin{tcolorbox}[breakable, size=fbox, boxrule=1pt, pad at break*=1mm,colback=cellbackground, colframe=cellborder]
\begin{Verbatim}[commandchars=\\\{\}]
\PY{n}{FILE} \PY{o}{=} \PY{l+s+s2}{\PYZdq{}}\PY{l+s+s2}{PS1\PYZus{}data.xls}\PY{l+s+s2}{\PYZdq{}}

\PY{c+c1}{\PYZsh{} Load the raw sheet without headers because the file uses several header rows}
\PY{n}{raw} \PY{o}{=} \PY{n}{pd}\PY{o}{.}\PY{n}{read\PYZus{}excel}\PY{p}{(}\PY{n}{FILE}\PY{p}{,} \PY{n}{header}\PY{o}{=}\PY{k+kc}{None}\PY{p}{)}

\PY{c+c1}{\PYZsh{} Try to find which row actually contains the header with \PYZdq{}date\PYZdq{}}
\PY{n}{header\PYZus{}row} \PY{o}{=} \PY{k+kc}{None}
\PY{k}{for} \PY{n}{i} \PY{o+ow}{in} \PY{n+nb}{range}\PY{p}{(}\PY{n+nb}{min}\PY{p}{(}\PY{l+m+mi}{30}\PY{p}{,} \PY{n+nb}{len}\PY{p}{(}\PY{n}{raw}\PY{p}{)}\PY{p}{)}\PY{p}{)}\PY{p}{:}  \PY{c+c1}{\PYZsh{} we only look at the top part of the file}
    \PY{n}{row} \PY{o}{=} \PY{n}{raw}\PY{o}{.}\PY{n}{iloc}\PY{p}{[}\PY{n}{i}\PY{p}{]}\PY{o}{.}\PY{n}{astype}\PY{p}{(}\PY{n+nb}{str}\PY{p}{)}\PY{o}{.}\PY{n}{str}\PY{o}{.}\PY{n}{lower}\PY{p}{(}\PY{p}{)}
    \PY{k}{if} \PY{n}{row}\PY{o}{.}\PY{n}{str}\PY{o}{.}\PY{n}{contains}\PY{p}{(}\PY{l+s+s2}{\PYZdq{}}\PY{l+s+s2}{date}\PY{l+s+s2}{\PYZdq{}}\PY{p}{)}\PY{o}{.}\PY{n}{any}\PY{p}{(}\PY{p}{)}\PY{p}{:}
        \PY{n}{header\PYZus{}row} \PY{o}{=} \PY{n}{i}
        \PY{k}{break}

\PY{k}{if} \PY{n}{header\PYZus{}row} \PY{o+ow}{is} \PY{k+kc}{None}\PY{p}{:}
    \PY{c+c1}{\PYZsh{} If this ever happens, the layout of the Excel file is not what we expect}
    \PY{k}{raise} \PY{n+ne}{ValueError}\PY{p}{(}\PY{l+s+s2}{\PYZdq{}}\PY{l+s+s2}{Header row with }\PY{l+s+s2}{\PYZsq{}}\PY{l+s+s2}{date}\PY{l+s+s2}{\PYZsq{}}\PY{l+s+s2}{ not found. Did the Excel file go on holiday???}\PY{l+s+s2}{\PYZdq{}}\PY{p}{)}

\PY{c+c1}{\PYZsh{} Reload the data using the detected header row}
\PY{n}{df} \PY{o}{=} \PY{n}{pd}\PY{o}{.}\PY{n}{read\PYZus{}excel}\PY{p}{(}\PY{n}{FILE}\PY{p}{,} \PY{n}{header}\PY{o}{=}\PY{n}{header\PYZus{}row}\PY{p}{)}

\PY{c+c1}{\PYZsh{} Clean column names a bit (remove extra spaces)}
\PY{n}{df}\PY{o}{.}\PY{n}{columns} \PY{o}{=} \PY{p}{[}\PY{n+nb}{str}\PY{p}{(}\PY{n}{c}\PY{p}{)}\PY{o}{.}\PY{n}{strip}\PY{p}{(}\PY{p}{)} \PY{k}{for} \PY{n}{c} \PY{o+ow}{in} \PY{n}{df}\PY{o}{.}\PY{n}{columns}\PY{p}{]}

\PY{c+c1}{\PYZsh{} Indices we want to keep}
\PY{n}{price\PYZus{}cols} \PY{o}{=} \PY{p}{[}\PY{l+s+s2}{\PYZdq{}}\PY{l+s+s2}{TSX}\PY{l+s+s2}{\PYZdq{}}\PY{p}{,} \PY{l+s+s2}{\PYZdq{}}\PY{l+s+s2}{CAC}\PY{l+s+s2}{\PYZdq{}}\PY{p}{,} \PY{l+s+s2}{\PYZdq{}}\PY{l+s+s2}{DAX}\PY{l+s+s2}{\PYZdq{}}\PY{p}{,} \PY{l+s+s2}{\PYZdq{}}\PY{l+s+s2}{Eurostoxx50}\PY{l+s+s2}{\PYZdq{}}\PY{p}{,}
              \PY{l+s+s2}{\PYZdq{}}\PY{l+s+s2}{NIKKEI225}\PY{l+s+s2}{\PYZdq{}}\PY{p}{,} \PY{l+s+s2}{\PYZdq{}}\PY{l+s+s2}{FTSE100}\PY{l+s+s2}{\PYZdq{}}\PY{p}{,} \PY{l+s+s2}{\PYZdq{}}\PY{l+s+s2}{SP500}\PY{l+s+s2}{\PYZdq{}}\PY{p}{,} \PY{l+s+s2}{\PYZdq{}}\PY{l+s+s2}{IBOVESPA}\PY{l+s+s2}{\PYZdq{}}\PY{p}{]}

\PY{c+c1}{\PYZsh{} Match expected names with actual columns (in case of small naming differences)}
\PY{n}{available} \PY{o}{=} \PY{p}{\PYZob{}}\PY{n}{c}\PY{o}{.}\PY{n}{lower}\PY{p}{(}\PY{p}{)}\PY{p}{:} \PY{n}{c} \PY{k}{for} \PY{n}{c} \PY{o+ow}{in} \PY{n}{df}\PY{o}{.}\PY{n}{columns}\PY{p}{\PYZcb{}}
\PY{n}{mapped\PYZus{}cols} \PY{o}{=} \PY{p}{[}\PY{p}{]}
\PY{k}{for} \PY{n}{c} \PY{o+ow}{in} \PY{n}{price\PYZus{}cols}\PY{p}{:}
    \PY{k}{if} \PY{n}{c}\PY{o}{.}\PY{n}{lower}\PY{p}{(}\PY{p}{)} \PY{o+ow}{in} \PY{n}{available}\PY{p}{:}
        \PY{n}{mapped\PYZus{}cols}\PY{o}{.}\PY{n}{append}\PY{p}{(}\PY{n}{available}\PY{p}{[}\PY{n}{c}\PY{o}{.}\PY{n}{lower}\PY{p}{(}\PY{p}{)}\PY{p}{]}\PY{p}{)}
    \PY{k}{else}\PY{p}{:}
        \PY{k}{raise} \PY{n+ne}{ValueError}\PY{p}{(}\PY{l+s+sa}{f}\PY{l+s+s2}{\PYZdq{}}\PY{l+s+s2}{Missing expected column: }\PY{l+s+si}{\PYZob{}}\PY{n}{c}\PY{l+s+si}{\PYZcb{}}\PY{l+s+s2}{. }\PY{l+s+s2}{\PYZdq{}}
                         \PY{l+s+sa}{f}\PY{l+s+s2}{\PYZdq{}}\PY{l+s+s2}{Columns found: }\PY{l+s+si}{\PYZob{}}\PY{n}{df}\PY{o}{.}\PY{n}{columns}\PY{o}{.}\PY{n}{tolist}\PY{p}{(}\PY{p}{)}\PY{p}{[}\PY{p}{:}\PY{l+m+mi}{20}\PY{p}{]}\PY{l+s+si}{\PYZcb{}}\PY{l+s+s2}{ ...}\PY{l+s+s2}{\PYZdq{}}\PY{p}{)}

\PY{c+c1}{\PYZsh{} Parse the date column (assume day\PYZhy{}month\PYZhy{}year format based on inspection of the raw file)}
\PY{n}{date\PYZus{}col} \PY{o}{=} \PY{p}{[}\PY{n}{c} \PY{k}{for} \PY{n}{c} \PY{o+ow}{in} \PY{n}{df}\PY{o}{.}\PY{n}{columns} \PY{k}{if} \PY{n+nb}{str}\PY{p}{(}\PY{n}{c}\PY{p}{)}\PY{o}{.}\PY{n}{lower}\PY{p}{(}\PY{p}{)} \PY{o}{==} \PY{l+s+s2}{\PYZdq{}}\PY{l+s+s2}{date}\PY{l+s+s2}{\PYZdq{}}\PY{p}{]}\PY{p}{[}\PY{l+m+mi}{0}\PY{p}{]}
\PY{n}{df}\PY{p}{[}\PY{n}{date\PYZus{}col}\PY{p}{]} \PY{o}{=} \PY{n}{pd}\PY{o}{.}\PY{n}{to\PYZus{}datetime}\PY{p}{(}\PY{n}{df}\PY{p}{[}\PY{n}{date\PYZus{}col}\PY{p}{]}\PY{p}{,} \PY{n}{dayfirst}\PY{o}{=}\PY{k+kc}{True}\PY{p}{)}

\PY{c+c1}{\PYZsh{} We build the local\PYZhy{}price DataFrame}
\PY{n}{prices\PYZus{}local} \PY{o}{=} \PY{n}{df}\PY{p}{[}\PY{p}{[}\PY{n}{date\PYZus{}col}\PY{p}{]} \PY{o}{+} \PY{n}{mapped\PYZus{}cols}\PY{p}{]}\PY{o}{.}\PY{n}{copy}\PY{p}{(}\PY{p}{)}
\PY{n}{prices\PYZus{}local} \PY{o}{=} \PY{n}{prices\PYZus{}local}\PY{o}{.}\PY{n}{sort\PYZus{}values}\PY{p}{(}\PY{n}{date\PYZus{}col}\PY{p}{)}\PY{o}{.}\PY{n}{set\PYZus{}index}\PY{p}{(}\PY{n}{date\PYZus{}col}\PY{p}{)}

\PY{c+c1}{\PYZsh{} We keep columns in the desired order}
\PY{n}{prices\PYZus{}local} \PY{o}{=} \PY{n}{prices\PYZus{}local}\PY{p}{[}\PY{n}{price\PYZus{}cols}\PY{p}{]}

\PY{c+c1}{\PYZsh{} We fix potential European decimal formats (\PYZdq{}13165,5\PYZdq{} → 13165.5) and cast to numeric}
\PY{k}{for} \PY{n}{c} \PY{o+ow}{in} \PY{n}{mapped\PYZus{}cols}\PY{p}{:}
    \PY{n}{prices\PYZus{}local}\PY{p}{[}\PY{n}{c}\PY{p}{]} \PY{o}{=} \PY{p}{(}
        \PY{n}{prices\PYZus{}local}\PY{p}{[}\PY{n}{c}\PY{p}{]}\PY{o}{.}\PY{n}{astype}\PY{p}{(}\PY{n+nb}{str}\PY{p}{)}
        \PY{o}{.}\PY{n}{str}\PY{o}{.}\PY{n}{replace}\PY{p}{(}\PY{l+s+s2}{\PYZdq{}}\PY{l+s+s2}{ }\PY{l+s+s2}{\PYZdq{}}\PY{p}{,} \PY{l+s+s2}{\PYZdq{}}\PY{l+s+s2}{\PYZdq{}}\PY{p}{,} \PY{n}{regex}\PY{o}{=}\PY{k+kc}{False}\PY{p}{)}
        \PY{o}{.}\PY{n}{str}\PY{o}{.}\PY{n}{replace}\PY{p}{(}\PY{l+s+s2}{\PYZdq{}}\PY{l+s+s2}{,}\PY{l+s+s2}{\PYZdq{}}\PY{p}{,} \PY{l+s+s2}{\PYZdq{}}\PY{l+s+s2}{.}\PY{l+s+s2}{\PYZdq{}}\PY{p}{,} \PY{n}{regex}\PY{o}{=}\PY{k+kc}{False}\PY{p}{)}
    \PY{p}{)}
    \PY{n}{prices\PYZus{}local}\PY{p}{[}\PY{n}{c}\PY{p}{]} \PY{o}{=} \PY{n}{pd}\PY{o}{.}\PY{n}{to\PYZus{}numeric}\PY{p}{(}\PY{n}{prices\PYZus{}local}\PY{p}{[}\PY{n}{c}\PY{p}{]}\PY{p}{,} \PY{n}{errors}\PY{o}{=}\PY{l+s+s2}{\PYZdq{}}\PY{l+s+s2}{coerce}\PY{l+s+s2}{\PYZdq{}}\PY{p}{)}

\PY{c+c1}{\PYZsh{} Drop rows that are entirely empty, just in case}
\PY{n}{prices\PYZus{}local} \PY{o}{=} \PY{n}{prices\PYZus{}local}\PY{o}{.}\PY{n}{dropna}\PY{p}{(}\PY{n}{how}\PY{o}{=}\PY{l+s+s2}{\PYZdq{}}\PY{l+s+s2}{all}\PY{l+s+s2}{\PYZdq{}}\PY{p}{)}

\PY{n}{prices\PYZus{}local}\PY{o}{.}\PY{n}{head}\PY{p}{(}\PY{p}{)}
\end{Verbatim}
\end{tcolorbox}

    \begin{Verbatim}[commandchars=\\\{\}]
/tmp/ipykernel\_9051/22146945.py:40: UserWarning: Parsing dates in \%m/\%d/\%Y
format when dayfirst=True was specified. Pass `dayfirst=False` or specify a
format to silence this warning.
  df[date\_col] = pd.to\_datetime(df[date\_col], dayfirst=True)
    \end{Verbatim}

            \begin{tcolorbox}[breakable, size=fbox, boxrule=.5pt, pad at break*=1mm, opacityfill=0]
\begin{Verbatim}[commandchars=\\\{\}]
                     TSX          CAC          DAX  Eurostoxx50     NIKKEI225  \textbackslash{}
date
2007-03-30  13165.500000  5634.160156  6917.029785  4181.029785  17287.650391
2007-04-02  13265.799805  5645.560059  6937.169922  4189.549805  17028.410156
2007-04-03  13361.200195  5711.910156  7045.560059  4246.299805  17244.050781
2007-04-04  13448.299805  5739.009766  7073.910156  4261.830078  17544.089844
2007-04-05  13425.000000  5741.379883  7099.910156  4271.540039  17491.419922

            FTSE100    SP500  IBOVESPA
date
2007-03-30  6308.03  1420.86     45805
2007-04-02  6315.53  1424.55     45597
2007-04-03  6366.11  1437.77     46288
2007-04-04  6364.70  1439.37     46554
2007-04-05  6397.34  1443.76     46647
\end{Verbatim}
\end{tcolorbox}
        
    The table above shows the first few observations of the cleaned price
data, in local currency for each index. Each row is a trading day and
each column corresponds to one equity index (Canadian TSX, French CAC,
German DAX, etc.).

The levels are very different across countries: for instance, the S\&P
500 is around 1,400, while the Ibovespa is above 45,000 over the same
dates. This is exactly why we work with returns rather than price
levels, and why we later convert everything to a common currency (USD)
before doing any portfolio analysis.

    \hypertarget{extracting-the-fx-data-from-the-excel-file}{%
\subsubsection{1.3 Extracting the FX data from the Excel
file}\label{extracting-the-fx-data-from-the-excel-file}}

The same Excel sheet also contains FX rates against the USD, stored to
the right of the equity index block. To express all index returns in a
common currency (USD), we first extract this FX block from the raw data.

Concretely, we:

\begin{itemize}
    \item locate the FX columns that come after the last equity index (\texttt{IBOVESPA}),
    \item take the rows where the currency codes (CAD, EUR, JPY, GBP, USD, BRL) are given,
    \item keep the underlying FX time series (USD per unit of local currency),
    \item and align their dates with the equity price data in \texttt{prices\_local}.
\end{itemize}

The resulting object \texttt{fx\_vals} is a first, ``raw'\,' FX table
that we will clean and then use to convert all indices into USD prices.

    \begin{tcolorbox}[breakable, size=fbox, boxrule=1pt, pad at break*=1mm,colback=cellbackground, colframe=cellborder]
\begin{Verbatim}[commandchars=\\\{\}]
\PY{c+c1}{\PYZsh{} FX block: extract the FX table from the raw sheet}

\PY{n}{fx\PYZus{}start\PYZus{}col} \PY{o}{=} \PY{n}{df}\PY{o}{.}\PY{n}{columns}\PY{o}{.}\PY{n}{get\PYZus{}loc}\PY{p}{(}\PY{l+s+s2}{\PYZdq{}}\PY{l+s+s2}{IBOVESPA}\PY{l+s+s2}{\PYZdq{}}\PY{p}{)} \PY{o}{+} \PY{l+m+mi}{2}   \PY{c+c1}{\PYZsh{} FX columns start a bit to the right}
\PY{n}{fx\PYZus{}end\PYZus{}col}   \PY{o}{=} \PY{n}{fx\PYZus{}start\PYZus{}col} \PY{o}{+} \PY{l+m+mi}{7}                     \PY{c+c1}{\PYZsh{} CAD, EUR, EUR, EUR, JPY, GBP, USD, BRL}

\PY{n}{fx\PYZus{}raw} \PY{o}{=} \PY{n}{raw}\PY{o}{.}\PY{n}{iloc}\PY{p}{[}\PY{n}{header\PYZus{}row}\PY{o}{\PYZhy{}}\PY{l+m+mi}{2}\PY{p}{:}\PY{p}{,} \PY{n}{fx\PYZus{}start\PYZus{}col}\PY{p}{:}\PY{n}{fx\PYZus{}end\PYZus{}col}\PY{o}{+}\PY{l+m+mi}{1}\PY{p}{]}\PY{o}{.}\PY{n}{copy}\PY{p}{(}\PY{p}{)}

\PY{c+c1}{\PYZsh{} First row in this block contains the currency codes}
\PY{n}{currencies} \PY{o}{=} \PY{n}{fx\PYZus{}raw}\PY{o}{.}\PY{n}{iloc}\PY{p}{[}\PY{l+m+mi}{0}\PY{p}{]}\PY{o}{.}\PY{n}{astype}\PY{p}{(}\PY{n+nb}{str}\PY{p}{)}\PY{o}{.}\PY{n}{str}\PY{o}{.}\PY{n}{strip}\PY{p}{(}\PY{p}{)}\PY{o}{.}\PY{n}{values}

\PY{c+c1}{\PYZsh{} FX values start a few rows below}
\PY{n}{fx\PYZus{}vals} \PY{o}{=} \PY{n}{fx\PYZus{}raw}\PY{o}{.}\PY{n}{iloc}\PY{p}{[}\PY{l+m+mi}{3}\PY{p}{:}\PY{p}{]}\PY{o}{.}\PY{n}{copy}\PY{p}{(}\PY{p}{)}
\PY{n}{fx\PYZus{}vals}\PY{o}{.}\PY{n}{columns} \PY{o}{=} \PY{n}{currencies}

\PY{c+c1}{\PYZsh{} We align FX data on the same date index as the equity prices}
\PY{n}{fx\PYZus{}vals}\PY{o}{.}\PY{n}{index} \PY{o}{=} \PY{n}{prices\PYZus{}local}\PY{o}{.}\PY{n}{index}
\end{Verbatim}
\end{tcolorbox}

    We then clean the FX series. We first convert all FX entries to numeric
values (fixing the comma vs.~dot issue), and deal with the fact that EUR
appears several times by keeping only one EUR column. In the end we
extract a neat \texttt{fx} dataframe with one column per currency (CAD,
EUR, JPY, GBP, USD, BRL), indexed by the same dates as the equity
indices.

    \begin{tcolorbox}[breakable, size=fbox, boxrule=1pt, pad at break*=1mm,colback=cellbackground, colframe=cellborder]
\begin{Verbatim}[commandchars=\\\{\}]
\PY{n}{fx\PYZus{}vals} \PY{o}{=} \PY{n}{fx\PYZus{}vals}\PY{o}{.}\PY{n}{apply}\PY{p}{(}\PY{k}{lambda} \PY{n}{s}\PY{p}{:} \PY{n}{pd}\PY{o}{.}\PY{n}{to\PYZus{}numeric}\PY{p}{(}
    \PY{n}{s}\PY{o}{.}\PY{n}{astype}\PY{p}{(}\PY{n+nb}{str}\PY{p}{)}\PY{o}{.}\PY{n}{str}\PY{o}{.}\PY{n}{replace}\PY{p}{(}\PY{l+s+s2}{\PYZdq{}}\PY{l+s+s2}{ }\PY{l+s+s2}{\PYZdq{}}\PY{p}{,} \PY{l+s+s2}{\PYZdq{}}\PY{l+s+s2}{\PYZdq{}}\PY{p}{,} \PY{n}{regex}\PY{o}{=}\PY{k+kc}{False}\PY{p}{)}\PY{o}{.}\PY{n}{str}\PY{o}{.}\PY{n}{replace}\PY{p}{(}\PY{l+s+s2}{\PYZdq{}}\PY{l+s+s2}{,}\PY{l+s+s2}{\PYZdq{}}\PY{p}{,} \PY{l+s+s2}{\PYZdq{}}\PY{l+s+s2}{.}\PY{l+s+s2}{\PYZdq{}}\PY{p}{,} \PY{n}{regex}\PY{o}{=}\PY{k+kc}{False}\PY{p}{)}\PY{p}{,}
    \PY{n}{errors}\PY{o}{=}\PY{l+s+s2}{\PYZdq{}}\PY{l+s+s2}{coerce}\PY{l+s+s2}{\PYZdq{}}
\PY{p}{)}\PY{p}{)}

\PY{k}{if} \PY{p}{(}\PY{n}{fx\PYZus{}vals}\PY{o}{.}\PY{n}{columns} \PY{o}{==} \PY{l+s+s2}{\PYZdq{}}\PY{l+s+s2}{EUR}\PY{l+s+s2}{\PYZdq{}}\PY{p}{)}\PY{o}{.}\PY{n}{any}\PY{p}{(}\PY{p}{)}\PY{p}{:}
    \PY{n}{eur\PYZus{}first} \PY{o}{=} \PY{n}{fx\PYZus{}vals}\PY{o}{.}\PY{n}{loc}\PY{p}{[}\PY{p}{:}\PY{p}{,} \PY{n}{fx\PYZus{}vals}\PY{o}{.}\PY{n}{columns} \PY{o}{==} \PY{l+s+s2}{\PYZdq{}}\PY{l+s+s2}{EUR}\PY{l+s+s2}{\PYZdq{}}\PY{p}{]}\PY{o}{.}\PY{n}{iloc}\PY{p}{[}\PY{p}{:}\PY{p}{,} \PY{l+m+mi}{0}\PY{p}{]}
    \PY{n}{fx\PYZus{}vals} \PY{o}{=} \PY{n}{fx\PYZus{}vals}\PY{o}{.}\PY{n}{drop}\PY{p}{(}\PY{n}{columns}\PY{o}{=}\PY{l+s+s2}{\PYZdq{}}\PY{l+s+s2}{EUR}\PY{l+s+s2}{\PYZdq{}}\PY{p}{)}
    \PY{n}{fx\PYZus{}vals}\PY{p}{[}\PY{l+s+s2}{\PYZdq{}}\PY{l+s+s2}{EUR}\PY{l+s+s2}{\PYZdq{}}\PY{p}{]} \PY{o}{=} \PY{n}{eur\PYZus{}first}

\PY{c+c1}{\PYZsh{}}
\PY{n}{fx} \PY{o}{=} \PY{n}{fx\PYZus{}vals}\PY{p}{[}\PY{p}{[}\PY{l+s+s2}{\PYZdq{}}\PY{l+s+s2}{CAD}\PY{l+s+s2}{\PYZdq{}}\PY{p}{,} \PY{l+s+s2}{\PYZdq{}}\PY{l+s+s2}{EUR}\PY{l+s+s2}{\PYZdq{}}\PY{p}{,} \PY{l+s+s2}{\PYZdq{}}\PY{l+s+s2}{JPY}\PY{l+s+s2}{\PYZdq{}}\PY{p}{,} \PY{l+s+s2}{\PYZdq{}}\PY{l+s+s2}{GBP}\PY{l+s+s2}{\PYZdq{}}\PY{p}{,} \PY{l+s+s2}{\PYZdq{}}\PY{l+s+s2}{USD}\PY{l+s+s2}{\PYZdq{}}\PY{p}{,} \PY{l+s+s2}{\PYZdq{}}\PY{l+s+s2}{BRL}\PY{l+s+s2}{\PYZdq{}}\PY{p}{]}\PY{p}{]}\PY{o}{.}\PY{n}{copy}\PY{p}{(}\PY{p}{)}
\end{Verbatim}
\end{tcolorbox}

    We now link each equity index to its corresponding FX rate and convert
everything into USD. Concretely, we create a simple mapping from index
(TSX, CAC, DAX, etc.) to its currency (CAD, EUR, JPY, \dots), build an
FX matrix aligned with the price data, and then multiply local prices by
the relevant FX rate. The resulting \texttt{prices\_usd} dataframe
contains all eight indices expressed in USD, and we store it as
\texttt{prices} for the rest of the assignment.

    \begin{tcolorbox}[breakable, size=fbox, boxrule=1pt, pad at break*=1mm,colback=cellbackground, colframe=cellborder]
\begin{Verbatim}[commandchars=\\\{\}]
\PY{c+c1}{\PYZsh{} Map each equity index to its FX currency and convert prices to USD}

\PY{n}{fx\PYZus{}map} \PY{o}{=} \PY{p}{\PYZob{}}
    \PY{l+s+s2}{\PYZdq{}}\PY{l+s+s2}{TSX}\PY{l+s+s2}{\PYZdq{}}\PY{p}{:} \PY{l+s+s2}{\PYZdq{}}\PY{l+s+s2}{CAD}\PY{l+s+s2}{\PYZdq{}}\PY{p}{,}
    \PY{l+s+s2}{\PYZdq{}}\PY{l+s+s2}{CAC}\PY{l+s+s2}{\PYZdq{}}\PY{p}{:} \PY{l+s+s2}{\PYZdq{}}\PY{l+s+s2}{EUR}\PY{l+s+s2}{\PYZdq{}}\PY{p}{,}
    \PY{l+s+s2}{\PYZdq{}}\PY{l+s+s2}{DAX}\PY{l+s+s2}{\PYZdq{}}\PY{p}{:} \PY{l+s+s2}{\PYZdq{}}\PY{l+s+s2}{EUR}\PY{l+s+s2}{\PYZdq{}}\PY{p}{,}
    \PY{l+s+s2}{\PYZdq{}}\PY{l+s+s2}{Eurostoxx50}\PY{l+s+s2}{\PYZdq{}}\PY{p}{:} \PY{l+s+s2}{\PYZdq{}}\PY{l+s+s2}{EUR}\PY{l+s+s2}{\PYZdq{}}\PY{p}{,}
    \PY{l+s+s2}{\PYZdq{}}\PY{l+s+s2}{NIKKEI225}\PY{l+s+s2}{\PYZdq{}}\PY{p}{:} \PY{l+s+s2}{\PYZdq{}}\PY{l+s+s2}{JPY}\PY{l+s+s2}{\PYZdq{}}\PY{p}{,}
    \PY{l+s+s2}{\PYZdq{}}\PY{l+s+s2}{FTSE100}\PY{l+s+s2}{\PYZdq{}}\PY{p}{:} \PY{l+s+s2}{\PYZdq{}}\PY{l+s+s2}{GBP}\PY{l+s+s2}{\PYZdq{}}\PY{p}{,}
    \PY{l+s+s2}{\PYZdq{}}\PY{l+s+s2}{SP500}\PY{l+s+s2}{\PYZdq{}}\PY{p}{:} \PY{l+s+s2}{\PYZdq{}}\PY{l+s+s2}{USD}\PY{l+s+s2}{\PYZdq{}}\PY{p}{,}      \PY{c+c1}{\PYZsh{} already quoted in USD}
    \PY{l+s+s2}{\PYZdq{}}\PY{l+s+s2}{IBOVESPA}\PY{l+s+s2}{\PYZdq{}}\PY{p}{:} \PY{l+s+s2}{\PYZdq{}}\PY{l+s+s2}{BRL}\PY{l+s+s2}{\PYZdq{}}\PY{p}{,}
\PY{p}{\PYZcb{}}

\PY{c+c1}{\PYZsh{} Build an FX matrix aligned with the equity price series}
\PY{n}{fx\PYZus{}per\PYZus{}index} \PY{o}{=} \PY{n}{pd}\PY{o}{.}\PY{n}{DataFrame}\PY{p}{(}\PY{n}{index}\PY{o}{=}\PY{n}{prices\PYZus{}local}\PY{o}{.}\PY{n}{index}\PY{p}{,}
                            \PY{n}{columns}\PY{o}{=}\PY{n}{prices\PYZus{}local}\PY{o}{.}\PY{n}{columns}\PY{p}{)}

\PY{k}{for} \PY{n}{idx} \PY{o+ow}{in} \PY{n}{prices\PYZus{}local}\PY{o}{.}\PY{n}{columns}\PY{p}{:}
    \PY{n}{fx\PYZus{}per\PYZus{}index}\PY{p}{[}\PY{n}{idx}\PY{p}{]} \PY{o}{=} \PY{n}{fx}\PY{p}{[}\PY{n}{fx\PYZus{}map}\PY{p}{[}\PY{n}{idx}\PY{p}{]}\PY{p}{]}

\PY{c+c1}{\PYZsh{} Convert local\PYZhy{}currency prices to USD}
\PY{n}{prices\PYZus{}usd} \PY{o}{=} \PY{n}{prices\PYZus{}local} \PY{o}{*} \PY{n}{fx\PYZus{}per\PYZus{}index}

\PY{c+c1}{\PYZsh{} From now on we work with USD prices only}
\PY{n}{prices} \PY{o}{=} \PY{n}{prices\PYZus{}usd}\PY{o}{.}\PY{n}{copy}\PY{p}{(}\PY{p}{)}

\PY{n}{prices}\PY{o}{.}\PY{n}{head}\PY{p}{(}\PY{p}{)}  \PY{c+c1}{\PYZsh{} quick sanity check}
\end{Verbatim}
\end{tcolorbox}

            \begin{tcolorbox}[breakable, size=fbox, boxrule=.5pt, pad at break*=1mm, opacityfill=0]
\begin{Verbatim}[commandchars=\\\{\}]
                     TSX          CAC          DAX  Eurostoxx50   NIKKEI225  \textbackslash{}
date
2007-03-30  11405.614953  7526.663284  9240.446256  5585.429328  146.737577
2007-04-02  11473.616783  7548.147672  9275.037809  5601.453227  144.605259
2007-04-03  11532.693226  7611.108859  9388.194687  5658.185998  145.074199
2007-04-04  11606.366871  7675.345922  9460.640369  5699.767284  147.721236
2007-04-05  11664.781125  7711.827200  9536.606421  5737.536852  147.365213

                 FTSE100    SP500      IBOVESPA
date
2007-03-30  12316.523195  1420.86  22289.537490
2007-04-02  12493.381446  1424.55  22313.211129
2007-04-03  12562.875108  1437.77  22758.235808
2007-04-04  12579.205809  1439.37  22947.677004
2007-04-05  12612.304631  1443.76  23008.306221
\end{Verbatim}
\end{tcolorbox}
        
    The table above shows the first few observations of our
\textbf{USD-denominated price
series}. Each column corresponds to one equity index, now expressed in
the same currency (USD):

\begin{itemize}
  \item TSX, CAC, DAX, Eurostoxx50, FTSE100 and SP500 are all in the \textbf{thousands of USD},
        as expected for large developed-market indices.
  \item NIKKEI225 appears around \textbf{145--150 USD} because the original JPY index level
        has been scaled by the JPY--USD exchange rate.
  \item IBOVESPA is in the \textbf{20,000+ USD} range after conversion from BRL.
\end{itemize}

The magnitudes look consistent with typical index levels once everything
is converted into USD, so we will use this \texttt{prices} dataframe as
the starting point to compute \textbf{daily USD log returns} in the next
step.

    \hypertarget{daily-log-returns-in-usd}{%
\subsubsection{1.4 Daily log returns in
USD}\label{daily-log-returns-in-usd}}

From the USD price series, we now compute \textbf{daily log returns} for
each index.

We take the log of the price ratio \(\log(P_t / P_{t-1})\) and drop the
first missing observation. The resulting \texttt{returns} dataframe will
be the main input for all the summary statistics and tests in the rest
of Question\textasciitilde1.

    \begin{tcolorbox}[breakable, size=fbox, boxrule=1pt, pad at break*=1mm,colback=cellbackground, colframe=cellborder]
\begin{Verbatim}[commandchars=\\\{\}]
\PY{c+c1}{\PYZsh{} Daily log returns (USD)}
\PY{n}{returns} \PY{o}{=} \PY{n}{np}\PY{o}{.}\PY{n}{log}\PY{p}{(}\PY{n}{prices} \PY{o}{/} \PY{n}{prices}\PY{o}{.}\PY{n}{shift}\PY{p}{(}\PY{l+m+mi}{1}\PY{p}{)}\PY{p}{)}\PY{o}{.}\PY{n}{dropna}\PY{p}{(}\PY{p}{)}

\PY{n}{returns}\PY{o}{.}\PY{n}{head}\PY{p}{(}\PY{p}{)}
\end{Verbatim}
\end{tcolorbox}

            \begin{tcolorbox}[breakable, size=fbox, boxrule=.5pt, pad at break*=1mm, opacityfill=0]
\begin{Verbatim}[commandchars=\\\{\}]
                 TSX       CAC       DAX  Eurostoxx50  NIKKEI225   FTSE100  \textbackslash{}
date
2007-04-02  0.005944  0.002850  0.003737     0.002865  -0.014638  0.014257
2007-04-03  0.005136  0.008307  0.012126     0.010077   0.003238  0.005547
2007-04-04  0.006368  0.008404  0.007687     0.007322   0.018082  0.001299
2007-04-05  0.005020  0.004742  0.007998     0.006605  -0.002413  0.002628
2007-04-10  0.003777  0.004326  0.009359     0.006982   0.006290  0.003346

               SP500  IBOVESPA
date
2007-04-02  0.002594  0.001062
2007-04-03  0.009237  0.019748
2007-04-04  0.001112  0.008290
2007-04-05  0.003045  0.002639
2007-04-10  0.003202  0.011431
\end{Verbatim}
\end{tcolorbox}
        
    The table shows the first few \textbf{daily USD log returns} for our
eight indices.\\
Each row corresponds to a trading day and each column to one market. The
numbers are small (mostly between \(-2\%\) and \(+2\%\)), which is
exactly what we expect for daily equity moves.

We already see some cross-sectional differences: for example, on 2 April
2007 the Nikkei drops by about (-1.5\%) while the European indices and
the S\&P\textasciitilde500 all post small positive returns. These
\texttt{returns} series are what we will use to compute summary
statistics, correlations and normality tests in the rest of the
question.

    We first look at the time series of \textbf{daily USD log returns} for
all eight indices on the same plot.\\
To get a sense of the typical scale of daily moves, we also add
horizontal bands at \(\pm 1\sigma\) and \(\pm 2\sigma\) (using the
S\&P\textasciitilde500 volatility as a reference).

This gives a global picture of how volatile each market is and how
crises show up in the data: large common spikes appear around 2008--2009
and again in 2020, while most of the time returns oscillate in a narrow
band around zero. The figure also suggests that shocks are largely
synchronised across markets, which will matter later when we look at
correlations and mean--variance portfolios.

    \begin{tcolorbox}[breakable, size=fbox, boxrule=1pt, pad at break*=1mm,colback=cellbackground, colframe=cellborder]
\begin{Verbatim}[commandchars=\\\{\}]
\PY{n}{plt}\PY{o}{.}\PY{n}{figure}\PY{p}{(}\PY{n}{figsize}\PY{o}{=}\PY{p}{(}\PY{l+m+mi}{10}\PY{p}{,} \PY{l+m+mi}{5}\PY{p}{)}\PY{p}{)}

\PY{c+c1}{\PYZsh{} We plot all daily USD log returns on the same axis}
\PY{n}{returns}\PY{o}{.}\PY{n}{plot}\PY{p}{(}\PY{n}{ax}\PY{o}{=}\PY{n}{plt}\PY{o}{.}\PY{n}{gca}\PY{p}{(}\PY{p}{)}\PY{p}{,} \PY{n}{linewidth}\PY{o}{=}\PY{l+m+mf}{0.7}\PY{p}{,} \PY{n}{alpha}\PY{o}{=}\PY{l+m+mf}{0.8}\PY{p}{)}

\PY{n}{plt}\PY{o}{.}\PY{n}{title}\PY{p}{(}\PY{l+s+s2}{\PYZdq{}}\PY{l+s+s2}{Daily log returns of equity indices (USD)}\PY{l+s+s2}{\PYZdq{}}\PY{p}{)}
\PY{n}{plt}\PY{o}{.}\PY{n}{ylabel}\PY{p}{(}\PY{l+s+s2}{\PYZdq{}}\PY{l+s+s2}{Log return}\PY{l+s+s2}{\PYZdq{}}\PY{p}{)}
\PY{n}{plt}\PY{o}{.}\PY{n}{xlabel}\PY{p}{(}\PY{l+s+s2}{\PYZdq{}}\PY{l+s+s2}{Date}\PY{l+s+s2}{\PYZdq{}}\PY{p}{)}

\PY{c+c1}{\PYZsh{} Here we compute overall daily volatility from SP500 (for the bands)}
\PY{n}{sigma} \PY{o}{=} \PY{n}{returns}\PY{p}{[}\PY{l+s+s2}{\PYZdq{}}\PY{l+s+s2}{SP500}\PY{l+s+s2}{\PYZdq{}}\PY{p}{]}\PY{o}{.}\PY{n}{std}\PY{p}{(}\PY{n}{ddof}\PY{o}{=}\PY{l+m+mi}{1}\PY{p}{)}

\PY{c+c1}{\PYZsh{} 0, ±1σ, ±2σ lines in BLUE}
\PY{n}{plt}\PY{o}{.}\PY{n}{axhline}\PY{p}{(}\PY{l+m+mi}{0}\PY{p}{,}        \PY{n}{color}\PY{o}{=}\PY{l+s+s2}{\PYZdq{}}\PY{l+s+s2}{blue}\PY{l+s+s2}{\PYZdq{}}\PY{p}{,} \PY{n}{linestyle}\PY{o}{=}\PY{l+s+s2}{\PYZdq{}}\PY{l+s+s2}{\PYZhy{}\PYZhy{}}\PY{l+s+s2}{\PYZdq{}}\PY{p}{,} \PY{n}{linewidth}\PY{o}{=}\PY{l+m+mi}{1}\PY{p}{)}
\PY{n}{plt}\PY{o}{.}\PY{n}{axhline}\PY{p}{(}\PY{n}{sigma}\PY{p}{,}    \PY{n}{color}\PY{o}{=}\PY{l+s+s2}{\PYZdq{}}\PY{l+s+s2}{blue}\PY{l+s+s2}{\PYZdq{}}\PY{p}{,} \PY{n}{linestyle}\PY{o}{=}\PY{l+s+s2}{\PYZdq{}}\PY{l+s+s2}{\PYZhy{}\PYZhy{}}\PY{l+s+s2}{\PYZdq{}}\PY{p}{,} \PY{n}{linewidth}\PY{o}{=}\PY{l+m+mi}{1}\PY{p}{)}
\PY{n}{plt}\PY{o}{.}\PY{n}{axhline}\PY{p}{(}\PY{o}{\PYZhy{}}\PY{n}{sigma}\PY{p}{,}   \PY{n}{color}\PY{o}{=}\PY{l+s+s2}{\PYZdq{}}\PY{l+s+s2}{blue}\PY{l+s+s2}{\PYZdq{}}\PY{p}{,} \PY{n}{linestyle}\PY{o}{=}\PY{l+s+s2}{\PYZdq{}}\PY{l+s+s2}{\PYZhy{}\PYZhy{}}\PY{l+s+s2}{\PYZdq{}}\PY{p}{,} \PY{n}{linewidth}\PY{o}{=}\PY{l+m+mi}{1}\PY{p}{)}
\PY{n}{plt}\PY{o}{.}\PY{n}{axhline}\PY{p}{(}\PY{l+m+mi}{2}\PY{o}{*}\PY{n}{sigma}\PY{p}{,}  \PY{n}{color}\PY{o}{=}\PY{l+s+s2}{\PYZdq{}}\PY{l+s+s2}{blue}\PY{l+s+s2}{\PYZdq{}}\PY{p}{,} \PY{n}{linestyle}\PY{o}{=}\PY{l+s+s2}{\PYZdq{}}\PY{l+s+s2}{\PYZhy{}\PYZhy{}}\PY{l+s+s2}{\PYZdq{}}\PY{p}{,} \PY{n}{linewidth}\PY{o}{=}\PY{l+m+mi}{1}\PY{p}{)}
\PY{n}{plt}\PY{o}{.}\PY{n}{axhline}\PY{p}{(}\PY{o}{\PYZhy{}}\PY{l+m+mi}{2}\PY{o}{*}\PY{n}{sigma}\PY{p}{,} \PY{n}{color}\PY{o}{=}\PY{l+s+s2}{\PYZdq{}}\PY{l+s+s2}{blue}\PY{l+s+s2}{\PYZdq{}}\PY{p}{,} \PY{n}{linestyle}\PY{o}{=}\PY{l+s+s2}{\PYZdq{}}\PY{l+s+s2}{\PYZhy{}\PYZhy{}}\PY{l+s+s2}{\PYZdq{}}\PY{p}{,} \PY{n}{linewidth}\PY{o}{=}\PY{l+m+mi}{1}\PY{p}{)}

\PY{c+c1}{\PYZsh{} We put the legend on the right so it doesn\PYZsq{}t bother our analysis }
\PY{n}{plt}\PY{o}{.}\PY{n}{legend}\PY{p}{(}\PY{n}{loc}\PY{o}{=}\PY{l+s+s2}{\PYZdq{}}\PY{l+s+s2}{center left}\PY{l+s+s2}{\PYZdq{}}\PY{p}{,} \PY{n}{bbox\PYZus{}to\PYZus{}anchor}\PY{o}{=}\PY{p}{(}\PY{l+m+mi}{1}\PY{p}{,} \PY{l+m+mf}{0.5}\PY{p}{)}\PY{p}{)}

\PY{n}{plt}\PY{o}{.}\PY{n}{grid}\PY{p}{(}\PY{n}{alpha}\PY{o}{=}\PY{l+m+mf}{0.2}\PY{p}{)}
\PY{n}{plt}\PY{o}{.}\PY{n}{tight\PYZus{}layout}\PY{p}{(}\PY{p}{)}
\PY{n}{plt}\PY{o}{.}\PY{n}{show}\PY{p}{(}\PY{p}{)}
\end{Verbatim}
\end{tcolorbox}

    \begin{center}
    \adjustimage{max size={0.9\linewidth}{0.9\paperheight}}{Work-7_files/Work-7_21_0.png}
    \end{center}
    { \hspace*{\fill} \\}
    
    The time-series plot of daily USD log returns shows that all indices
fluctuate around zero, with frequent small moves and occasional large
spikes. Most observations lie between the \(\pm 1\sigma\) bands, and a
large majority remain within \(\pm 2\sigma\), which is in line with the
idea of ``normal'\,' day-to-day volatility.

However, during crisis periods (2008--2009, some episodes around
2011--2016, and especially 2020 with Covid) many points clearly break
through the \(\pm 2\sigma\) lines. These outliers highlight how extreme
those episodes are compared with the usual daily noise. The fact that
several indices cross the bands at the same time also confirms that
markets tend to move together on stressful days, which is consistent
with the high cross-correlations we report later.

    To look more closely at this time-varying volatility, we now focus on a
single market (the S\&P 500) and compute a 15-day rolling standard
deviation of its daily log returns. This rolling volatility measure
makes the volatility clustering more visible and is directly relevant
for assessing the i.i.d.~Gaussian assumption.

    \begin{tcolorbox}[breakable, size=fbox, boxrule=1pt, pad at break*=1mm,colback=cellbackground, colframe=cellborder]
\begin{Verbatim}[commandchars=\\\{\}]
\PY{n}{roll} \PY{o}{=} \PY{n}{returns}\PY{o}{.}\PY{n}{rolling}\PY{p}{(}\PY{l+m+mi}{15}\PY{p}{)}\PY{o}{.}\PY{n}{std}\PY{p}{(}\PY{p}{)}

\PY{n}{plt}\PY{o}{.}\PY{n}{figure}\PY{p}{(}\PY{n}{figsize}\PY{o}{=}\PY{p}{(}\PY{l+m+mi}{9}\PY{p}{,}\PY{l+m+mi}{5}\PY{p}{)}\PY{p}{)}
\PY{n}{plt}\PY{o}{.}\PY{n}{plot}\PY{p}{(}\PY{n}{roll}\PY{o}{.}\PY{n}{index}\PY{p}{,} \PY{n}{roll}\PY{p}{[}\PY{l+s+s2}{\PYZdq{}}\PY{l+s+s2}{SP500}\PY{l+s+s2}{\PYZdq{}}\PY{p}{]}\PY{p}{)}
\PY{n}{plt}\PY{o}{.}\PY{n}{title}\PY{p}{(}\PY{l+s+s2}{\PYZdq{}}\PY{l+s+s2}{SP500 rolling volatility (15\PYZhy{}day window, USD log returns)}\PY{l+s+s2}{\PYZdq{}}\PY{p}{)}
\PY{n}{plt}\PY{o}{.}\PY{n}{xlabel}\PY{p}{(}\PY{l+s+s2}{\PYZdq{}}\PY{l+s+s2}{Date}\PY{l+s+s2}{\PYZdq{}}\PY{p}{)}
\PY{n}{plt}\PY{o}{.}\PY{n}{ylabel}\PY{p}{(}\PY{l+s+s2}{\PYZdq{}}\PY{l+s+s2}{Rolling std}\PY{l+s+s2}{\PYZdq{}}\PY{p}{)}
\PY{n}{plt}\PY{o}{.}\PY{n}{tight\PYZus{}layout}\PY{p}{(}\PY{p}{)}
\PY{n}{plt}\PY{o}{.}\PY{n}{show}\PY{p}{(}\PY{p}{)}
\end{Verbatim}
\end{tcolorbox}

    \begin{center}
    \adjustimage{max size={0.9\linewidth}{0.9\paperheight}}{Work-7_files/Work-7_24_0.png}
    \end{center}
    { \hspace*{\fill} \\}
    
    The 15-day rolling volatility of S\&P 500 returns confirms very strong
time variation in risk.

Volatility is low and fairly stable in ``normal'\,' periods (around
0.5--1\% per day), but it jumps to much higher levels during crises:
above 4--6\% in 2008--2009 and even higher around March 2020.

These long stretches of calm followed by volatility bursts are clearly
inconsistent with the constant-variance assumption behind an
i.i.d.~Gaussian model, and they suggest that any realistic model should
allow for volatility clustering (for instance through GARCH-type
dynamics).

    \hypertarget{summary-statistics-of-daily-usd-returns}{%
\subsubsection{1.5 Summary statistics of daily USD
returns}\label{summary-statistics-of-daily-usd-returns}}

To summarise the behaviour of the USD daily log returns, we now compute
a set of basic moments for each index: mean, standard deviation,
skewness, (excess) kurtosis and the first-order autocorrelation. This
gives us a compact view of average performance, risk and departures from
normality before moving on to more formal tests.

    \begin{tcolorbox}[breakable, size=fbox, boxrule=1pt, pad at break*=1mm,colback=cellbackground, colframe=cellborder]
\begin{Verbatim}[commandchars=\\\{\}]
\PY{n}{summary} \PY{o}{=} \PY{n}{pd}\PY{o}{.}\PY{n}{DataFrame}\PY{p}{(}\PY{n}{index}\PY{o}{=}\PY{n}{returns}\PY{o}{.}\PY{n}{columns}\PY{p}{)}

\PY{n}{summary}\PY{p}{[}\PY{l+s+s2}{\PYZdq{}}\PY{l+s+s2}{mean}\PY{l+s+s2}{\PYZdq{}}\PY{p}{]} \PY{o}{=} \PY{n}{returns}\PY{o}{.}\PY{n}{mean}\PY{p}{(}\PY{p}{)}
\PY{n}{summary}\PY{p}{[}\PY{l+s+s2}{\PYZdq{}}\PY{l+s+s2}{std}\PY{l+s+s2}{\PYZdq{}}\PY{p}{]} \PY{o}{=} \PY{n}{returns}\PY{o}{.}\PY{n}{std}\PY{p}{(}\PY{n}{ddof}\PY{o}{=}\PY{l+m+mi}{1}\PY{p}{)}
\PY{n}{summary}\PY{p}{[}\PY{l+s+s2}{\PYZdq{}}\PY{l+s+s2}{skew}\PY{l+s+s2}{\PYZdq{}}\PY{p}{]} \PY{o}{=} \PY{n}{returns}\PY{o}{.}\PY{n}{apply}\PY{p}{(}\PY{n}{skew}\PY{p}{)}

\PY{c+c1}{\PYZsh{} \PYZdq{}kurtosis\PYZdq{} in the usual sense (not centred at 3)}
\PY{n}{summary}\PY{p}{[}\PY{l+s+s2}{\PYZdq{}}\PY{l+s+s2}{kurtosis}\PY{l+s+s2}{\PYZdq{}}\PY{p}{]} \PY{o}{=} \PY{n}{returns}\PY{o}{.}\PY{n}{apply}\PY{p}{(}\PY{k}{lambda} \PY{n}{x}\PY{p}{:} \PY{n}{kurtosis}\PY{p}{(}\PY{n}{x}\PY{p}{,} \PY{n}{fisher}\PY{o}{=}\PY{k+kc}{False}\PY{p}{)}\PY{p}{)}

\PY{c+c1}{\PYZsh{} excess kurtosis = kurtosis − 3}
\PY{n}{summary}\PY{p}{[}\PY{l+s+s2}{\PYZdq{}}\PY{l+s+s2}{excess\PYZus{}kurtosis}\PY{l+s+s2}{\PYZdq{}}\PY{p}{]} \PY{o}{=} \PY{n}{returns}\PY{o}{.}\PY{n}{apply}\PY{p}{(}\PY{k}{lambda} \PY{n}{x}\PY{p}{:} \PY{n}{kurtosis}\PY{p}{(}\PY{n}{x}\PY{p}{,} \PY{n}{fisher}\PY{o}{=}\PY{k+kc}{True}\PY{p}{)}\PY{p}{)}

\PY{c+c1}{\PYZsh{} first\PYZhy{}order autocorrelation of daily returns}
\PY{n}{summary}\PY{p}{[}\PY{l+s+s2}{\PYZdq{}}\PY{l+s+s2}{autocorr\PYZus{}lag1}\PY{l+s+s2}{\PYZdq{}}\PY{p}{]} \PY{o}{=} \PY{n}{returns}\PY{o}{.}\PY{n}{apply}\PY{p}{(}\PY{k}{lambda} \PY{n}{x}\PY{p}{:} \PY{n}{x}\PY{o}{.}\PY{n}{autocorr}\PY{p}{(}\PY{n}{lag}\PY{o}{=}\PY{l+m+mi}{1}\PY{p}{)}\PY{p}{)}

\PY{n}{summary}
\end{Verbatim}
\end{tcolorbox}

            \begin{tcolorbox}[breakable, size=fbox, boxrule=.5pt, pad at break*=1mm, opacityfill=0]
\begin{Verbatim}[commandchars=\\\{\}]
                 mean       std      skew   kurtosis  excess\_kurtosis  \textbackslash{}
TSX          0.000117  0.015251 -0.865886  15.285325        12.285325
CAC          0.000025  0.018352  0.372059  20.554237        17.554237
DAX          0.000207  0.018058  0.410799  19.888278        16.888278
Eurostoxx50 -0.000041  0.018481  0.283613  19.952566        16.952566
NIKKEI225    0.000165  0.016163 -0.552175  19.710934        16.710934
FTSE100     -0.000065  0.015374 -0.274320  13.079991        10.079991
SP500        0.000374  0.013581 -0.637448  15.001456        12.001456
IBOVESPA    -0.000060  0.024794 -0.154722  14.059129        11.059129

             autocorr\_lag1
TSX               0.007233
CAC              -0.075810
DAX              -0.058452
Eurostoxx50      -0.077204
NIKKEI225        -0.199259
FTSE100          -0.019601
SP500            -0.154566
IBOVESPA          0.012955
\end{Verbatim}
\end{tcolorbox}
        
    This table already tells us a lot about the behaviour of the eight
indices.

First, the \textbf{average daily returns} are tiny compared to
volatility: they are on the order of a few basis points per day, with
the S\&P\textasciitilde500 having the highest mean (about 0.037\% per
day) and some markets such as Eurostoxx50 or FTSE100 even slightly
negative over the sample. At the daily horizon, the ``signal'\,' in the
mean is clearly dominated by noise.

By contrast, the \textbf{standard deviations} are much larger and differ
quite a lot across markets. Ibovespa stands out as the most volatile
index (around 2.5\% per day), followed by the Euro area indices (CAC,
DAX, Eurostoxx50, with volatilities around 1.8--1.9\%). TSX and
S\&P\textasciitilde500 are a bit less volatile, which matches the usual
intuition that emerging markets and some European markets are riskier
than large developed US/Canadian equities.

The \textbf{shape of the return distribution} is very far from Gaussian.
Skewness is often negative (TSX, Nikkei, FTSE, S\&P\textasciitilde500,
Ibovespa), which means that large negative moves are more frequent than
large positive ones. Above all, the excess kurtosis is extremely high
for every index (between roughly 10 and 18, instead of 0 under a normal
distribution). This points to very fat tails and frequent extreme
events, which is hard to reconcile with a simple Gaussian i.i.d.~model.

Finally, the \textbf{lag-1 autocorrelations} of returns are all close to
zero, sometimes slightly negative (European indices, Nikkei), sometimes
slightly positive (TSX, Ibovespa). This is broadly consistent with
weak-form market efficiency: daily returns themselves are almost
uncorrelated over time, even though their distribution is very
non-normal.

    Overall, the i.i.d.~Gaussian assumption is not really supported by the
data: independence at the daily frequency is roughly plausible, but the
normality part is strongly rejected because of the heavy tails and
skewness.

    \hypertarget{correlation-matrix-of-daily-usd-returns}{%
\subsubsection{1.6 Correlation matrix of daily USD
returns}\label{correlation-matrix-of-daily-usd-returns}}

We now compute the correlation matrix of daily USD log returns across
the eight indices. This gives a compact view of how strongly the
different equity markets move together and will be useful later when we
study diversification and mean--variance portfolios.

    \begin{tcolorbox}[breakable, size=fbox, boxrule=1pt, pad at break*=1mm,colback=cellbackground, colframe=cellborder]
\begin{Verbatim}[commandchars=\\\{\}]
\PY{n}{corr} \PY{o}{=} \PY{n}{returns}\PY{o}{.}\PY{n}{corr}\PY{p}{(}\PY{p}{)}
\PY{n}{corr}
\end{Verbatim}
\end{tcolorbox}

            \begin{tcolorbox}[breakable, size=fbox, boxrule=.5pt, pad at break*=1mm, opacityfill=0]
\begin{Verbatim}[commandchars=\\\{\}]
                  TSX       CAC       DAX  Eurostoxx50  NIKKEI225   FTSE100  \textbackslash{}
TSX          1.000000  0.677227  0.659697     0.661918   0.222340  0.716389
CAC          0.677227  1.000000  0.950523     0.984838   0.260700  0.841713
DAX          0.659697  0.950523  1.000000     0.966589   0.246152  0.809851
Eurostoxx50  0.661918  0.984838  0.966589     1.000000   0.248228  0.830241
NIKKEI225    0.222340  0.260700  0.246152     0.248228   1.000000  0.343642
FTSE100      0.716389  0.841713  0.809851     0.830241   0.343642  1.000000
SP500        0.771410  0.612419  0.615929     0.611192   0.099309  0.616687
IBOVESPA     0.685627  0.567331  0.540996     0.553301   0.238317  0.617852

                SP500  IBOVESPA
TSX          0.771410  0.685627
CAC          0.612419  0.567331
DAX          0.615929  0.540996
Eurostoxx50  0.611192  0.553301
NIKKEI225    0.099309  0.238317
FTSE100      0.616687  0.617852
SP500        1.000000  0.610080
IBOVESPA     0.610080  1.000000
\end{Verbatim}
\end{tcolorbox}
        
    The correlation matrix confirms that most developed markets move very
closely together in USD. CAC, DAX and Eurostoxx50 are almost perfectly
correlated, and TSX, FTSE100, SP500 and IBOVESPA all show fairly high
correlations with this European block (around 0.6--0.8). The Nikkei is
the only real outlier, with much lower correlations (around 0.2--0.3,
and only about 0.1 with the S\&P\textasciitilde500), so Japan is the
main source of diversification in this set of indices.

    To better visualise the dependence structure across markets, we now plot
the correlation matrix of daily USD log returns as a heatmap. This gives
us a quick, graphical view of which indices tend to move together and
which ones offer more diversification.

    \begin{tcolorbox}[breakable, size=fbox, boxrule=1pt, pad at break*=1mm,colback=cellbackground, colframe=cellborder]
\begin{Verbatim}[commandchars=\\\{\}]
\PY{c+c1}{\PYZsh{} Correlation matrix of daily USD log returns}
\PY{n}{corr} \PY{o}{=} \PY{n}{returns}\PY{o}{.}\PY{n}{corr}\PY{p}{(}\PY{p}{)}

\PY{n}{plt}\PY{o}{.}\PY{n}{figure}\PY{p}{(}\PY{n}{figsize}\PY{o}{=}\PY{p}{(}\PY{l+m+mi}{7}\PY{p}{,} \PY{l+m+mi}{6}\PY{p}{)}\PY{p}{)}
\PY{n}{im} \PY{o}{=} \PY{n}{plt}\PY{o}{.}\PY{n}{imshow}\PY{p}{(}\PY{n}{corr}\PY{p}{,} \PY{n}{cmap}\PY{o}{=}\PY{l+s+s2}{\PYZdq{}}\PY{l+s+s2}{viridis}\PY{l+s+s2}{\PYZdq{}}\PY{p}{,} \PY{n}{vmin}\PY{o}{=}\PY{o}{\PYZhy{}}\PY{l+m+mi}{1}\PY{p}{,} \PY{n}{vmax}\PY{o}{=}\PY{l+m+mi}{1}\PY{p}{)}

\PY{c+c1}{\PYZsh{} Ticks \PYZam{} labels}
\PY{n}{assets} \PY{o}{=} \PY{n}{corr}\PY{o}{.}\PY{n}{columns}\PY{o}{.}\PY{n}{tolist}\PY{p}{(}\PY{p}{)}
\PY{n}{plt}\PY{o}{.}\PY{n}{xticks}\PY{p}{(}\PY{n+nb}{range}\PY{p}{(}\PY{n+nb}{len}\PY{p}{(}\PY{n}{assets}\PY{p}{)}\PY{p}{)}\PY{p}{,} \PY{n}{assets}\PY{p}{,} \PY{n}{rotation}\PY{o}{=}\PY{l+m+mi}{45}\PY{p}{,} \PY{n}{ha}\PY{o}{=}\PY{l+s+s2}{\PYZdq{}}\PY{l+s+s2}{right}\PY{l+s+s2}{\PYZdq{}}\PY{p}{)}
\PY{n}{plt}\PY{o}{.}\PY{n}{yticks}\PY{p}{(}\PY{n+nb}{range}\PY{p}{(}\PY{n+nb}{len}\PY{p}{(}\PY{n}{assets}\PY{p}{)}\PY{p}{)}\PY{p}{,} \PY{n}{assets}\PY{p}{)}

\PY{n}{plt}\PY{o}{.}\PY{n}{title}\PY{p}{(}\PY{l+s+s2}{\PYZdq{}}\PY{l+s+s2}{Correlation matrix of daily USD log returns}\PY{l+s+s2}{\PYZdq{}}\PY{p}{)}
\PY{n}{plt}\PY{o}{.}\PY{n}{colorbar}\PY{p}{(}\PY{n}{im}\PY{p}{,} \PY{n}{label}\PY{o}{=}\PY{l+s+s2}{\PYZdq{}}\PY{l+s+s2}{Correlation}\PY{l+s+s2}{\PYZdq{}}\PY{p}{)}

\PY{n}{plt}\PY{o}{.}\PY{n}{tight\PYZus{}layout}\PY{p}{(}\PY{p}{)}
\PY{n}{plt}\PY{o}{.}\PY{n}{show}\PY{p}{(}\PY{p}{)}
\end{Verbatim}
\end{tcolorbox}

    \begin{center}
    \adjustimage{max size={0.9\linewidth}{0.9\paperheight}}{Work-7_files/Work-7_34_0.png}
    \end{center}
    { \hspace*{\fill} \\}
    
    The heatmap makes these patterns easy to see: the European indices, the
UK, North America and Brazil form a bright yellow ``cluster'' of highly
correlated markets. In contrast, the darker row and column for the
Nikkei show that Japanese equities are noticeably less connected to the
rest, which is exactly what we would look for when trying to diversify a
global equity portfolio.

    It is worth noting that these high correlations are obtained
\emph{after} converting all indices to USD. The strong within-Europe
block therefore cannot be attributed to a shared quotation currency
only, but rather to common economic and policy shocks that affect
European equity markets in a very similar way.

    \hypertarget{sp500-empirical-distribution-vs-fitted-normal}{%
\subsubsection{1.7 SP500 empirical distribution vs fitted
normal}\label{sp500-empirical-distribution-vs-fitted-normal}}

To get a more visual feeling for the (non-)normality of returns, we
focus on the S\&P 500 and compare its empirical distribution to a fitted
Gaussian.

We plot the histogram of daily USD log returns together with the normal
density using the sample mean and standard deviation. This helps us see
where the normal approximation breaks down (in particular in the tails
and around the centre), before moving on to formal normality tests.

    \begin{tcolorbox}[breakable, size=fbox, boxrule=1pt, pad at break*=1mm,colback=cellbackground, colframe=cellborder]
\begin{Verbatim}[commandchars=\\\{\}]
\PY{n}{asset} \PY{o}{=} \PY{l+s+s2}{\PYZdq{}}\PY{l+s+s2}{SP500}\PY{l+s+s2}{\PYZdq{}}
\PY{n}{x} \PY{o}{=} \PY{n}{returns}\PY{p}{[}\PY{n}{asset}\PY{p}{]}\PY{o}{.}\PY{n}{dropna}\PY{p}{(}\PY{p}{)}

\PY{n}{mu\PYZus{}hat}\PY{p}{,} \PY{n}{sigma\PYZus{}hat} \PY{o}{=} \PY{n}{x}\PY{o}{.}\PY{n}{mean}\PY{p}{(}\PY{p}{)}\PY{p}{,} \PY{n}{x}\PY{o}{.}\PY{n}{std}\PY{p}{(}\PY{n}{ddof}\PY{o}{=}\PY{l+m+mi}{1}\PY{p}{)}

\PY{n}{plt}\PY{o}{.}\PY{n}{figure}\PY{p}{(}\PY{n}{figsize}\PY{o}{=}\PY{p}{(}\PY{l+m+mi}{8}\PY{p}{,}\PY{l+m+mi}{5}\PY{p}{)}\PY{p}{)}
\PY{n}{plt}\PY{o}{.}\PY{n}{hist}\PY{p}{(}\PY{n}{x}\PY{p}{,} \PY{n}{bins}\PY{o}{=}\PY{l+m+mi}{60}\PY{p}{,} \PY{n}{density}\PY{o}{=}\PY{k+kc}{True}\PY{p}{,} \PY{n}{alpha}\PY{o}{=}\PY{l+m+mf}{0.6}\PY{p}{)}
\PY{n}{xx} \PY{o}{=} \PY{n}{np}\PY{o}{.}\PY{n}{linspace}\PY{p}{(}\PY{n}{x}\PY{o}{.}\PY{n}{min}\PY{p}{(}\PY{p}{)}\PY{p}{,} \PY{n}{x}\PY{o}{.}\PY{n}{max}\PY{p}{(}\PY{p}{)}\PY{p}{,} \PY{l+m+mi}{400}\PY{p}{)}
\PY{n}{plt}\PY{o}{.}\PY{n}{plot}\PY{p}{(}\PY{n}{xx}\PY{p}{,} \PY{n}{norm}\PY{o}{.}\PY{n}{pdf}\PY{p}{(}\PY{n}{xx}\PY{p}{,} \PY{n}{loc}\PY{o}{=}\PY{n}{mu\PYZus{}hat}\PY{p}{,} \PY{n}{scale}\PY{o}{=}\PY{n}{sigma\PYZus{}hat}\PY{p}{)}\PY{p}{)}
\PY{n}{plt}\PY{o}{.}\PY{n}{title}\PY{p}{(}\PY{l+s+sa}{f}\PY{l+s+s2}{\PYZdq{}}\PY{l+s+si}{\PYZob{}}\PY{n}{asset}\PY{l+s+si}{\PYZcb{}}\PY{l+s+s2}{: daily log returns (USD) vs fitted normal}\PY{l+s+s2}{\PYZdq{}}\PY{p}{)}
\PY{n}{plt}\PY{o}{.}\PY{n}{xlabel}\PY{p}{(}\PY{l+s+s2}{\PYZdq{}}\PY{l+s+s2}{Daily return}\PY{l+s+s2}{\PYZdq{}}\PY{p}{)}
\PY{n}{plt}\PY{o}{.}\PY{n}{ylabel}\PY{p}{(}\PY{l+s+s2}{\PYZdq{}}\PY{l+s+s2}{Density}\PY{l+s+s2}{\PYZdq{}}\PY{p}{)}
\PY{n}{plt}\PY{o}{.}\PY{n}{tight\PYZus{}layout}\PY{p}{(}\PY{p}{)}
\PY{n}{plt}\PY{o}{.}\PY{n}{show}\PY{p}{(}\PY{p}{)}
\end{Verbatim}
\end{tcolorbox}

    \begin{center}
    \adjustimage{max size={0.9\linewidth}{0.9\paperheight}}{Work-7_files/Work-7_38_0.png}
    \end{center}
    { \hspace*{\fill} \\}
    
    On the S\&P 500, the fitted normal curve roughly matches the centre of
the histogram, but the fit is clearly not perfect. The empirical
distribution is slightly left-skewed and much more peaked with fatter
tails than the Gaussian: most observations are very close to zero, but
we also see a few extreme moves that occur more often than a normal
model would predict.

Visually this already suggests that the i.i.d.~normal assumption is
questionable, at least for the tail behaviour. To make this diagnosis
more precise, we now turn to a formal Jarque--Bera normality test on all
indices.

    \hypertarget{jarquebera-normality-test}{%
\subsubsection{1.8 Jarque--Bera normality
test}\label{jarquebera-normality-test}}

Visual inspection of the S\&P 500 returns already suggests deviations
from normality, but this is still subjective. We therefore run a
Jarque--Bera test on each index. The test statistic combines skewness
and excess kurtosis, and under the null hypothesis of normality it
follows (asymptotically) a \(\chi^2\) distribution with 2 degrees of
freedom.

In the code below we compute, for every index, the Jarque--Bera
statistic and its p-value and collect the results in a small table. This
allows us to see systematically whether the Gaussian assumption can be
rejected for our daily USD log returns.

    \begin{tcolorbox}[breakable, size=fbox, boxrule=1pt, pad at break*=1mm,colback=cellbackground, colframe=cellborder]
\begin{Verbatim}[commandchars=\\\{\}]
\PY{n}{jb} \PY{o}{=} \PY{n}{returns}\PY{o}{.}\PY{n}{apply}\PY{p}{(}\PY{k}{lambda} \PY{n}{s}\PY{p}{:} \PY{n}{jarque\PYZus{}bera}\PY{p}{(}\PY{n}{s}\PY{o}{.}\PY{n}{dropna}\PY{p}{(}\PY{p}{)}\PY{p}{)}\PY{o}{.}\PY{n}{statistic}\PY{p}{)}
\PY{n}{jb\PYZus{}p} \PY{o}{=} \PY{n}{returns}\PY{o}{.}\PY{n}{apply}\PY{p}{(}\PY{k}{lambda} \PY{n}{s}\PY{p}{:} \PY{n}{jarque\PYZus{}bera}\PY{p}{(}\PY{n}{s}\PY{o}{.}\PY{n}{dropna}\PY{p}{(}\PY{p}{)}\PY{p}{)}\PY{o}{.}\PY{n}{pvalue}\PY{p}{)}

\PY{n}{jb\PYZus{}table} \PY{o}{=} \PY{n}{pd}\PY{o}{.}\PY{n}{DataFrame}\PY{p}{(}\PY{p}{\PYZob{}}\PY{l+s+s2}{\PYZdq{}}\PY{l+s+s2}{JB\PYZus{}stat}\PY{l+s+s2}{\PYZdq{}}\PY{p}{:} \PY{n}{jb}\PY{p}{,} \PY{l+s+s2}{\PYZdq{}}\PY{l+s+s2}{p\PYZus{}value}\PY{l+s+s2}{\PYZdq{}}\PY{p}{:} \PY{n}{jb\PYZus{}p}\PY{p}{\PYZcb{}}\PY{p}{)}
\PY{n}{jb\PYZus{}table}
\end{Verbatim}
\end{tcolorbox}

            \begin{tcolorbox}[breakable, size=fbox, boxrule=.5pt, pad at break*=1mm, opacityfill=0]
\begin{Verbatim}[commandchars=\\\{\}]
                  JB\_stat  p\_value
TSX          20780.314426      0.0
CAC          41675.166886      0.0
DAX          38595.008923      0.0
Eurostoxx50  38841.018140      0.0
NIKKEI225    37864.111306      0.0
FTSE100      13757.475020      0.0
SP500        19664.141129      0.0
IBOVESPA     16524.010894      0.0
\end{Verbatim}
\end{tcolorbox}
        
    The Jarque--Bera results are completely unambiguous. For all eight
indices the JB statistics are extremely large (of order
\(10^4\)--\(10^5\)) and the p-values are effectively zero. At any usual
significance level (even 1\% or 0.1\%) we therefore \textbf{strongly
reject} the null hypothesis that daily USD log returns are normally
distributed.

Combined with the large excess kurtosis and negative skewness we found
earlier, this confirms that the Gaussian i.i.d.~assumption is not
realistic for these equity index returns: the distributions are far too
heavy-tailed and asymmetric to be well approximated by a normal law.

    \hypertarget{robustness-removing-crisis-years}{%
\subsubsection{1.9 Robustness: removing crisis
years}\label{robustness-removing-crisis-years}}

Because the strong rejection of normality could in principle be driven
by a few extreme crisis days, we run a simple robustness check. We
remove the main crisis periods (2008--2009 and 2020) from the sample,
recompute daily USD log returns on the trimmed dataset, and run the
Jarque--Bera test again for each index. This allows us to see whether
the non-Gaussianity is entirely due to tail events or remains a
structural feature of equity returns.

    \begin{tcolorbox}[breakable, size=fbox, boxrule=1pt, pad at break*=1mm,colback=cellbackground, colframe=cellborder]
\begin{Verbatim}[commandchars=\\\{\}]
\PY{c+c1}{\PYZsh{} We remove main crisis years (2008, 2009 and 2020)}
\PY{n}{crisis\PYZus{}years} \PY{o}{=} \PY{p}{[}\PY{l+m+mi}{2008}\PY{p}{,} \PY{l+m+mi}{2009}\PY{p}{,} \PY{l+m+mi}{2020}\PY{p}{]}
\PY{n}{mask\PYZus{}crisis} \PY{o}{=} \PY{n}{returns}\PY{o}{.}\PY{n}{index}\PY{o}{.}\PY{n}{year}\PY{o}{.}\PY{n}{isin}\PY{p}{(}\PY{n}{crisis\PYZus{}years}\PY{p}{)}

\PY{n}{returns\PYZus{}trim} \PY{o}{=} \PY{n}{returns}\PY{o}{.}\PY{n}{loc}\PY{p}{[}\PY{o}{\PYZti{}}\PY{n}{mask\PYZus{}crisis}\PY{p}{]}\PY{o}{.}\PY{n}{copy}\PY{p}{(}\PY{p}{)}

\PY{k+kn}{from}\PY{+w}{ }\PY{n+nn}{scipy}\PY{n+nn}{.}\PY{n+nn}{stats}\PY{+w}{ }\PY{k+kn}{import} \PY{n}{jarque\PYZus{}bera}

\PY{n}{jb\PYZus{}trim\PYZus{}stat} \PY{o}{=} \PY{n}{returns\PYZus{}trim}\PY{o}{.}\PY{n}{apply}\PY{p}{(}\PY{k}{lambda} \PY{n}{s}\PY{p}{:} \PY{n}{jarque\PYZus{}bera}\PY{p}{(}\PY{n}{s}\PY{o}{.}\PY{n}{dropna}\PY{p}{(}\PY{p}{)}\PY{p}{)}\PY{o}{.}\PY{n}{statistic}\PY{p}{)}
\PY{n}{jb\PYZus{}trim\PYZus{}p} \PY{o}{=} \PY{n}{returns\PYZus{}trim}\PY{o}{.}\PY{n}{apply}\PY{p}{(}\PY{k}{lambda} \PY{n}{s}\PY{p}{:} \PY{n}{jarque\PYZus{}bera}\PY{p}{(}\PY{n}{s}\PY{o}{.}\PY{n}{dropna}\PY{p}{(}\PY{p}{)}\PY{p}{)}\PY{o}{.}\PY{n}{pvalue}\PY{p}{)}

\PY{n}{jb\PYZus{}trim\PYZus{}table} \PY{o}{=} \PY{n}{pd}\PY{o}{.}\PY{n}{DataFrame}\PY{p}{(}
    \PY{p}{\PYZob{}}\PY{l+s+s2}{\PYZdq{}}\PY{l+s+s2}{JB\PYZus{}stat\PYZus{}trim}\PY{l+s+s2}{\PYZdq{}}\PY{p}{:} \PY{n}{jb\PYZus{}trim\PYZus{}stat}\PY{p}{,} \PY{l+s+s2}{\PYZdq{}}\PY{l+s+s2}{p\PYZus{}value\PYZus{}trim}\PY{l+s+s2}{\PYZdq{}}\PY{p}{:} \PY{n}{jb\PYZus{}trim\PYZus{}p}\PY{p}{\PYZcb{}}
\PY{p}{)}
\PY{n}{jb\PYZus{}trim\PYZus{}table}
\end{Verbatim}
\end{tcolorbox}

            \begin{tcolorbox}[breakable, size=fbox, boxrule=.5pt, pad at break*=1mm, opacityfill=0]
\begin{Verbatim}[commandchars=\\\{\}]
             JB\_stat\_trim   p\_value\_trim
TSX           1339.646168  1.257566e-291
CAC           3279.052062   0.000000e+00
DAX           1469.757549   0.000000e+00
Eurostoxx50   3921.160086   0.000000e+00
NIKKEI225     4414.006357   0.000000e+00
FTSE100       2733.185183   0.000000e+00
SP500         1806.351927   0.000000e+00
IBOVESPA       685.158441  1.658574e-149
\end{Verbatim}
\end{tcolorbox}
        
    The trimmed-sample Jarque--Bera statistics remain very large (of order
\(10^3\) for most indices, with p-values effectively equal to zero).
Even though the test statistics are lower than in the full sample, we
still strongly reject the null of normality for all markets, including
Ibovespa where the JB statistic is around 685 with a p-value of about
\(10^{-149}\). This confirms that non-Gaussian features are not driven
solely by the 2008--2009 and 2020 crisis periods, but are a persistent
property of daily equity index returns.

    \hypertarget{testing-autocorrelation-ljungbox}{%
\subsubsection{1.10 Testing autocorrelation
(Ljung--Box)}\label{testing-autocorrelation-ljungbox}}

To check whether returns are really i.i.d.~over time, we run Ljung--Box
tests on each index. For every return series we compute the test at lags
5 and 10 and store both the test statistics and the p-values in a small
summary table. This tells us whether we can reject the null hypothesis
of ``no autocorrelation up to lag 5/10'\,' for each market.

    \begin{tcolorbox}[breakable, size=fbox, boxrule=1pt, pad at break*=1mm,colback=cellbackground, colframe=cellborder]
\begin{Verbatim}[commandchars=\\\{\}]
\PY{n}{lb\PYZus{}results} \PY{o}{=} \PY{p}{\PYZob{}}\PY{p}{\PYZcb{}}
\PY{k}{for} \PY{n}{col} \PY{o+ow}{in} \PY{n}{returns}\PY{o}{.}\PY{n}{columns}\PY{p}{:}
    \PY{n}{res} \PY{o}{=} \PY{n}{acorr\PYZus{}ljungbox}\PY{p}{(}\PY{n}{returns}\PY{p}{[}\PY{n}{col}\PY{p}{]}\PY{o}{.}\PY{n}{dropna}\PY{p}{(}\PY{p}{)}\PY{p}{,} \PY{n}{lags}\PY{o}{=}\PY{p}{[}\PY{l+m+mi}{5}\PY{p}{,} \PY{l+m+mi}{10}\PY{p}{]}\PY{p}{,} \PY{n}{return\PYZus{}df}\PY{o}{=}\PY{k+kc}{True}\PY{p}{)}
    \PY{n}{lb\PYZus{}results}\PY{p}{[}\PY{n}{col}\PY{p}{]} \PY{o}{=} \PY{p}{\PYZob{}}
        \PY{l+s+s2}{\PYZdq{}}\PY{l+s+s2}{LB\PYZus{}stat\PYZus{}5}\PY{l+s+s2}{\PYZdq{}}\PY{p}{:} \PY{n}{res}\PY{p}{[}\PY{l+s+s2}{\PYZdq{}}\PY{l+s+s2}{lb\PYZus{}stat}\PY{l+s+s2}{\PYZdq{}}\PY{p}{]}\PY{o}{.}\PY{n}{iloc}\PY{p}{[}\PY{l+m+mi}{0}\PY{p}{]}\PY{p}{,}
        \PY{l+s+s2}{\PYZdq{}}\PY{l+s+s2}{p\PYZus{}value\PYZus{}5}\PY{l+s+s2}{\PYZdq{}}\PY{p}{:} \PY{n}{res}\PY{p}{[}\PY{l+s+s2}{\PYZdq{}}\PY{l+s+s2}{lb\PYZus{}pvalue}\PY{l+s+s2}{\PYZdq{}}\PY{p}{]}\PY{o}{.}\PY{n}{iloc}\PY{p}{[}\PY{l+m+mi}{0}\PY{p}{]}\PY{p}{,}
        \PY{l+s+s2}{\PYZdq{}}\PY{l+s+s2}{LB\PYZus{}stat\PYZus{}10}\PY{l+s+s2}{\PYZdq{}}\PY{p}{:} \PY{n}{res}\PY{p}{[}\PY{l+s+s2}{\PYZdq{}}\PY{l+s+s2}{lb\PYZus{}stat}\PY{l+s+s2}{\PYZdq{}}\PY{p}{]}\PY{o}{.}\PY{n}{iloc}\PY{p}{[}\PY{l+m+mi}{1}\PY{p}{]}\PY{p}{,}
        \PY{l+s+s2}{\PYZdq{}}\PY{l+s+s2}{p\PYZus{}value\PYZus{}10}\PY{l+s+s2}{\PYZdq{}}\PY{p}{:} \PY{n}{res}\PY{p}{[}\PY{l+s+s2}{\PYZdq{}}\PY{l+s+s2}{lb\PYZus{}pvalue}\PY{l+s+s2}{\PYZdq{}}\PY{p}{]}\PY{o}{.}\PY{n}{iloc}\PY{p}{[}\PY{l+m+mi}{1}\PY{p}{]}\PY{p}{,}
    \PY{p}{\PYZcb{}}

\PY{n}{lb\PYZus{}table} \PY{o}{=} \PY{n}{pd}\PY{o}{.}\PY{n}{DataFrame}\PY{p}{(}\PY{n}{lb\PYZus{}results}\PY{p}{)}\PY{o}{.}\PY{n}{T}
\PY{n}{lb\PYZus{}table}
\end{Verbatim}
\end{tcolorbox}

            \begin{tcolorbox}[breakable, size=fbox, boxrule=.5pt, pad at break*=1mm, opacityfill=0]
\begin{Verbatim}[commandchars=\\\{\}]
              LB\_stat\_5     p\_value\_5  LB\_stat\_10    p\_value\_10
TSX           17.759759  3.263065e-03   35.610376  9.818733e-05
CAC           52.736517  3.809433e-10   69.887446  4.661267e-11
DAX           35.893740  9.974462e-07   50.878012  1.839312e-07
Eurostoxx50   48.875857  2.352759e-09   62.487962  1.222083e-09
NIKKEI225    130.566808  1.804204e-26  134.983663  4.483116e-24
FTSE100       35.961387  9.668529e-07   43.837120  3.520954e-06
SP500         91.833893  2.766657e-18  112.634488  1.568458e-19
IBOVESPA       5.147162  3.981858e-01   18.314960  4.987750e-02
\end{Verbatim}
\end{tcolorbox}
        
    The Ljung--Box results point in the same direction as before: daily
returns are not i.i.d.~white noise.

\begin{itemize}
  \item For almost all indices, the p-values at lags 5 and 10 are essentially zero. We clearly
  reject the null of ``no autocorrelation up to lag 5/10''. This is especially strong for
  the Nikkei and the S\&P 500, where the test statistics are very large.
  \item Ibovespa is the only borderline case: at lag 5 the p-value is around 0.40 (we do \emph{not}
  reject), but by lag 10 it drops to about 5\%, so even there we see some dependence once
  we look further out.
\end{itemize}

Overall, the Ljung--Box tests confirm that daily USD returns display
significant serial dependence, which is another reason why the
i.i.d.~Gaussian assumption is not really appropriate for these series.

    \hypertarget{conclusion-for-question-1}{%
\subsubsection{1.11 Conclusion for Question
1}\label{conclusion-for-question-1}}

Overall, the daily USD log returns of our eight indices do not behave
like i.i.d.~Gaussian variables.

Means are tiny compared to volatilities, returns are heavily correlated
across markets, distributions are clearly skewed with very fat tails,
and both the Jarque--Bera and Ljung--Box tests strongly reject normality
and independence.

In the rest of the assignment we will still rely on the Gaussian
i.i.d.~framework to build mean--variance portfolios, but with the clear
caveat that it is a modelling simplification rather than a realistic
description of the data.

    \hypertarget{question-2-mean-variance-mv-analysis-1}{%
\subsection{Question 2 -- Mean-Variance (MV) analysis
\#1}\label{question-2-mean-variance-mv-analysis-1}}

In this question, we implement the traditional Markowitz mean--variance
framework using historical estimates of expected returns and the
covariance matrix (the ``plug-in'\,' approach from the course notes).

We:

\begin{enumerate}
    \item Estimate annualized expected returns and the annualized covariance matrix from daily log returns.
    \item Compute MV-efficient portfolios across a grid of target returns (efficient frontier).
    \item Identify the tangency portfolio (maximum Sharpe ratio).
    \item Repeat the same exercise with a no short-sale constraint (non-negative weights).
\end{enumerate}

    Since the assignment does not provide a particular value, we assume a
constant annual risk-free rate of \((r_f = 0)\) for the whole analysis.
This implies that excess returns coincide with expected returns, so the
shape of the efficient frontier and the composition of the tangency
portfolios are unaffected by choosing a small positive benchmark rate
instead.

    \begin{tcolorbox}[breakable, size=fbox, boxrule=1pt, pad at break*=1mm,colback=cellbackground, colframe=cellborder]
\begin{Verbatim}[commandchars=\\\{\}]
\PY{k+kn}{import}\PY{+w}{ }\PY{n+nn}{quadprog}
\end{Verbatim}
\end{tcolorbox}

    \hypertarget{estimate-inputs-plug-in-rule}{%
\subsubsection{2.1 Estimate inputs (plug-in
rule)}\label{estimate-inputs-plug-in-rule}}

We use historical daily log returns (in USD) to estimate:

\begin{itemize}
\item
  Annualized expected returns: \(\mu_{ann} = 252 \cdot \mu_{daily}\)
\item
  Annualized covariance matrix:
  \(\Sigma_{ann} = 252 \cdot \Sigma_{daily}\)
\end{itemize}

These are the inputs needed for mean--variance optimization.

    \begin{tcolorbox}[breakable, size=fbox, boxrule=1pt, pad at break*=1mm,colback=cellbackground, colframe=cellborder]
\begin{Verbatim}[commandchars=\\\{\}]
\PY{c+c1}{\PYZsh{} returns is our daily log returns DataFrame from Q1}
\PY{n}{assets} \PY{o}{=} \PY{n}{returns}\PY{o}{.}\PY{n}{columns}\PY{o}{.}\PY{n}{tolist}\PY{p}{(}\PY{p}{)}
\PY{n}{n} \PY{o}{=} \PY{n+nb}{len}\PY{p}{(}\PY{n}{assets}\PY{p}{)}

\PY{n}{mu\PYZus{}daily} \PY{o}{=} \PY{n}{returns}\PY{o}{.}\PY{n}{mean}\PY{p}{(}\PY{p}{)}\PY{o}{.}\PY{n}{values}                  \PY{c+c1}{\PYZsh{} shape (n,)}
\PY{n}{Sigma\PYZus{}daily} \PY{o}{=} \PY{n}{returns}\PY{o}{.}\PY{n}{cov}\PY{p}{(}\PY{p}{)}\PY{o}{.}\PY{n}{values}                \PY{c+c1}{\PYZsh{} shape (n, n)}

\PY{n}{ann\PYZus{}factor} \PY{o}{=} \PY{l+m+mi}{252}
\PY{n}{mu} \PY{o}{=} \PY{n}{mu\PYZus{}daily} \PY{o}{*} \PY{n}{ann\PYZus{}factor}                        \PY{c+c1}{\PYZsh{} annualized expected returns}
\PY{n}{Sigma} \PY{o}{=} \PY{n}{Sigma\PYZus{}daily} \PY{o}{*} \PY{n}{ann\PYZus{}factor}                  \PY{c+c1}{\PYZsh{} annualized covariance}

\PY{c+c1}{\PYZsh{} Risk\PYZhy{}free rate (not provided in the assignment, we set it to 0 and state it)}
\PY{n}{rf} \PY{o}{=} \PY{l+m+mf}{0.0}

\PY{n}{mu\PYZus{}excess} \PY{o}{=} \PY{n}{mu} \PY{o}{\PYZhy{}} \PY{n}{rf} \PY{o}{*} \PY{n}{np}\PY{o}{.}\PY{n}{ones}\PY{p}{(}\PY{n}{n}\PY{p}{)}

\PY{n}{mu\PYZus{}series} \PY{o}{=} \PY{n}{pd}\PY{o}{.}\PY{n}{Series}\PY{p}{(}\PY{n}{mu}\PY{p}{,} \PY{n}{index}\PY{o}{=}\PY{n}{assets}\PY{p}{,} \PY{n}{name}\PY{o}{=}\PY{l+s+s2}{\PYZdq{}}\PY{l+s+s2}{ann\PYZus{}mean}\PY{l+s+s2}{\PYZdq{}}\PY{p}{)}
\PY{n}{sigma\PYZus{}series} \PY{o}{=} \PY{n}{pd}\PY{o}{.}\PY{n}{Series}\PY{p}{(}\PY{n}{np}\PY{o}{.}\PY{n}{sqrt}\PY{p}{(}\PY{n}{np}\PY{o}{.}\PY{n}{diag}\PY{p}{(}\PY{n}{Sigma}\PY{p}{)}\PY{p}{)}\PY{p}{,} \PY{n}{index}\PY{o}{=}\PY{n}{assets}\PY{p}{,} \PY{n}{name}\PY{o}{=}\PY{l+s+s2}{\PYZdq{}}\PY{l+s+s2}{ann\PYZus{}vol}\PY{l+s+s2}{\PYZdq{}}\PY{p}{)}

\PY{n}{mu\PYZus{}series}\PY{p}{,} \PY{n}{sigma\PYZus{}series}
\end{Verbatim}
\end{tcolorbox}

            \begin{tcolorbox}[breakable, size=fbox, boxrule=.5pt, pad at break*=1mm, opacityfill=0]
\begin{Verbatim}[commandchars=\\\{\}]
(TSX            0.029466
 CAC            0.006178
 DAX            0.052057
 Eurostoxx50   -0.010314
 NIKKEI225      0.041593
 FTSE100       -0.016259
 SP500          0.094338
 IBOVESPA      -0.014995
 Name: ann\_mean, dtype: float64,
 TSX            0.242104
 CAC            0.291322
 DAX            0.286656
 Eurostoxx50    0.293384
 NIKKEI225      0.256576
 FTSE100        0.244062
 SP500          0.215596
 IBOVESPA       0.393585
 Name: ann\_vol, dtype: float64)
\end{Verbatim}
\end{tcolorbox}
        
    The table above reports the annualised sample means and volatilities of
daily USD log returns for the eight indices.

A first observation is that average returns are quite modest compared to
risk. Over the full 2007--2021 sample, the S\&P\textasciitilde500 stands
out with an annualised mean close to 9\%, while markets such as the DAX
and the Nikkei are in the 4--5\% range, and the TSX around 3\%. In
contrast, several indices have \emph{negative} average returns over this
period (Eurostoxx50, FTSE100, Ibovespa), which is plausible given that
our window includes the global financial crisis, the European debt
crisis and the Covid shock.

Volatilities are much larger and display strong cross-sectional
differences. The Ibovespa is by far the most volatile market, with an
annualised standard deviation close to 40\%, followed by the European
indices and the TSX, all around 25--30\%. The S\&P\textasciitilde500 is
relatively less volatile (around 21\%), but still much riskier than a
typical bond investment. Overall, this table already suggests that any
mean--variance allocation will be driven more by differences in
volatility and correlation than by small differences in average returns.

    \hypertarget{mapping-markowitz-to-quadprog}{%
\subsubsection{\texorpdfstring{2.2 Mapping Markowitz to
\texttt{quadprog}}{2.2 Mapping Markowitz to quadprog}}\label{mapping-markowitz-to-quadprog}}

The function \texttt{quadprog.solve\_qp} solves quadratic programs of
the form

\[
\min_w \; \frac{1}{2} w^\top G w - a^\top w
\quad \text{s.t.} \quad C^\top w \ge b.
\]

In the assignment, the problem is written as

\[
\min_x \; \frac{1}{2} x^\top H x + f^\top x
\quad \text{s.t.} \quad A x \le b.
\]

The two formulations are equivalent up to a change of notation and a
sign convention. We can set \[
G = H, \qquad a = -f, \qquad C = -A, \qquad b_{\text{qp}} = -b,
\] so that the constraint (A x \le b) is the same as ((-A)x \ge -b),
i.e.~\(C^\top w \ge b_{\text{qp}}\) in the \texttt{quadprog} notation.

In our mean--variance setting we take (x = w) (portfolio weights) and (H
= 2\Sigma). Then \[
\min_w \; \frac{1}{2} w^\top G w
\quad \text{with} \quad G = 2\Sigma
\] corresponds to minimising portfolio variance. For minimum-variance
portfolios we set (a = 0).

To trace the efficient frontier, we impose two equality constraints (as
in the course notes): \[
\mathbf{1}^\top w = 1
\quad \text{and} \quad
\mu^\top w = \mu_{\text{target}}.
\] In \texttt{quadprog}, these are implemented via \texttt{meq\ =\ 2},
meaning that the first two columns of (C) (and the corresponding entries
of (\(b_{\text{qp}}\)) are treated as equalities. Additional inequality
constraints, such as the no short-sale constraint \(w_i \ge 0\), can
then be added as extra columns of (C) with the appropriate entries in
\(b_{\text{qp}}\).

    \begin{tcolorbox}[breakable, size=fbox, boxrule=1pt, pad at break*=1mm,colback=cellbackground, colframe=cellborder]
\begin{Verbatim}[commandchars=\\\{\}]
\PY{k}{def}\PY{+w}{ }\PY{n+nf}{solve\PYZus{}mv\PYZus{}target\PYZus{}return\PYZus{}quadprog}\PY{p}{(}\PY{n}{mu\PYZus{}vec}\PY{p}{,} \PY{n}{Sigma\PYZus{}mat}\PY{p}{,} \PY{n}{target\PYZus{}return}\PY{p}{,} \PY{n}{no\PYZus{}short}\PY{o}{=}\PY{k+kc}{False}\PY{p}{)}\PY{p}{:}
\PY{+w}{    }\PY{l+s+sd}{\PYZdq{}\PYZdq{}\PYZdq{}}
\PY{l+s+sd}{    Minimize variance subject to:}
\PY{l+s+sd}{      1\PYZsq{} w = 1}
\PY{l+s+sd}{      mu\PYZsq{} w = target\PYZus{}return}
\PY{l+s+sd}{    Optionally: w \PYZgt{}= 0 (no\PYZus{}short=True)}
\PY{l+s+sd}{    Uses quadprog: min 1/2 w\PYZsq{} G w \PYZhy{} a\PYZsq{} w  s.t.  C\PYZsq{} w \PYZgt{}= b}
\PY{l+s+sd}{    \PYZdq{}\PYZdq{}\PYZdq{}}
    \PY{n}{n} \PY{o}{=} \PY{n+nb}{len}\PY{p}{(}\PY{n}{mu\PYZus{}vec}\PY{p}{)}
    \PY{n}{G} \PY{o}{=} \PY{l+m+mf}{2.0} \PY{o}{*} \PY{n}{Sigma\PYZus{}mat}
    \PY{n}{a} \PY{o}{=} \PY{n}{np}\PY{o}{.}\PY{n}{zeros}\PY{p}{(}\PY{n}{n}\PY{p}{)}

    \PY{n}{ones} \PY{o}{=} \PY{n}{np}\PY{o}{.}\PY{n}{ones}\PY{p}{(}\PY{n}{n}\PY{p}{)}

    \PY{c+c1}{\PYZsh{} Build constraints C\PYZsq{} w gt or equal to b}
    \PY{c+c1}{\PYZsh{} First two constraints are equalities: mu\PYZsq{}w = target, 1\PYZsq{}w = 1}
    \PY{k}{if} \PY{n}{no\PYZus{}short}\PY{p}{:}
        \PY{n}{C} \PY{o}{=} \PY{n}{np}\PY{o}{.}\PY{n}{column\PYZus{}stack}\PY{p}{(}\PY{p}{[}\PY{n}{mu\PYZus{}vec}\PY{p}{,} \PY{n}{ones}\PY{p}{,} \PY{n}{np}\PY{o}{.}\PY{n}{eye}\PY{p}{(}\PY{n}{n}\PY{p}{)}\PY{p}{]}\PY{p}{)}     \PY{c+c1}{\PYZsh{} shape (n, 2+n)}
        \PY{n}{b} \PY{o}{=} \PY{n}{np}\PY{o}{.}\PY{n}{concatenate}\PY{p}{(}\PY{p}{[}\PY{p}{[}\PY{n}{target\PYZus{}return}\PY{p}{,} \PY{l+m+mf}{1.0}\PY{p}{]}\PY{p}{,} \PY{n}{np}\PY{o}{.}\PY{n}{zeros}\PY{p}{(}\PY{n}{n}\PY{p}{)}\PY{p}{]}\PY{p}{)}
        \PY{n}{meq} \PY{o}{=} \PY{l+m+mi}{2}
    \PY{k}{else}\PY{p}{:}
        \PY{n}{C} \PY{o}{=} \PY{n}{np}\PY{o}{.}\PY{n}{column\PYZus{}stack}\PY{p}{(}\PY{p}{[}\PY{n}{mu\PYZus{}vec}\PY{p}{,} \PY{n}{ones}\PY{p}{]}\PY{p}{)}                \PY{c+c1}{\PYZsh{} shape (n, 2)}
        \PY{n}{b} \PY{o}{=} \PY{n}{np}\PY{o}{.}\PY{n}{array}\PY{p}{(}\PY{p}{[}\PY{n}{target\PYZus{}return}\PY{p}{,} \PY{l+m+mf}{1.0}\PY{p}{]}\PY{p}{)}
        \PY{n}{meq} \PY{o}{=} \PY{l+m+mi}{2}

    \PY{n}{sol} \PY{o}{=} \PY{n}{quadprog}\PY{o}{.}\PY{n}{solve\PYZus{}qp}\PY{p}{(}\PY{n}{G}\PY{p}{,} \PY{n}{a}\PY{p}{,} \PY{n}{C}\PY{p}{,} \PY{n}{b}\PY{p}{,} \PY{n}{meq}\PY{p}{)}
    \PY{n}{w} \PY{o}{=} \PY{n}{sol}\PY{p}{[}\PY{l+m+mi}{0}\PY{p}{]}
    \PY{k}{return} \PY{n}{w}


\PY{k}{def}\PY{+w}{ }\PY{n+nf}{portfolio\PYZus{}stats}\PY{p}{(}\PY{n}{w}\PY{p}{,} \PY{n}{mu\PYZus{}vec}\PY{p}{,} \PY{n}{Sigma\PYZus{}mat}\PY{p}{,} \PY{n}{rf}\PY{o}{=}\PY{l+m+mf}{0.0}\PY{p}{)}\PY{p}{:}
\PY{+w}{    }\PY{l+s+sd}{\PYZdq{}\PYZdq{}\PYZdq{}}
\PY{l+s+sd}{    Return (ann\PYZus{}return, ann\PYZus{}vol, ann\PYZus{}sharpe)}
\PY{l+s+sd}{    \PYZdq{}\PYZdq{}\PYZdq{}}
    \PY{n}{ret} \PY{o}{=} \PY{n+nb}{float}\PY{p}{(}\PY{n}{w} \PY{o}{@} \PY{n}{mu\PYZus{}vec}\PY{p}{)}
    \PY{n}{vol} \PY{o}{=} \PY{n+nb}{float}\PY{p}{(}\PY{n}{np}\PY{o}{.}\PY{n}{sqrt}\PY{p}{(}\PY{n}{w} \PY{o}{@} \PY{n}{Sigma\PYZus{}mat} \PY{o}{@} \PY{n}{w}\PY{p}{)}\PY{p}{)}
    \PY{n}{sharpe} \PY{o}{=} \PY{p}{(}\PY{n}{ret} \PY{o}{\PYZhy{}} \PY{n}{rf}\PY{p}{)} \PY{o}{/} \PY{n}{vol} \PY{k}{if} \PY{n}{vol} \PY{o}{\PYZgt{}} \PY{l+m+mi}{0} \PY{k}{else} \PY{n}{np}\PY{o}{.}\PY{n}{nan}
    \PY{k}{return} \PY{n}{ret}\PY{p}{,} \PY{n}{vol}\PY{p}{,} \PY{n}{sharpe}
\end{Verbatim}
\end{tcolorbox}

    The two helper functions defined above will be used throughout the
mean--variance analysis.

\begin{itemize}
\item
  \texttt{solve\_mv\_target\_return\_quadprog} solves the Markowitz
  problem of \textbf{minimum variance for a given target return}.\\
  Given annualised expected returns \(\mu\), covariance matrix
  \(\Sigma\) and a target return \(\mu_{\text{target}}\), it solves \[
  \min_{w} \; \frac{1}{2} w^\top (2\Sigma) w
  \quad \text{s.t.} \quad
  \mathbf{1}^\top w = 1, \quad
  \mu^\top w = \mu_{\text{target}},
  \] and, when \texttt{no\_short=True}, it additionally imposes \[
  w_i \ge 0 \quad \text{for all } i.
  \] Internally, the function maps this problem to the
  \texttt{quadprog.solve\_qp} format
  \(\min_w \tfrac12 w^\top G w - a^\top w \;\text{s.t.}\; C^\top w \ge b\)
  and returns the optimal weights \(w^\star\).
\item
  \texttt{portfolio\_stats} is a small helper that takes a weight vector
  \(w\) and computes its annualised expected return, volatility and
  Sharpe ratio (with risk-free rate \(r_f\)): \[
  R_p = w^\top \mu, \qquad
  \sigma_p = \sqrt{w^\top \Sigma w}, \qquad
  \text{Sharpe} = \frac{R_p - r_f}{\sigma_p}.
  \]
\end{itemize}

These two routines will be used next to construct the mean-variance
efficient frontier and to identify the tangency portfolio.

    \hypertarget{efficient-frontier-unconstrained}{%
\subsubsection{2.3 Efficient frontier
(unconstrained)}\label{efficient-frontier-unconstrained}}

We compute the mean--variance efficient frontier by solving a sequence
of minimum-variance problems for a grid of target returns. Using the
plug-in estimates \(\mu\) and \(\Sigma\) from Section 2.1, we now build
the mean--variance efficient frontier without any constraints on
short-selling.

First, we create a grid of target annualized returns between the minimum
and maximum sample means of the assets. For each target return in this
grid, we call \texttt{solve\_mv\_target\_return\_quadprog} to find the
portfolio with the lowest variance that

\begin{itemize}
\tightlist
\item
  is fully invested (\(1^\top w = 1\)), and\\
\item
  has expected return equal to the chosen target.
\end{itemize}

For every solution we record the portfolio's annualized return,
volatility, and Sharpe ratio, which gives us the set of points used to
trace the unconstrained efficient frontier.

In addition to the frontier, we compute the \textbf{tangency portfolio}
using the closed-form formula from the course notes, \[
w_T = \frac{\Sigma^{-1} \mu_{\text{excess}}}{\mathbf{1}^\top \Sigma^{-1} \mu_{\text{excess}}},
\] where \(\mu_{\text{excess}}\) denotes the vector of excess expected
returns. Since we set \(r_f = 0\) in this assignment, we have
\(\mu_{\text{excess}} = \mu\). This gives us the portfolio with the
maximum Sharpe ratio when short-selling is allowed.

    \begin{tcolorbox}[breakable, size=fbox, boxrule=1pt, pad at break*=1mm,colback=cellbackground, colframe=cellborder]
\begin{Verbatim}[commandchars=\\\{\}]
\PY{c+c1}{\PYZsh{} Grid of target returns (annualized)}
\PY{c+c1}{\PYZsh{} We take a reasonable range based on asset annualized means}
\PY{n}{target\PYZus{}grid} \PY{o}{=} \PY{n}{np}\PY{o}{.}\PY{n}{linspace}\PY{p}{(}\PY{n}{mu}\PY{o}{.}\PY{n}{min}\PY{p}{(}\PY{p}{)}\PY{p}{,} \PY{n}{mu}\PY{o}{.}\PY{n}{max}\PY{p}{(}\PY{p}{)}\PY{p}{,} \PY{l+m+mi}{60}\PY{p}{)}

\PY{n}{frontier\PYZus{}uncon} \PY{o}{=} \PY{p}{[}\PY{p}{]}
\PY{n}{weights\PYZus{}uncon} \PY{o}{=} \PY{p}{[}\PY{p}{]}

\PY{k}{for} \PY{n}{t\PYZus{}ret} \PY{o+ow}{in} \PY{n}{target\PYZus{}grid}\PY{p}{:}
    \PY{n}{w} \PY{o}{=} \PY{n}{solve\PYZus{}mv\PYZus{}target\PYZus{}return\PYZus{}quadprog}\PY{p}{(}\PY{n}{mu}\PY{p}{,} \PY{n}{Sigma}\PY{p}{,} \PY{n}{t\PYZus{}ret}\PY{p}{,} \PY{n}{no\PYZus{}short}\PY{o}{=}\PY{k+kc}{False}\PY{p}{)}
    \PY{n}{ret}\PY{p}{,} \PY{n}{vol}\PY{p}{,} \PY{n}{sh} \PY{o}{=} \PY{n}{portfolio\PYZus{}stats}\PY{p}{(}\PY{n}{w}\PY{p}{,} \PY{n}{mu}\PY{p}{,} \PY{n}{Sigma}\PY{p}{,} \PY{n}{rf}\PY{o}{=}\PY{n}{rf}\PY{p}{)}
    \PY{n}{frontier\PYZus{}uncon}\PY{o}{.}\PY{n}{append}\PY{p}{(}\PY{p}{(}\PY{n}{ret}\PY{p}{,} \PY{n}{vol}\PY{p}{,} \PY{n}{sh}\PY{p}{)}\PY{p}{)}
    \PY{n}{weights\PYZus{}uncon}\PY{o}{.}\PY{n}{append}\PY{p}{(}\PY{n}{w}\PY{p}{)}

\PY{n}{frontier\PYZus{}uncon} \PY{o}{=} \PY{n}{np}\PY{o}{.}\PY{n}{array}\PY{p}{(}\PY{n}{frontier\PYZus{}uncon}\PY{p}{)}  \PY{c+c1}{\PYZsh{} columns: ret, vol, sharpe}

\PY{c+c1}{\PYZsh{} Tangency portfolio (based on course notes)}
\PY{c+c1}{\PYZsh{} w\PYZus{}T = Σ\PYZca{}\PYZob{}\PYZhy{}1\PYZcb{} μ\PYZus{}excess / (1\PYZsq{} Σ\PYZca{}\PYZob{}\PYZhy{}1\PYZcb{} μ\PYZus{}excess)}

\PY{c+c1}{\PYZsh{} Solve Σ x = μ\PYZus{}excess instead of explicitly inverting Σ}
\PY{n}{numer} \PY{o}{=} \PY{n}{np}\PY{o}{.}\PY{n}{linalg}\PY{o}{.}\PY{n}{solve}\PY{p}{(}\PY{n}{Sigma}\PY{p}{,} \PY{n}{mu\PYZus{}excess}\PY{p}{)}     \PY{c+c1}{\PYZsh{} = Σ\PYZca{}\PYZob{}\PYZhy{}1\PYZcb{} μ\PYZus{}excess}
\PY{n}{denom} \PY{o}{=} \PY{n}{np}\PY{o}{.}\PY{n}{ones}\PY{p}{(}\PY{n}{n}\PY{p}{)} \PY{o}{@} \PY{n}{numer}
\PY{n}{w\PYZus{}tan\PYZus{}uncon} \PY{o}{=} \PY{n}{numer} \PY{o}{/} \PY{n}{denom}

\PY{n}{tan\PYZus{}ret\PYZus{}uncon}\PY{p}{,} \PY{n}{tan\PYZus{}vol\PYZus{}uncon}\PY{p}{,} \PY{n}{tan\PYZus{}sh\PYZus{}uncon} \PY{o}{=} \PY{n}{portfolio\PYZus{}stats}\PY{p}{(}\PY{n}{w\PYZus{}tan\PYZus{}uncon}\PY{p}{,} \PY{n}{mu}\PY{p}{,} \PY{n}{Sigma}\PY{p}{,} \PY{n}{rf}\PY{o}{=}\PY{n}{rf}\PY{p}{)}

\PY{n}{w\PYZus{}tan\PYZus{}uncon}\PY{p}{,} \PY{n}{tan\PYZus{}ret\PYZus{}uncon}\PY{p}{,} \PY{n}{tan\PYZus{}vol\PYZus{}uncon}\PY{p}{,} \PY{n}{tan\PYZus{}sh\PYZus{}uncon}
\end{Verbatim}
\end{tcolorbox}

            \begin{tcolorbox}[breakable, size=fbox, boxrule=.5pt, pad at break*=1mm, opacityfill=0]
\begin{Verbatim}[commandchars=\\\{\}]
(array([-0.48237118,  3.07841857,  4.02783384, -6.68241971,  0.44634165,
        -0.8347226 ,  1.73123959, -0.28432016]),
 0.48312455032742646,
 0.42391369373721177,
 1.139676678212995)
\end{Verbatim}
\end{tcolorbox}
        
    The tangency portfolio without short-sale constraints has the following
properties:

\begin{itemize}
\item
  The weight vector \(w_T\) satisfies \(\sum_i w_{T,i} = 1\), but
  several positions are very large in absolute value (both long and
  short).\\
  Numerically, the leverage is high, with
  \(\sum_i |w_{T,i}| \approx 17.6\), so this is a strongly levered
  long--short portfolio.
\item
  Annualised expected return: \(R_T \approx 48.3\%\).
\item
  Annualised volatility: \(\sigma_T \approx 42.4\%\).
\item
  Sharpe ratio (with \(r_f = 0\)): \(\text{Sharpe}_T \approx 1.14\).
\end{itemize}

These results are consistent with the Markowitz theory: once
short-selling is allowed without additional constraints, the tangency
portfolio exploits small differences in risk-adjusted performance across
indices, which leads to very concentrated and highly levered positions.
In practice, such a portfolio would be hard to implement for an investor
facing leverage or risk limits, which motivates repeating the exercise
with a non-negativity constraint on the weights.

    \hypertarget{plot-assets-efficient-frontier-and-tangency-portfolio-unconstrained}{%
\subsubsection{2.4 Plot: assets, efficient frontier, and tangency
portfolio
(unconstrained)}\label{plot-assets-efficient-frontier-and-tangency-portfolio-unconstrained}}

We now visualise the results in annualised units (assuming \(r_f = 0\)):

\begin{itemize}
\tightlist
\item
  Individual assets, plotted as points in the \((\sigma, \mu)\) plane,
  where \(\sigma\) is annualised volatility and \(\mu\) is the
  annualised expected return.\\
\item
  The mean--variance efficient frontier, obtained by solving the
  minimum-variance problem for a grid of target returns
  \(\mu_{\text{target}}\).\\
\item
  The tangency portfolio, located at \((\sigma_T, R_T)\), which
  corresponds to the maximum Sharpe ratio
  \(\text{Sharpe}_T = \dfrac{R_T - r_f}{\sigma_T}\) under unrestricted
  short-selling.
\end{itemize}

This plot allows us to compare the historical risk--return profiles of
the individual indices with the set of portfolios that are optimal in
the Markowitz sense, and to highlight where the tangency portfolio lies
on the unconstrained efficient frontier.

    \begin{tcolorbox}[breakable, size=fbox, boxrule=1pt, pad at break*=1mm,colback=cellbackground, colframe=cellborder]
\begin{Verbatim}[commandchars=\\\{\}]
\PY{c+c1}{\PYZsh{} 2.3 Efficient frontier (unconstrained) – recomputed up to tangency return}

\PY{c+c1}{\PYZsh{} Grid of target returns (annualised)}
\PY{c+c1}{\PYZsh{} We go from the minimum sample mean up to slightly above the tangency return}
\PY{n}{target\PYZus{}grid} \PY{o}{=} \PY{n}{np}\PY{o}{.}\PY{n}{linspace}\PY{p}{(}\PY{n}{mu}\PY{o}{.}\PY{n}{min}\PY{p}{(}\PY{p}{)}\PY{p}{,} \PY{n}{tan\PYZus{}ret\PYZus{}uncon} \PY{o}{*} \PY{l+m+mf}{1.1}\PY{p}{,} \PY{l+m+mi}{80}\PY{p}{)}

\PY{n}{frontier\PYZus{}uncon} \PY{o}{=} \PY{p}{[}\PY{p}{]}
\PY{n}{weights\PYZus{}uncon} \PY{o}{=} \PY{p}{[}\PY{p}{]}

\PY{k}{for} \PY{n}{t\PYZus{}ret} \PY{o+ow}{in} \PY{n}{target\PYZus{}grid}\PY{p}{:}
    \PY{n}{w} \PY{o}{=} \PY{n}{solve\PYZus{}mv\PYZus{}target\PYZus{}return\PYZus{}quadprog}\PY{p}{(}\PY{n}{mu}\PY{p}{,} \PY{n}{Sigma}\PY{p}{,} \PY{n}{t\PYZus{}ret}\PY{p}{,} \PY{n}{no\PYZus{}short}\PY{o}{=}\PY{k+kc}{False}\PY{p}{)}
    \PY{n}{ret}\PY{p}{,} \PY{n}{vol}\PY{p}{,} \PY{n}{sh} \PY{o}{=} \PY{n}{portfolio\PYZus{}stats}\PY{p}{(}\PY{n}{w}\PY{p}{,} \PY{n}{mu}\PY{p}{,} \PY{n}{Sigma}\PY{p}{,} \PY{n}{rf}\PY{o}{=}\PY{n}{rf}\PY{p}{)}
    \PY{n}{frontier\PYZus{}uncon}\PY{o}{.}\PY{n}{append}\PY{p}{(}\PY{p}{(}\PY{n}{ret}\PY{p}{,} \PY{n}{vol}\PY{p}{,} \PY{n}{sh}\PY{p}{)}\PY{p}{)}
    \PY{n}{weights\PYZus{}uncon}\PY{o}{.}\PY{n}{append}\PY{p}{(}\PY{n}{w}\PY{p}{)}

\PY{n}{frontier\PYZus{}uncon} \PY{o}{=} \PY{n}{np}\PY{o}{.}\PY{n}{array}\PY{p}{(}\PY{n}{frontier\PYZus{}uncon}\PY{p}{)}  \PY{c+c1}{\PYZsh{} columns: [ret, vol, sharpe]}
\end{Verbatim}
\end{tcolorbox}

    \begin{tcolorbox}[breakable, size=fbox, boxrule=1pt, pad at break*=1mm,colback=cellbackground, colframe=cellborder]
\begin{Verbatim}[commandchars=\\\{\}]
\PY{n}{plt}\PY{o}{.}\PY{n}{figure}\PY{p}{(}\PY{n}{figsize}\PY{o}{=}\PY{p}{(}\PY{l+m+mi}{9}\PY{p}{,} \PY{l+m+mi}{6}\PY{p}{)}\PY{p}{)}

\PY{c+c1}{\PYZsh{} 1) Individual assets}
\PY{n}{plt}\PY{o}{.}\PY{n}{scatter}\PY{p}{(}\PY{n}{sigma\PYZus{}series}\PY{o}{.}\PY{n}{values}\PY{p}{,} \PY{n}{mu\PYZus{}series}\PY{o}{.}\PY{n}{values}\PY{p}{,} \PY{n}{label}\PY{o}{=}\PY{l+s+s2}{\PYZdq{}}\PY{l+s+s2}{Assets}\PY{l+s+s2}{\PYZdq{}}\PY{p}{)}
\PY{k}{for} \PY{n}{i}\PY{p}{,} \PY{n}{name} \PY{o+ow}{in} \PY{n+nb}{enumerate}\PY{p}{(}\PY{n}{assets}\PY{p}{)}\PY{p}{:}
    \PY{n}{plt}\PY{o}{.}\PY{n}{annotate}\PY{p}{(}\PY{n}{name}\PY{p}{,}
                 \PY{p}{(}\PY{n}{sigma\PYZus{}series}\PY{o}{.}\PY{n}{values}\PY{p}{[}\PY{n}{i}\PY{p}{]}\PY{p}{,} \PY{n}{mu\PYZus{}series}\PY{o}{.}\PY{n}{values}\PY{p}{[}\PY{n}{i}\PY{p}{]}\PY{p}{)}\PY{p}{,}
                 \PY{n}{fontsize}\PY{o}{=}\PY{l+m+mi}{9}\PY{p}{)}

\PY{c+c1}{\PYZsh{} 2) Unconstrained efficient frontier}
\PY{n}{plt}\PY{o}{.}\PY{n}{plot}\PY{p}{(}\PY{n}{frontier\PYZus{}uncon}\PY{p}{[}\PY{p}{:}\PY{p}{,} \PY{l+m+mi}{1}\PY{p}{]}\PY{p}{,} \PY{n}{frontier\PYZus{}uncon}\PY{p}{[}\PY{p}{:}\PY{p}{,} \PY{l+m+mi}{0}\PY{p}{]}\PY{p}{,}
         \PY{n}{label}\PY{o}{=}\PY{l+s+s2}{\PYZdq{}}\PY{l+s+s2}{Efficient frontier (unconstrained)}\PY{l+s+s2}{\PYZdq{}}\PY{p}{)}

\PY{c+c1}{\PYZsh{} 3) Tangency portfolio}
\PY{n}{plt}\PY{o}{.}\PY{n}{scatter}\PY{p}{(}\PY{p}{[}\PY{n}{tan\PYZus{}vol\PYZus{}uncon}\PY{p}{]}\PY{p}{,} \PY{p}{[}\PY{n}{tan\PYZus{}ret\PYZus{}uncon}\PY{p}{]}\PY{p}{,}
            \PY{n}{marker}\PY{o}{=}\PY{l+s+s2}{\PYZdq{}}\PY{l+s+s2}{*}\PY{l+s+s2}{\PYZdq{}}\PY{p}{,} \PY{n}{s}\PY{o}{=}\PY{l+m+mi}{200}\PY{p}{,} \PY{n}{label}\PY{o}{=}\PY{l+s+s2}{\PYZdq{}}\PY{l+s+s2}{Tangency (unconstrained)}\PY{l+s+s2}{\PYZdq{}}\PY{p}{)}

\PY{c+c1}{\PYZsh{} 4) Capital Market Line (rf = 0 =\PYZgt{} line from origin to tangency)}
\PY{n}{vol\PYZus{}line} \PY{o}{=} \PY{n}{np}\PY{o}{.}\PY{n}{linspace}\PY{p}{(}\PY{l+m+mi}{0}\PY{p}{,} \PY{n}{tan\PYZus{}vol\PYZus{}uncon} \PY{o}{*} \PY{l+m+mf}{1.05}\PY{p}{,} \PY{l+m+mi}{50}\PY{p}{)}
\PY{n}{ret\PYZus{}line} \PY{o}{=} \PY{p}{(}\PY{n}{tan\PYZus{}ret\PYZus{}uncon} \PY{o}{/} \PY{n}{tan\PYZus{}vol\PYZus{}uncon}\PY{p}{)} \PY{o}{*} \PY{n}{vol\PYZus{}line}  \PY{c+c1}{\PYZsh{} slope = Sharpe\PYZus{}T}
\PY{n}{plt}\PY{o}{.}\PY{n}{plot}\PY{p}{(}\PY{n}{vol\PYZus{}line}\PY{p}{,} \PY{n}{ret\PYZus{}line}\PY{p}{,} \PY{n}{linestyle}\PY{o}{=}\PY{l+s+s2}{\PYZdq{}}\PY{l+s+s2}{\PYZhy{}\PYZhy{}}\PY{l+s+s2}{\PYZdq{}}\PY{p}{,} \PY{n}{label}\PY{o}{=}\PY{l+s+s2}{\PYZdq{}}\PY{l+s+s2}{Capital Market Line}\PY{l+s+s2}{\PYZdq{}}\PY{p}{)}

\PY{n}{plt}\PY{o}{.}\PY{n}{title}\PY{p}{(}\PY{l+s+s2}{\PYZdq{}}\PY{l+s+s2}{Mean\PYZhy{}Variance Efficient Frontier (Unconstrained)}\PY{l+s+s2}{\PYZdq{}}\PY{p}{)}
\PY{n}{plt}\PY{o}{.}\PY{n}{xlabel}\PY{p}{(}\PY{l+s+s2}{\PYZdq{}}\PY{l+s+s2}{Annualized volatility}\PY{l+s+s2}{\PYZdq{}}\PY{p}{)}
\PY{n}{plt}\PY{o}{.}\PY{n}{ylabel}\PY{p}{(}\PY{l+s+s2}{\PYZdq{}}\PY{l+s+s2}{Annualized expected return}\PY{l+s+s2}{\PYZdq{}}\PY{p}{)}
\PY{n}{plt}\PY{o}{.}\PY{n}{legend}\PY{p}{(}\PY{p}{)}
\PY{n}{plt}\PY{o}{.}\PY{n}{tight\PYZus{}layout}\PY{p}{(}\PY{p}{)}
\PY{n}{plt}\PY{o}{.}\PY{n}{show}\PY{p}{(}\PY{p}{)}
\end{Verbatim}
\end{tcolorbox}

    \begin{center}
    \adjustimage{max size={0.9\linewidth}{0.9\paperheight}}{Work-7_files/Work-7_64_0.png}
    \end{center}
    { \hspace*{\fill} \\}
    
    The plot shows the \textbf{unconstrained Markowitz world} in one
picture.

The blue dots are the individual indices \((\sigma_i, \mu_i)\). Most of
them sit clearly \textbf{below} the blue efficient frontier: for the
same volatility, a suitably diversified portfolio can deliver a higher
expected return than any single index.

The blue curve is the \textbf{mean--variance efficient frontier}
obtained from our quadratic program with\\
\(1^\top w = 1\) and \(\mu^\top w = \mu_{\text{target}}\). The orange
star is the \textbf{tangency portfolio}: it lies on the frontier and
achieves the highest Sharpe ratio. Its location, far above and to the
right of the indices, is consistent with the very levered long--short
weights we obtained numerically.

Finally, the dashed orange line is the \textbf{Capital Market Line}.
With \(r_f = 0\), it runs from the origin through the tangency
portfolio. Any combination of the risk-free asset and the tangency
portfolio lies on this line and dominates the rest of the frontier in
terms of expected return per unit of risk.

    \hypertarget{efficient-frontier-with-no-short-sale-constraint-w-ge-0}{%
\subsubsection{\texorpdfstring{2.5 Efficient frontier with no short-sale
constraint
(\(w \ge 0\))}{2.5 Efficient frontier with no short-sale constraint (w \textbackslash ge 0)}}\label{efficient-frontier-with-no-short-sale-constraint-w-ge-0}}

We now recompute the mean--variance efficient frontier \textbf{imposing
non-negative weights}, i.e. \(w_i \ge 0\) for all assets. The
optimization problem becomes:

\begin{itemize}
\tightlist
\item
  fully invested: \(1^\top w = 1\)\\
\item
  target return: \(\mu^\top w = \mu_{\text{target}}\)\\
\item
  no short-selling: \(w_i \ge 0\) for all \(i\).
\end{itemize}

For each target return on the same grid as before, we solve the
constrained minimum-variance problem and record the corresponding
portfolio return, volatility and Sharpe ratio.

The \textbf{constrained tangency portfolio} is then defined as the
portfolio on this no-short-sale frontier that attains the
\textbf{highest Sharpe ratio} (with \(r_f = 0\)).

    Because we impose a no short-sale constraint (\(w \ge 0\)), not every
target return on the grid is necessarily feasible. In practice, for some
high \(\mu_{\text{target}}\) the optimizer cannot find a portfolio that
satisfies simultaneously: - \(1^\top w = 1\), -
\(\mu^\top w = \mu_{\text{target}}\), - and \(w_i \ge 0\) for all \(i\).

We therefore solve the optimization over the whole grid of target
returns and retain only the portfolios for which the constraints admit a
valid solution.

    \begin{tcolorbox}[breakable, size=fbox, boxrule=1pt, pad at break*=1mm,colback=cellbackground, colframe=cellborder]
\begin{Verbatim}[commandchars=\\\{\}]
\PY{c+c1}{\PYZsh{} Efficient frontier with no short\PYZhy{}selling (w \PYZgt{}= 0)}
\PY{c+c1}{\PYZsh{} We keep only feasible targets, and define the constrained tangency}
\PY{c+c1}{\PYZsh{} as the portfolio on this frontier with the highest Sharpe ratio.}
\PY{n}{frontier\PYZus{}noshort} \PY{o}{=} \PY{p}{[}\PY{p}{]}
\PY{n}{weights\PYZus{}noshort} \PY{o}{=} \PY{p}{[}\PY{p}{]}
\PY{n}{targets\PYZus{}ok} \PY{o}{=} \PY{p}{[}\PY{p}{]}

\PY{k}{for} \PY{n}{t\PYZus{}ret} \PY{o+ow}{in} \PY{n}{target\PYZus{}grid}\PY{p}{:}
    \PY{k}{try}\PY{p}{:}
        \PY{n}{w} \PY{o}{=} \PY{n}{solve\PYZus{}mv\PYZus{}target\PYZus{}return\PYZus{}quadprog}\PY{p}{(}\PY{n}{mu}\PY{p}{,} \PY{n}{Sigma}\PY{p}{,} \PY{n}{t\PYZus{}ret}\PY{p}{,} \PY{n}{no\PYZus{}short}\PY{o}{=}\PY{k+kc}{True}\PY{p}{)}
        \PY{n}{ret}\PY{p}{,} \PY{n}{vol}\PY{p}{,} \PY{n}{sh} \PY{o}{=} \PY{n}{portfolio\PYZus{}stats}\PY{p}{(}\PY{n}{w}\PY{p}{,} \PY{n}{mu}\PY{p}{,} \PY{n}{Sigma}\PY{p}{,} \PY{n}{rf}\PY{o}{=}\PY{n}{rf}\PY{p}{)}
        \PY{n}{frontier\PYZus{}noshort}\PY{o}{.}\PY{n}{append}\PY{p}{(}\PY{p}{(}\PY{n}{ret}\PY{p}{,} \PY{n}{vol}\PY{p}{,} \PY{n}{sh}\PY{p}{)}\PY{p}{)}
        \PY{n}{weights\PYZus{}noshort}\PY{o}{.}\PY{n}{append}\PY{p}{(}\PY{n}{w}\PY{p}{)}
        \PY{n}{targets\PYZus{}ok}\PY{o}{.}\PY{n}{append}\PY{p}{(}\PY{n}{t\PYZus{}ret}\PY{p}{)}
    \PY{k}{except} \PY{n+ne}{ValueError}\PY{p}{:}
        \PY{c+c1}{\PYZsh{} infeasible target under no\PYZhy{}short constraints \PYZhy{}\PYZgt{} skip}
        \PY{k}{continue}

\PY{n}{frontier\PYZus{}noshort} \PY{o}{=} \PY{n}{np}\PY{o}{.}\PY{n}{array}\PY{p}{(}\PY{n}{frontier\PYZus{}noshort}\PY{p}{)}
\PY{n}{targets\PYZus{}ok} \PY{o}{=} \PY{n}{np}\PY{o}{.}\PY{n}{array}\PY{p}{(}\PY{n}{targets\PYZus{}ok}\PY{p}{)}

\PY{c+c1}{\PYZsh{} Tangency under no\PYZhy{}short: pick the feasible portfolio with the highest Sharpe ratio}
\PY{n}{idx\PYZus{}best} \PY{o}{=} \PY{n}{np}\PY{o}{.}\PY{n}{nanargmax}\PY{p}{(}\PY{n}{frontier\PYZus{}noshort}\PY{p}{[}\PY{p}{:}\PY{p}{,} \PY{l+m+mi}{2}\PY{p}{]}\PY{p}{)}
\PY{n}{w\PYZus{}tan\PYZus{}noshort} \PY{o}{=} \PY{n}{weights\PYZus{}noshort}\PY{p}{[}\PY{n}{idx\PYZus{}best}\PY{p}{]}
\PY{n}{tan\PYZus{}ret\PYZus{}ns}\PY{p}{,} \PY{n}{tan\PYZus{}vol\PYZus{}ns}\PY{p}{,} \PY{n}{tan\PYZus{}sh\PYZus{}ns} \PY{o}{=} \PY{n}{portfolio\PYZus{}stats}\PY{p}{(}\PY{n}{w\PYZus{}tan\PYZus{}noshort}\PY{p}{,} \PY{n}{mu}\PY{p}{,} \PY{n}{Sigma}\PY{p}{,} \PY{n}{rf}\PY{o}{=}\PY{n}{rf}\PY{p}{)}

\PY{n}{w\PYZus{}tan\PYZus{}noshort}\PY{p}{,} \PY{n}{tan\PYZus{}ret\PYZus{}ns}\PY{p}{,} \PY{n}{tan\PYZus{}vol\PYZus{}ns}\PY{p}{,} \PY{n}{tan\PYZus{}sh\PYZus{}ns}
\end{Verbatim}
\end{tcolorbox}

            \begin{tcolorbox}[breakable, size=fbox, boxrule=.5pt, pad at break*=1mm, opacityfill=0]
\begin{Verbatim}[commandchars=\\\{\}]
(array([ 5.30900506e-17, -2.55695953e-16, -1.53780126e-18,  9.62876608e-17,
         1.25196383e-01, -1.00424866e-16,  8.74803617e-01,  0.00000000e+00]),
 0.08773417759756252,
 0.19443955322941314,
 0.45121569217991214)
\end{Verbatim}
\end{tcolorbox}
        
    The tangency portfolio \textbf{with no short-selling} has the following
properties:

\begin{itemize}
\item
  The weight vector is essentially long-only in two indices.\\
  Numerically, we obtain\\
  \[
  w_T \approx 
  (0,\ 0,\ 0,\ 0,\ 0.125,\ 0,\ 0.875,\ 0),
  \] up to tiny numerical rounding errors (of order \(10^{-16}\)).

  So the constrained tangency portfolio is roughly \textbf{12.5\% in
  NIKKEI225} and \textbf{87.5\% in SP500}, with all other weights very
  close to zero. It still satisfies the budget constraint
  \(\sum_i w_{T,i} = 1\) and remains fully invested, but without any
  leverage.
\item
  Annualised expected return:\\
  \[
  R_T \approx 8.8\%.
  \]
\item
  Annualised volatility:\\
  \[
  \sigma_T \approx 19.4\%.
  \]
\item
  Sharpe ratio (with \(r_f = 0\)):\\
  \[
  \text{Sharpe}_T \approx 0.45.
  \]
\end{itemize}

Compared to the \textbf{unconstrained} tangency portfolio (which was
highly levered and long--short with \(\text{Sharpe} \approx 1.14\)), the
\textbf{no-short} tangency portfolio is much more realistic from a
practical point of view: it is long-only, far less risky in terms of
volatility and leverage, but this comes at the cost of a lower Sharpe
ratio.

Intuitively, once short-selling is forbidden, the optimizer can no
longer exploit extreme long--short positions, and ends up concentrating
on the indices with the best standalone risk--return profile (here
mainly the SP500, with a small allocation to NIKKEI225 for
diversification).

    We clean the tangency weights by setting to zero all entries with
absolute value below a small numerical tolerance (e.g.~\(10^{-10}\)), so
that the portfolio vector is easier to read. We then check that the
cleaned weights still sum to one, which confirms that the fully invested
constraint \(\mathbf{1}^\top w = 1\) remains satisfied after rounding.

    \begin{tcolorbox}[breakable, size=fbox, boxrule=1pt, pad at break*=1mm,colback=cellbackground, colframe=cellborder]
\begin{Verbatim}[commandchars=\\\{\}]
\PY{n}{tol} \PY{o}{=} \PY{l+m+mf}{1e\PYZhy{}10}
\PY{n}{w\PYZus{}tan\PYZus{}noshort\PYZus{}clean} \PY{o}{=} \PY{n}{w\PYZus{}tan\PYZus{}noshort}\PY{o}{.}\PY{n}{copy}\PY{p}{(}\PY{p}{)}
\PY{n}{w\PYZus{}tan\PYZus{}noshort\PYZus{}clean}\PY{p}{[}\PY{n}{np}\PY{o}{.}\PY{n}{abs}\PY{p}{(}\PY{n}{w\PYZus{}tan\PYZus{}noshort\PYZus{}clean}\PY{p}{)} \PY{o}{\PYZlt{}} \PY{n}{tol}\PY{p}{]} \PY{o}{=} \PY{l+m+mf}{0.0}

\PY{n}{w\PYZus{}tan\PYZus{}noshort\PYZus{}clean}\PY{p}{,} \PY{n}{w\PYZus{}tan\PYZus{}noshort\PYZus{}clean}\PY{o}{.}\PY{n}{sum}\PY{p}{(}\PY{p}{)}
\end{Verbatim}
\end{tcolorbox}

            \begin{tcolorbox}[breakable, size=fbox, boxrule=.5pt, pad at break*=1mm, opacityfill=0]
\begin{Verbatim}[commandchars=\\\{\}]
(array([0.        , 0.        , 0.        , 0.        , 0.12519638,
        0.        , 0.87480362, 0.        ]),
 np.float64(1.0000000000000002))
\end{Verbatim}
\end{tcolorbox}
        
    \begin{tcolorbox}[breakable, size=fbox, boxrule=1pt, pad at break*=1mm,colback=cellbackground, colframe=cellborder]
\begin{Verbatim}[commandchars=\\\{\}]
\PY{n+nb}{print}\PY{p}{(}\PY{l+s+s2}{\PYZdq{}}\PY{l+s+s2}{Sum weights:}\PY{l+s+s2}{\PYZdq{}}\PY{p}{,} \PY{n}{w\PYZus{}tan\PYZus{}noshort}\PY{o}{.}\PY{n}{sum}\PY{p}{(}\PY{p}{)}\PY{p}{)}
\PY{n+nb}{print}\PY{p}{(}\PY{l+s+s2}{\PYZdq{}}\PY{l+s+s2}{Min weight:}\PY{l+s+s2}{\PYZdq{}}\PY{p}{,} \PY{n}{w\PYZus{}tan\PYZus{}noshort}\PY{o}{.}\PY{n}{min}\PY{p}{(}\PY{p}{)}\PY{p}{)}
\end{Verbatim}
\end{tcolorbox}

    \begin{Verbatim}[commandchars=\\\{\}]
Sum weights: 1.0000000000000002
Min weight: -2.5569595323416374e-16
    \end{Verbatim}

    The cleaned no-short-sale tangency portfolio has weights

\(w_T \approx (0,\;0,\;0,\;0,\;0.1252,\;0,\;0.8748,\;0)\),

which means that the optimal constrained portfolio is effectively a
two-asset portfolio: about \(12.5\%\) in NIKKEI225 and \(87.5\%\) in
SP500, with zero weight on all the other indices. The cleaned weights
still satisfy

\(\sum_i w_i \approx 1.0000000000000002 \simeq 1\),

and the raw solution gives

\(\text{Sum weights} = 1.0\),
\(\min_i w_i \approx -2.8 \times 10^{-16}\).

This tiny negative value is just a numerical rounding error and can
safely be treated as zero, so the non-negativity constraints
\(w_i \ge 0\) and the fully-invested constraint
\(\mathbf{1}^\top w = 1\) are both satisfied in practice.

    Now we clean the no-short tangency portfolio by forcing very small
numerical values to zero: any weight with absolute value below
\(10^{-10}\) is treated as zero. This gets rid of tiny negative numbers
that only come from floating point rounding. We then convert the cleaned
weight vector into a \texttt{pandas} Series with the asset names as the
index, so that we can clearly see how the portfolio is allocated.

    \begin{tcolorbox}[breakable, size=fbox, boxrule=1pt, pad at break*=1mm,colback=cellbackground, colframe=cellborder]
\begin{Verbatim}[commandchars=\\\{\}]
\PY{n}{w\PYZus{}tan\PYZus{}noshort\PYZus{}clean} \PY{o}{=} \PY{n}{np}\PY{o}{.}\PY{n}{where}\PY{p}{(}\PY{n}{np}\PY{o}{.}\PY{n}{abs}\PY{p}{(}\PY{n}{w\PYZus{}tan\PYZus{}noshort}\PY{p}{)} \PY{o}{\PYZlt{}} \PY{l+m+mf}{1e\PYZhy{}10}\PY{p}{,} \PY{l+m+mi}{0}\PY{p}{,} \PY{n}{w\PYZus{}tan\PYZus{}noshort}\PY{p}{)}
\PY{n}{w\PYZus{}tan\PYZus{}noshort\PYZus{}clean} \PY{o}{=} \PY{n}{pd}\PY{o}{.}\PY{n}{Series}\PY{p}{(}\PY{n}{w\PYZus{}tan\PYZus{}noshort\PYZus{}clean}\PY{p}{,} \PY{n}{index}\PY{o}{=}\PY{n}{assets}\PY{p}{)}
\PY{n}{w\PYZus{}tan\PYZus{}noshort\PYZus{}clean}
\end{Verbatim}
\end{tcolorbox}

            \begin{tcolorbox}[breakable, size=fbox, boxrule=.5pt, pad at break*=1mm, opacityfill=0]
\begin{Verbatim}[commandchars=\\\{\}]
TSX            0.000000
CAC            0.000000
DAX            0.000000
Eurostoxx50    0.000000
NIKKEI225      0.125196
FTSE100        0.000000
SP500          0.874804
IBOVESPA       0.000000
dtype: float64
\end{Verbatim}
\end{tcolorbox}
        
    The cleaned no--short-sale tangency portfolio allocates:

\begin{itemize}
\tightlist
\item
  \textbf{87.48\%} of wealth to \textbf{SP500}
\item
  \textbf{12.52\%} of wealth to \textbf{NIKKEI225}
\item
  \textbf{0\%} to all other indices (TSX, CAC, DAX, Eurostoxx50,
  FTSE100, IBOVESPA)
\end{itemize}

In other words, the optimal constrained tangency portfolio is
effectively a \textbf{two-asset long-only portfolio}, heavily tilted
towards the US equity market (SP500) with a smaller satellite position
in Japanese equities (NIKKEI225). This confirms that, once short-selling
is ruled out, the optimizer concentrates the entire allocation on the
two indices with the best risk--return trade-off in the sample.

    \hypertarget{plot-unconstrained-vs-no-short-efficient-frontiers}{%
\subsubsection{2.6 Plot: unconstrained vs no-short efficient
frontiers}\label{plot-unconstrained-vs-no-short-efficient-frontiers}}

We now put everything together in a single figure.\\
We plot the individual indices using their annualized mean return and
volatility, then overlay the two efficient frontiers:

\begin{itemize}
\tightlist
\item
  the unconstrained frontier (short-selling allowed), and\\
\item
  the frontier with the no short-sale constraint (\(w \ge 0\)).
\end{itemize}

On top of that, we mark the tangency portfolio in each case, so we can
directly see how the shape of the frontier and the location of the
optimal portfolio change once we rule out short positions.

    \begin{tcolorbox}[breakable, size=fbox, boxrule=1pt, pad at break*=1mm,colback=cellbackground, colframe=cellborder]
\begin{Verbatim}[commandchars=\\\{\}]
\PY{n}{plt}\PY{o}{.}\PY{n}{figure}\PY{p}{(}\PY{n}{figsize}\PY{o}{=}\PY{p}{(}\PY{l+m+mi}{9}\PY{p}{,} \PY{l+m+mi}{6}\PY{p}{)}\PY{p}{)}

\PY{c+c1}{\PYZsh{} Assets}
\PY{n}{plt}\PY{o}{.}\PY{n}{scatter}\PY{p}{(}\PY{n}{sigma\PYZus{}series}\PY{o}{.}\PY{n}{values}\PY{p}{,} \PY{n}{mu\PYZus{}series}\PY{o}{.}\PY{n}{values}\PY{p}{,} \PY{n}{label}\PY{o}{=}\PY{l+s+s2}{\PYZdq{}}\PY{l+s+s2}{Assets}\PY{l+s+s2}{\PYZdq{}}\PY{p}{)}
\PY{k}{for} \PY{n}{i}\PY{p}{,} \PY{n}{name} \PY{o+ow}{in} \PY{n+nb}{enumerate}\PY{p}{(}\PY{n}{assets}\PY{p}{)}\PY{p}{:}
    \PY{n}{plt}\PY{o}{.}\PY{n}{annotate}\PY{p}{(}\PY{n}{name}\PY{p}{,}
                 \PY{p}{(}\PY{n}{sigma\PYZus{}series}\PY{o}{.}\PY{n}{values}\PY{p}{[}\PY{n}{i}\PY{p}{]}\PY{p}{,} \PY{n}{mu\PYZus{}series}\PY{o}{.}\PY{n}{values}\PY{p}{[}\PY{n}{i}\PY{p}{]}\PY{p}{)}\PY{p}{,}
                 \PY{n}{fontsize}\PY{o}{=}\PY{l+m+mi}{9}\PY{p}{)}

\PY{c+c1}{\PYZsh{} Frontiers}
\PY{n}{plt}\PY{o}{.}\PY{n}{plot}\PY{p}{(}\PY{n}{frontier\PYZus{}uncon}\PY{p}{[}\PY{p}{:}\PY{p}{,} \PY{l+m+mi}{1}\PY{p}{]}\PY{p}{,} \PY{n}{frontier\PYZus{}uncon}\PY{p}{[}\PY{p}{:}\PY{p}{,} \PY{l+m+mi}{0}\PY{p}{]}\PY{p}{,}
         \PY{n}{label}\PY{o}{=}\PY{l+s+s2}{\PYZdq{}}\PY{l+s+s2}{Unconstrained frontier}\PY{l+s+s2}{\PYZdq{}}\PY{p}{)}
\PY{n}{plt}\PY{o}{.}\PY{n}{plot}\PY{p}{(}\PY{n}{frontier\PYZus{}noshort}\PY{p}{[}\PY{p}{:}\PY{p}{,} \PY{l+m+mi}{1}\PY{p}{]}\PY{p}{,} \PY{n}{frontier\PYZus{}noshort}\PY{p}{[}\PY{p}{:}\PY{p}{,} \PY{l+m+mi}{0}\PY{p}{]}\PY{p}{,}
         \PY{n}{label}\PY{o}{=}\PY{l+s+s2}{\PYZdq{}}\PY{l+s+s2}{No\PYZhy{}short frontier}\PY{l+s+s2}{\PYZdq{}}\PY{p}{)}

\PY{c+c1}{\PYZsh{} Tangency points}
\PY{n}{plt}\PY{o}{.}\PY{n}{scatter}\PY{p}{(}\PY{p}{[}\PY{n}{tan\PYZus{}vol\PYZus{}uncon}\PY{p}{]}\PY{p}{,} \PY{p}{[}\PY{n}{tan\PYZus{}ret\PYZus{}uncon}\PY{p}{]}\PY{p}{,}
            \PY{n}{marker}\PY{o}{=}\PY{l+s+s2}{\PYZdq{}}\PY{l+s+s2}{*}\PY{l+s+s2}{\PYZdq{}}\PY{p}{,} \PY{n}{s}\PY{o}{=}\PY{l+m+mi}{200}\PY{p}{,} \PY{n}{label}\PY{o}{=}\PY{l+s+s2}{\PYZdq{}}\PY{l+s+s2}{Tangency (unconstrained)}\PY{l+s+s2}{\PYZdq{}}\PY{p}{)}
\PY{n}{plt}\PY{o}{.}\PY{n}{scatter}\PY{p}{(}\PY{p}{[}\PY{n}{tan\PYZus{}vol\PYZus{}ns}\PY{p}{]}\PY{p}{,} \PY{p}{[}\PY{n}{tan\PYZus{}ret\PYZus{}ns}\PY{p}{]}\PY{p}{,}
            \PY{n}{marker}\PY{o}{=}\PY{l+s+s2}{\PYZdq{}}\PY{l+s+s2}{*}\PY{l+s+s2}{\PYZdq{}}\PY{p}{,} \PY{n}{s}\PY{o}{=}\PY{l+m+mi}{200}\PY{p}{,} \PY{n}{label}\PY{o}{=}\PY{l+s+s2}{\PYZdq{}}\PY{l+s+s2}{Tangency (no\PYZhy{}short)}\PY{l+s+s2}{\PYZdq{}}\PY{p}{)}

\PY{n}{plt}\PY{o}{.}\PY{n}{title}\PY{p}{(}\PY{l+s+s2}{\PYZdq{}}\PY{l+s+s2}{Efficient Frontier: Unconstrained vs No Short\PYZhy{}Sale}\PY{l+s+s2}{\PYZdq{}}\PY{p}{)}
\PY{n}{plt}\PY{o}{.}\PY{n}{xlabel}\PY{p}{(}\PY{l+s+s2}{\PYZdq{}}\PY{l+s+s2}{Annualized volatility}\PY{l+s+s2}{\PYZdq{}}\PY{p}{)}
\PY{n}{plt}\PY{o}{.}\PY{n}{ylabel}\PY{p}{(}\PY{l+s+s2}{\PYZdq{}}\PY{l+s+s2}{Annualized expected return}\PY{l+s+s2}{\PYZdq{}}\PY{p}{)}
\PY{n}{plt}\PY{o}{.}\PY{n}{legend}\PY{p}{(}\PY{p}{)}

\PY{c+c1}{\PYZsh{} We keep a large scale to see both tangencies}
\PY{n}{plt}\PY{o}{.}\PY{n}{xlim}\PY{p}{(}\PY{l+m+mf}{0.15}\PY{p}{,} \PY{l+m+mf}{0.45}\PY{p}{)}
\PY{n}{plt}\PY{o}{.}\PY{n}{ylim}\PY{p}{(}\PY{o}{\PYZhy{}}\PY{l+m+mf}{0.05}\PY{p}{,} \PY{l+m+mf}{0.55}\PY{p}{)}

\PY{n}{plt}\PY{o}{.}\PY{n}{tight\PYZus{}layout}\PY{p}{(}\PY{p}{)}
\PY{n}{plt}\PY{o}{.}\PY{n}{show}\PY{p}{(}\PY{p}{)}
\end{Verbatim}
\end{tcolorbox}

    \begin{center}
    \adjustimage{max size={0.9\linewidth}{0.9\paperheight}}{Work-7_files/Work-7_78_0.png}
    \end{center}
    { \hspace*{\fill} \\}
    
    This plot compares the Markowitz world with and without short-selling.

\begin{itemize}
\item
  The \textbf{blue curve} is the unconstrained frontier. The
  \textbf{orange curve} is the no-short frontier. The blue one sits
  clearly \textbf{above} the orange one, which shows that allowing short
  positions (and leverage) gives strictly better risk--return
  trade-offs.
\item
  The \textbf{orange star} (unconstrained tangency) is very far to the
  right, with high volatility and a very high expected return. This is
  the crazy long--short portfolio with large positive and negative
  weights we found earlier.
\item
  The \textbf{green star} (no-short tangency) is much closer to the
  equity indices. With \(w_i \ge 0\), the optimal portfolio looks like a
  standard long-only equity mix: lower volatility, lower expected
  return, and a lower Sharpe compared to the unconstrained case.
\end{itemize}

So visually we see the trade-off: allowing short-selling boosts the
efficient frontier and the Sharpe ratio, but at the cost of very extreme
portfolios.

    We now add a second figure to \textbf{zoom in} on the realistic region
of the plot (volatility around 15--30\%, returns below 15\%). On the
full-scale graph, the unconstrained tangency portfolio is so extreme
that it compresses everything else.

The zoomed figure lets us clearly compare the unconstrained and no-short
frontiers around the indices, and see where the no-short tangency
portfolio sits relative to the actual assets.

    \begin{tcolorbox}[breakable, size=fbox, boxrule=1pt, pad at break*=1mm,colback=cellbackground, colframe=cellborder]
\begin{Verbatim}[commandchars=\\\{\}]
\PY{c+c1}{\PYZsh{} Second figure: zoom on realistic region (no\PYZhy{}short frontier and assets)}
\PY{n}{plt}\PY{o}{.}\PY{n}{figure}\PY{p}{(}\PY{n}{figsize}\PY{o}{=}\PY{p}{(}\PY{l+m+mi}{9}\PY{p}{,} \PY{l+m+mi}{6}\PY{p}{)}\PY{p}{)}

\PY{c+c1}{\PYZsh{} Assets}
\PY{n}{plt}\PY{o}{.}\PY{n}{scatter}\PY{p}{(}\PY{n}{sigma\PYZus{}series}\PY{o}{.}\PY{n}{values}\PY{p}{,} \PY{n}{mu\PYZus{}series}\PY{o}{.}\PY{n}{values}\PY{p}{,} \PY{n}{label}\PY{o}{=}\PY{l+s+s2}{\PYZdq{}}\PY{l+s+s2}{Assets}\PY{l+s+s2}{\PYZdq{}}\PY{p}{)}
\PY{k}{for} \PY{n}{i}\PY{p}{,} \PY{n}{name} \PY{o+ow}{in} \PY{n+nb}{enumerate}\PY{p}{(}\PY{n}{assets}\PY{p}{)}\PY{p}{:}
    \PY{n}{plt}\PY{o}{.}\PY{n}{annotate}\PY{p}{(}\PY{n}{name}\PY{p}{,}
                 \PY{p}{(}\PY{n}{sigma\PYZus{}series}\PY{o}{.}\PY{n}{values}\PY{p}{[}\PY{n}{i}\PY{p}{]}\PY{p}{,} \PY{n}{mu\PYZus{}series}\PY{o}{.}\PY{n}{values}\PY{p}{[}\PY{n}{i}\PY{p}{]}\PY{p}{)}\PY{p}{,}
                 \PY{n}{fontsize}\PY{o}{=}\PY{l+m+mi}{9}\PY{p}{)}

\PY{c+c1}{\PYZsh{} Frontiers (we can keep both but the zoom value the no\PYZhy{}short)}
\PY{n}{plt}\PY{o}{.}\PY{n}{plot}\PY{p}{(}\PY{n}{frontier\PYZus{}uncon}\PY{p}{[}\PY{p}{:}\PY{p}{,} \PY{l+m+mi}{1}\PY{p}{]}\PY{p}{,} \PY{n}{frontier\PYZus{}uncon}\PY{p}{[}\PY{p}{:}\PY{p}{,} \PY{l+m+mi}{0}\PY{p}{]}\PY{p}{,}
         \PY{n}{label}\PY{o}{=}\PY{l+s+s2}{\PYZdq{}}\PY{l+s+s2}{Unconstrained frontier}\PY{l+s+s2}{\PYZdq{}}\PY{p}{)}
\PY{n}{plt}\PY{o}{.}\PY{n}{plot}\PY{p}{(}\PY{n}{frontier\PYZus{}noshort}\PY{p}{[}\PY{p}{:}\PY{p}{,} \PY{l+m+mi}{1}\PY{p}{]}\PY{p}{,} \PY{n}{frontier\PYZus{}noshort}\PY{p}{[}\PY{p}{:}\PY{p}{,} \PY{l+m+mi}{0}\PY{p}{]}\PY{p}{,}
         \PY{n}{label}\PY{o}{=}\PY{l+s+s2}{\PYZdq{}}\PY{l+s+s2}{No\PYZhy{}short frontier}\PY{l+s+s2}{\PYZdq{}}\PY{p}{)}

\PY{c+c1}{\PYZsh{} Tangency points}

\PY{n}{plt}\PY{o}{.}\PY{n}{scatter}\PY{p}{(}\PY{p}{[}\PY{n}{tan\PYZus{}vol\PYZus{}ns}\PY{p}{]}\PY{p}{,} \PY{p}{[}\PY{n}{tan\PYZus{}ret\PYZus{}ns}\PY{p}{]}\PY{p}{,}
            \PY{n}{marker}\PY{o}{=}\PY{l+s+s2}{\PYZdq{}}\PY{l+s+s2}{*}\PY{l+s+s2}{\PYZdq{}}\PY{p}{,} \PY{n}{s}\PY{o}{=}\PY{l+m+mi}{200}\PY{p}{,} \PY{n}{label}\PY{o}{=}\PY{l+s+s2}{\PYZdq{}}\PY{l+s+s2}{Tangency (no\PYZhy{}short)}\PY{l+s+s2}{\PYZdq{}}\PY{p}{)}

\PY{n}{plt}\PY{o}{.}\PY{n}{title}\PY{p}{(}\PY{l+s+s2}{\PYZdq{}}\PY{l+s+s2}{Efficient Frontier (Zoom on Realistic Region)}\PY{l+s+s2}{\PYZdq{}}\PY{p}{)}
\PY{n}{plt}\PY{o}{.}\PY{n}{xlabel}\PY{p}{(}\PY{l+s+s2}{\PYZdq{}}\PY{l+s+s2}{Annualized volatility}\PY{l+s+s2}{\PYZdq{}}\PY{p}{)}
\PY{n}{plt}\PY{o}{.}\PY{n}{ylabel}\PY{p}{(}\PY{l+s+s2}{\PYZdq{}}\PY{l+s+s2}{Annualized expected return}\PY{l+s+s2}{\PYZdq{}}\PY{p}{)}
\PY{n}{plt}\PY{o}{.}\PY{n}{legend}\PY{p}{(}\PY{p}{)}

\PY{c+c1}{\PYZsh{} Zoom on the zone where indexes live and the tangency no\PYZhy{}short}
\PY{n}{plt}\PY{o}{.}\PY{n}{xlim}\PY{p}{(}\PY{l+m+mf}{0.15}\PY{p}{,} \PY{l+m+mf}{0.35}\PY{p}{)}
\PY{n}{plt}\PY{o}{.}\PY{n}{ylim}\PY{p}{(}\PY{o}{\PYZhy{}}\PY{l+m+mf}{0.05}\PY{p}{,} \PY{l+m+mf}{0.15}\PY{p}{)}

\PY{n}{plt}\PY{o}{.}\PY{n}{tight\PYZus{}layout}\PY{p}{(}\PY{p}{)}
\PY{n}{plt}\PY{o}{.}\PY{n}{show}\PY{p}{(}\PY{p}{)}
\end{Verbatim}
\end{tcolorbox}

    \begin{center}
    \adjustimage{max size={0.9\linewidth}{0.9\paperheight}}{Work-7_files/Work-7_81_0.png}
    \end{center}
    { \hspace*{\fill} \\}
    
    On the zoomed plot we focus on the ``realistic'' region (volatility
between about 15\% and 35\%, returns between −5\% and 15\%). This is why
\textbf{IBOVESPA does not appear} here: its annual volatility is around
40\%, so it lies outside the chosen \(x\)-axis range.

In this region we clearly see that:

\begin{itemize}
\item
  The \textbf{no-short frontier} (orange) is below the
  \textbf{unconstrained frontier} (blue), as expected: once we forbid
  short positions, the set of feasible portfolios shrinks and
  risk--return trade-offs become worse.
\item
  The \textbf{no-short tangency portfolio} (orange star) sits around
  19\% volatility and 8.8\% expected return. It lies above all
  individual indices in this zoom, especially above the SP500 point,
  which shows that a long-only combination of NIKKEI225 and SP500 can
  still dominate any single market index in terms of Sharpe ratio.
\end{itemize}

Overall, the zoomed figure makes it easier to compare the two frontiers
and to see that the long-only efficient set is much closer to the cloud
of indices than the extreme unconstrained frontier.

    \hypertarget{portfolio-weights}{%
\subsubsection{2.7 Portfolio weights}\label{portfolio-weights}}

We now report and compare the tangency portfolio weights for:

\begin{itemize}
\tightlist
\item
  the unconstrained case (short-selling allowed), and\\
\item
  the no short-sale case (weights constrained to be non-negative).
\end{itemize}

This allows us to see how the no-short constraint changes the
composition and concentration of the optimal portfolio.

    To make the comparison more transparent, we collect the tangency
portfolio weights in a single table. The first column reports the
weights of the unconstrained tangency portfolio (short-selling allowed),
and the second column shows the tangency portfolio when we impose the no
short-sale constraint \(w_i \ge 0\). In both cases, the weights sum to
one.

    \begin{tcolorbox}[breakable, size=fbox, boxrule=1pt, pad at break*=1mm,colback=cellbackground, colframe=cellborder]
\begin{Verbatim}[commandchars=\\\{\}]
\PY{c+c1}{\PYZsh{} Reminder: w\PYZus{}tan\PYZus{}uncon, w\PYZus{}tan\PYZus{}noshort et w\PYZus{}tan\PYZus{}noshort\PYZus{}clean already computed above}

\PY{c+c1}{\PYZsh{} Table of tangency portfolio weights}
\PY{n}{weights\PYZus{}df} \PY{o}{=} \PY{n}{pd}\PY{o}{.}\PY{n}{DataFrame}\PY{p}{(}\PY{p}{\PYZob{}}
    \PY{l+s+s2}{\PYZdq{}}\PY{l+s+s2}{Tangency\PYZus{}Unconstrained}\PY{l+s+s2}{\PYZdq{}}\PY{p}{:} \PY{n}{w\PYZus{}tan\PYZus{}uncon}\PY{p}{,}
    \PY{l+s+s2}{\PYZdq{}}\PY{l+s+s2}{Tangency\PYZus{}NoShort}\PY{l+s+s2}{\PYZdq{}}\PY{p}{:} \PY{n}{w\PYZus{}tan\PYZus{}noshort\PYZus{}clean}  \PY{c+c1}{\PYZsh{} cleaned version (we set tiny values to 0)}
\PY{p}{\PYZcb{}}\PY{p}{,} \PY{n}{index}\PY{o}{=}\PY{n}{assets}\PY{p}{)}

\PY{c+c1}{\PYZsh{} Rounded version for a better presentation!}
\PY{n}{weights\PYZus{}df\PYZus{}rounded} \PY{o}{=} \PY{n}{weights\PYZus{}df}\PY{o}{.}\PY{n}{round}\PY{p}{(}\PY{l+m+mi}{4}\PY{p}{)}
\PY{n}{weights\PYZus{}df\PYZus{}rounded}
\end{Verbatim}
\end{tcolorbox}

            \begin{tcolorbox}[breakable, size=fbox, boxrule=.5pt, pad at break*=1mm, opacityfill=0]
\begin{Verbatim}[commandchars=\\\{\}]
             Tangency\_Unconstrained  Tangency\_NoShort
TSX                         -0.4824            0.0000
CAC                          3.0784            0.0000
DAX                          4.0278            0.0000
Eurostoxx50                 -6.6824            0.0000
NIKKEI225                    0.4463            0.1252
FTSE100                     -0.8347            0.0000
SP500                        1.7312            0.8748
IBOVESPA                    -0.2843            0.0000
\end{Verbatim}
\end{tcolorbox}
        
    We now compare the tangency portfolios with and without short-selling:

\begin{itemize}
\item
  In the \textbf{unconstrained case}, the optimizer uses very
  \textbf{large long and short positions}.\\
  For example, it takes big long positions in \textbf{CAC} and
  \textbf{DAX}, but also large short positions in \textbf{Eurostoxx50},
  \textbf{TSX}, \textbf{FTSE100} and \textbf{IBOVESPA}. The weights
  still sum to 1, but \(\sum_i |w_i|\) is much larger than 1, which
  means the portfolio is highly levered. This is a typical feature of
  the Markowitz solution when we do not impose any practical constraint.
\item
  With the \textbf{no short-sale constraint} (\(w_i \ge 0\)), the
  tangency portfolio becomes much simpler and more realistic: it is
  essentially a combination of \textbf{SP500} (about \(87\%\)) and
  \textbf{NIKKEI225} (about \(13\%\)), with zero weight on all other
  indices. Under our plug-in estimates, these two markets offer the best
  trade-off between risk and return once we rule out short positions.
\end{itemize}

Overall, the comparison clearly shows that short-selling leads to very
extreme and levered Markowitz portfolios, while the no-short constraint
produces a more concentrated but more plausible allocation.

    \hypertarget{conclusion-for-question-2}{%
\subsubsection{2.8 Conclusion for Question
2}\label{conclusion-for-question-2}}

In this part, we used historical daily USD returns to estimate
annualised means \(\mu\) and the covariance matrix \(\Sigma\), and we
mapped the Markowitz problem into the \texttt{quadprog.solve\_qp} format
with the constraints \(1^\top w = 1\) and
\(\mu^\top w = \mu_{\text{target}}\).

On this basis we traced the mean--variance efficient frontier and found
the \textbf{unconstrained tangency portfolio}, with annualised return
\(R_T \approx 48.3\%\), volatility \(\sigma_T \approx 42.4\%\) and
Sharpe ratio \(\text{Sharpe}_T \approx 1.14\). It looks extremely
attractive on paper, but it achieves this performance by taking very
large long and short positions (leverage \(\sum_i |w_i| \approx 17.6\)),
which would be hard to implement in practice because of margin
requirements, risk limits and typical regulations for institutional
investors.

When we repeat the exercise under the \textbf{no short-sale constraint}
\(w_i \ge 0\), the efficient frontier shifts down and the tangency
portfolio becomes a much more realistic long-only allocation: roughly
\(87.5\%\) in SP500 and \(12.5\%\) in NIKKEI225, with annualised return
\(R_T \approx 8.8\%\), volatility \(\sigma_T \approx 19.4\%\) and Sharpe
ratio \(\text{Sharpe}_T \approx 0.45\). The constrained portfolio is
less ``impressive'\,' in terms of Sharpe ratio, but it is far closer to
what an actual fund (e.g.~a pension fund with capital and regulatory
constraints) could hold.

Overall, Question 2 shows clearly how adding simple constraints like
\(w_i \ge 0\) dramatically changes both the efficient frontier and the
composition of the ``optimal'\,' portfolio.

    \hypertarget{question-3-mean-variance-mv-analysis-2}{%
\subsection{Question 3 -- Mean-Variance (MV) analysis
\#2}\label{question-3-mean-variance-mv-analysis-2}}

In this part we repeat the mean--variance analysis using rolling 5-year
windows.\\
The idea is to mimic an investor who periodically re-estimates expected
returns and covariances using only the most recent 5 years of data, and
then updates their optimal portfolio.

Concretely, we start with the window from 2010--2014, then roll the
window forward one year at a time until the last window 2017--2021. For
each window, we compute the mean--variance tangency portfolio (using the
same plug-in estimates and \texttt{quadprog} mapping as in Question 2),
and at the end we compare these time-varying MV weights to the
corresponding market cap-based weights.

    \begin{tcolorbox}[breakable, size=fbox, boxrule=1pt, pad at break*=1mm,colback=cellbackground, colframe=cellborder]
\begin{Verbatim}[commandchars=\\\{\}]
\PY{c+c1}{\PYZsh{} Define 5\PYZhy{}year rolling windows by calendar year}
\PY{n}{window\PYZus{}years} \PY{o}{=} \PY{p}{[}\PY{p}{(}\PY{l+m+mi}{2010}\PY{p}{,} \PY{l+m+mi}{2014}\PY{p}{)}\PY{p}{,}
                \PY{p}{(}\PY{l+m+mi}{2011}\PY{p}{,} \PY{l+m+mi}{2015}\PY{p}{)}\PY{p}{,}
                \PY{p}{(}\PY{l+m+mi}{2012}\PY{p}{,} \PY{l+m+mi}{2016}\PY{p}{)}\PY{p}{,}
                \PY{p}{(}\PY{l+m+mi}{2013}\PY{p}{,} \PY{l+m+mi}{2017}\PY{p}{)}\PY{p}{,}
                \PY{p}{(}\PY{l+m+mi}{2014}\PY{p}{,} \PY{l+m+mi}{2018}\PY{p}{)}\PY{p}{,}
                \PY{p}{(}\PY{l+m+mi}{2015}\PY{p}{,} \PY{l+m+mi}{2019}\PY{p}{)}\PY{p}{,}
                \PY{p}{(}\PY{l+m+mi}{2016}\PY{p}{,} \PY{l+m+mi}{2020}\PY{p}{)}\PY{p}{,}
                \PY{p}{(}\PY{l+m+mi}{2017}\PY{p}{,} \PY{l+m+mi}{2021}\PY{p}{)}\PY{p}{]}

\PY{n}{window\PYZus{}labels} \PY{o}{=} \PY{p}{[}\PY{l+s+sa}{f}\PY{l+s+s2}{\PYZdq{}}\PY{l+s+si}{\PYZob{}}\PY{n}{start}\PY{l+s+si}{\PYZcb{}}\PY{l+s+s2}{\PYZhy{}}\PY{l+s+si}{\PYZob{}}\PY{n}{end}\PY{l+s+si}{\PYZcb{}}\PY{l+s+s2}{\PYZdq{}} \PY{k}{for} \PY{p}{(}\PY{n}{start}\PY{p}{,} \PY{n}{end}\PY{p}{)} \PY{o+ow}{in} \PY{n}{window\PYZus{}years}\PY{p}{]}
\PY{n}{window\PYZus{}labels}
\end{Verbatim}
\end{tcolorbox}

            \begin{tcolorbox}[breakable, size=fbox, boxrule=.5pt, pad at break*=1mm, opacityfill=0]
\begin{Verbatim}[commandchars=\\\{\}]
['2010-2014',
 '2011-2015',
 '2012-2016',
 '2013-2017',
 '2014-2018',
 '2015-2019',
 '2016-2020',
 '2017-2021']
\end{Verbatim}
\end{tcolorbox}
        
    \hypertarget{rolling-window-set-up-and-estimation}{%
\subsubsection{3.1 Rolling-window set-up and
estimation}\label{rolling-window-set-up-and-estimation}}

For each 5-year window, we:

\begin{enumerate}
\def\labelenumi{\arabic{enumi}.}
\tightlist
\item
  Restrict the daily log-return data to that period.
\item
  Recompute the annualized expected returns and covariance matrix.
\item
  Compute the unconstrained tangency portfolio using the closed-form
  formula from the course notes.
\item
  Compute the no-short tangency portfolio by searching along the
  constrained efficient frontier and picking the portfolio with the
  highest Sharpe ratio.
\end{enumerate}

We then collect the weights across windows to see how unstable (or
stable) the mean--variance portfolios are over time.

    \begin{tcolorbox}[breakable, size=fbox, boxrule=1pt, pad at break*=1mm,colback=cellbackground, colframe=cellborder]
\begin{Verbatim}[commandchars=\\\{\}]
\PY{n}{tan\PYZus{}uncon\PYZus{}weights} \PY{o}{=} \PY{p}{[}\PY{p}{]}
\PY{n}{tan\PYZus{}noshort\PYZus{}weights} \PY{o}{=} \PY{p}{[}\PY{p}{]}
\PY{n}{tan\PYZus{}uncon\PYZus{}stats} \PY{o}{=} \PY{p}{[}\PY{p}{]}   \PY{c+c1}{\PYZsh{} (ret, vol, sharpe)}
\PY{n}{tan\PYZus{}noshort\PYZus{}stats} \PY{o}{=} \PY{p}{[}\PY{p}{]}

\PY{k}{for} \PY{p}{(}\PY{n}{start\PYZus{}year}\PY{p}{,} \PY{n}{end\PYZus{}year}\PY{p}{)}\PY{p}{,} \PY{n}{label} \PY{o+ow}{in} \PY{n+nb}{zip}\PY{p}{(}\PY{n}{window\PYZus{}years}\PY{p}{,} \PY{n}{window\PYZus{}labels}\PY{p}{)}\PY{p}{:}
    \PY{c+c1}{\PYZsh{} Select 5\PYZhy{}year window }
    \PY{n}{start\PYZus{}date} \PY{o}{=} \PY{l+s+sa}{f}\PY{l+s+s2}{\PYZdq{}}\PY{l+s+si}{\PYZob{}}\PY{n}{start\PYZus{}year}\PY{l+s+si}{\PYZcb{}}\PY{l+s+s2}{\PYZhy{}01\PYZhy{}01}\PY{l+s+s2}{\PYZdq{}}
    \PY{n}{end\PYZus{}date} \PY{o}{=} \PY{l+s+sa}{f}\PY{l+s+s2}{\PYZdq{}}\PY{l+s+si}{\PYZob{}}\PY{n}{end\PYZus{}year}\PY{l+s+si}{\PYZcb{}}\PY{l+s+s2}{\PYZhy{}12\PYZhy{}31}\PY{l+s+s2}{\PYZdq{}}
    \PY{n}{window\PYZus{}data} \PY{o}{=} \PY{n}{returns}\PY{o}{.}\PY{n}{loc}\PY{p}{[}\PY{n}{start\PYZus{}date}\PY{p}{:}\PY{n}{end\PYZus{}date}\PY{p}{]}
    
    \PY{c+c1}{\PYZsh{} If for some reason window is empty, skip}
    \PY{k}{if} \PY{n}{window\PYZus{}data}\PY{o}{.}\PY{n}{empty}\PY{p}{:}
        \PY{n+nb}{print}\PY{p}{(}\PY{l+s+sa}{f}\PY{l+s+s2}{\PYZdq{}}\PY{l+s+s2}{Warning: no data for window }\PY{l+s+si}{\PYZob{}}\PY{n}{label}\PY{l+s+si}{\PYZcb{}}\PY{l+s+s2}{\PYZdq{}}\PY{p}{)}
        \PY{n}{tan\PYZus{}uncon\PYZus{}weights}\PY{o}{.}\PY{n}{append}\PY{p}{(}\PY{n}{np}\PY{o}{.}\PY{n}{nan} \PY{o}{*} \PY{n}{np}\PY{o}{.}\PY{n}{ones}\PY{p}{(}\PY{n}{n}\PY{p}{)}\PY{p}{)}
        \PY{n}{tan\PYZus{}noshort\PYZus{}weights}\PY{o}{.}\PY{n}{append}\PY{p}{(}\PY{n}{np}\PY{o}{.}\PY{n}{nan} \PY{o}{*} \PY{n}{np}\PY{o}{.}\PY{n}{ones}\PY{p}{(}\PY{n}{n}\PY{p}{)}\PY{p}{)}
        \PY{n}{tan\PYZus{}uncon\PYZus{}stats}\PY{o}{.}\PY{n}{append}\PY{p}{(}\PY{p}{(}\PY{n}{np}\PY{o}{.}\PY{n}{nan}\PY{p}{,} \PY{n}{np}\PY{o}{.}\PY{n}{nan}\PY{p}{,} \PY{n}{np}\PY{o}{.}\PY{n}{nan}\PY{p}{)}\PY{p}{)}
        \PY{n}{tan\PYZus{}noshort\PYZus{}stats}\PY{o}{.}\PY{n}{append}\PY{p}{(}\PY{p}{(}\PY{n}{np}\PY{o}{.}\PY{n}{nan}\PY{p}{,} \PY{n}{np}\PY{o}{.}\PY{n}{nan}\PY{p}{,} \PY{n}{np}\PY{o}{.}\PY{n}{nan}\PY{p}{)}\PY{p}{)}
        \PY{k}{continue}
    
    \PY{c+c1}{\PYZsh{} We recompute daily mean and covariance for this window}
    \PY{n}{mu\PYZus{}daily\PYZus{}w} \PY{o}{=} \PY{n}{window\PYZus{}data}\PY{o}{.}\PY{n}{mean}\PY{p}{(}\PY{p}{)}\PY{o}{.}\PY{n}{values}
    \PY{n}{Sigma\PYZus{}daily\PYZus{}w} \PY{o}{=} \PY{n}{window\PYZus{}data}\PY{o}{.}\PY{n}{cov}\PY{p}{(}\PY{p}{)}\PY{o}{.}\PY{n}{values}
    
    \PY{n}{mu\PYZus{}w} \PY{o}{=} \PY{n}{mu\PYZus{}daily\PYZus{}w} \PY{o}{*} \PY{n}{ann\PYZus{}factor}
    \PY{n}{Sigma\PYZus{}w} \PY{o}{=} \PY{n}{Sigma\PYZus{}daily\PYZus{}w} \PY{o}{*} \PY{n}{ann\PYZus{}factor}
    \PY{n}{mu\PYZus{}excess\PYZus{}w} \PY{o}{=} \PY{n}{mu\PYZus{}w} \PY{o}{\PYZhy{}} \PY{n}{rf} \PY{o}{*} \PY{n}{np}\PY{o}{.}\PY{n}{ones}\PY{p}{(}\PY{n}{n}\PY{p}{)}
    
    \PY{c+c1}{\PYZsh{} Unconstrained tangency (closed form)}
    \PY{n}{invSigma\PYZus{}w} \PY{o}{=} \PY{n}{np}\PY{o}{.}\PY{n}{linalg}\PY{o}{.}\PY{n}{inv}\PY{p}{(}\PY{n}{Sigma\PYZus{}w}\PY{p}{)}
    \PY{n}{numer\PYZus{}w} \PY{o}{=} \PY{n}{invSigma\PYZus{}w} \PY{o}{@} \PY{n}{mu\PYZus{}excess\PYZus{}w}
    \PY{n}{denom\PYZus{}w} \PY{o}{=} \PY{n}{np}\PY{o}{.}\PY{n}{ones}\PY{p}{(}\PY{n}{n}\PY{p}{)} \PY{o}{@} \PY{n}{numer\PYZus{}w}
    \PY{n}{w\PYZus{}tan\PYZus{}uncon\PYZus{}w} \PY{o}{=} \PY{n}{numer\PYZus{}w} \PY{o}{/} \PY{n}{denom\PYZus{}w}
    
    \PY{n}{ret\PYZus{}u}\PY{p}{,} \PY{n}{vol\PYZus{}u}\PY{p}{,} \PY{n}{sh\PYZus{}u} \PY{o}{=} \PY{n}{portfolio\PYZus{}stats}\PY{p}{(}\PY{n}{w\PYZus{}tan\PYZus{}uncon\PYZus{}w}\PY{p}{,} \PY{n}{mu\PYZus{}w}\PY{p}{,} \PY{n}{Sigma\PYZus{}w}\PY{p}{,} \PY{n}{rf}\PY{o}{=}\PY{n}{rf}\PY{p}{)}
    
    \PY{n}{tan\PYZus{}uncon\PYZus{}weights}\PY{o}{.}\PY{n}{append}\PY{p}{(}\PY{n}{w\PYZus{}tan\PYZus{}uncon\PYZus{}w}\PY{p}{)}
    \PY{n}{tan\PYZus{}uncon\PYZus{}stats}\PY{o}{.}\PY{n}{append}\PY{p}{(}\PY{p}{(}\PY{n}{ret\PYZus{}u}\PY{p}{,} \PY{n}{vol\PYZus{}u}\PY{p}{,} \PY{n}{sh\PYZus{}u}\PY{p}{)}\PY{p}{)}
    
    \PY{c+c1}{\PYZsh{} No\PYZhy{}short tangency: search along constrained frontier}
    \PY{n}{target\PYZus{}grid\PYZus{}w} \PY{o}{=} \PY{n}{np}\PY{o}{.}\PY{n}{linspace}\PY{p}{(}\PY{n}{mu\PYZus{}w}\PY{o}{.}\PY{n}{min}\PY{p}{(}\PY{p}{)}\PY{p}{,} \PY{n}{mu\PYZus{}w}\PY{o}{.}\PY{n}{max}\PY{p}{(}\PY{p}{)}\PY{p}{,} \PY{l+m+mi}{60}\PY{p}{)}
    \PY{n}{frontier\PYZus{}noshort\PYZus{}w} \PY{o}{=} \PY{p}{[}\PY{p}{]}
    \PY{n}{weights\PYZus{}noshort\PYZus{}w} \PY{o}{=} \PY{p}{[}\PY{p}{]}
    
    \PY{k}{for} \PY{n}{t\PYZus{}ret} \PY{o+ow}{in} \PY{n}{target\PYZus{}grid\PYZus{}w}\PY{p}{:}
        \PY{k}{try}\PY{p}{:}
            \PY{n}{w\PYZus{}ns} \PY{o}{=} \PY{n}{solve\PYZus{}mv\PYZus{}target\PYZus{}return\PYZus{}quadprog}\PY{p}{(}\PY{n}{mu\PYZus{}w}\PY{p}{,} \PY{n}{Sigma\PYZus{}w}\PY{p}{,} \PY{n}{t\PYZus{}ret}\PY{p}{,} \PY{n}{no\PYZus{}short}\PY{o}{=}\PY{k+kc}{True}\PY{p}{)}
            \PY{n}{r\PYZus{}ns}\PY{p}{,} \PY{n}{v\PYZus{}ns}\PY{p}{,} \PY{n}{s\PYZus{}ns} \PY{o}{=} \PY{n}{portfolio\PYZus{}stats}\PY{p}{(}\PY{n}{w\PYZus{}ns}\PY{p}{,} \PY{n}{mu\PYZus{}w}\PY{p}{,} \PY{n}{Sigma\PYZus{}w}\PY{p}{,} \PY{n}{rf}\PY{o}{=}\PY{n}{rf}\PY{p}{)}
            \PY{n}{frontier\PYZus{}noshort\PYZus{}w}\PY{o}{.}\PY{n}{append}\PY{p}{(}\PY{p}{(}\PY{n}{r\PYZus{}ns}\PY{p}{,} \PY{n}{v\PYZus{}ns}\PY{p}{,} \PY{n}{s\PYZus{}ns}\PY{p}{)}\PY{p}{)}
            \PY{n}{weights\PYZus{}noshort\PYZus{}w}\PY{o}{.}\PY{n}{append}\PY{p}{(}\PY{n}{w\PYZus{}ns}\PY{p}{)}
        \PY{k}{except} \PY{n+ne}{ValueError}\PY{p}{:}
            \PY{c+c1}{\PYZsh{} Infeasible target under no\PYZhy{}short \PYZhy{}\PYZgt{} skip}
            \PY{k}{continue}
    
    \PY{k}{if} \PY{n+nb}{len}\PY{p}{(}\PY{n}{frontier\PYZus{}noshort\PYZus{}w}\PY{p}{)} \PY{o}{==} \PY{l+m+mi}{0}\PY{p}{:}
        \PY{c+c1}{\PYZsh{} No feasible portfolio found (unlikely, but just in case)}
        \PY{n+nb}{print}\PY{p}{(}\PY{l+s+sa}{f}\PY{l+s+s2}{\PYZdq{}}\PY{l+s+s2}{Warning: no feasible no\PYZhy{}short solution for window }\PY{l+s+si}{\PYZob{}}\PY{n}{label}\PY{l+s+si}{\PYZcb{}}\PY{l+s+s2}{\PYZdq{}}\PY{p}{)}
        \PY{n}{tan\PYZus{}noshort\PYZus{}weights}\PY{o}{.}\PY{n}{append}\PY{p}{(}\PY{n}{np}\PY{o}{.}\PY{n}{nan} \PY{o}{*} \PY{n}{np}\PY{o}{.}\PY{n}{ones}\PY{p}{(}\PY{n}{n}\PY{p}{)}\PY{p}{)}
        \PY{n}{tan\PYZus{}noshort\PYZus{}stats}\PY{o}{.}\PY{n}{append}\PY{p}{(}\PY{p}{(}\PY{n}{np}\PY{o}{.}\PY{n}{nan}\PY{p}{,} \PY{n}{np}\PY{o}{.}\PY{n}{nan}\PY{p}{,} \PY{n}{np}\PY{o}{.}\PY{n}{nan}\PY{p}{)}\PY{p}{)}
    \PY{k}{else}\PY{p}{:}
        \PY{n}{frontier\PYZus{}noshort\PYZus{}w} \PY{o}{=} \PY{n}{np}\PY{o}{.}\PY{n}{array}\PY{p}{(}\PY{n}{frontier\PYZus{}noshort\PYZus{}w}\PY{p}{)}
        \PY{n}{idx\PYZus{}best} \PY{o}{=} \PY{n}{np}\PY{o}{.}\PY{n}{nanargmax}\PY{p}{(}\PY{n}{frontier\PYZus{}noshort\PYZus{}w}\PY{p}{[}\PY{p}{:}\PY{p}{,} \PY{l+m+mi}{2}\PY{p}{]}\PY{p}{)}  \PY{c+c1}{\PYZsh{} max Sharpe}
        \PY{n}{w\PYZus{}tan\PYZus{}noshort\PYZus{}w} \PY{o}{=} \PY{n}{weights\PYZus{}noshort\PYZus{}w}\PY{p}{[}\PY{n}{idx\PYZus{}best}\PY{p}{]}
        \PY{n}{r\PYZus{}ns}\PY{p}{,} \PY{n}{v\PYZus{}ns}\PY{p}{,} \PY{n}{s\PYZus{}ns} \PY{o}{=} \PY{n}{frontier\PYZus{}noshort\PYZus{}w}\PY{p}{[}\PY{n}{idx\PYZus{}best}\PY{p}{]}
        
        \PY{n}{tan\PYZus{}noshort\PYZus{}weights}\PY{o}{.}\PY{n}{append}\PY{p}{(}\PY{n}{w\PYZus{}tan\PYZus{}noshort\PYZus{}w}\PY{p}{)}
        \PY{n}{tan\PYZus{}noshort\PYZus{}stats}\PY{o}{.}\PY{n}{append}\PY{p}{(}\PY{p}{(}\PY{n}{r\PYZus{}ns}\PY{p}{,} \PY{n}{v\PYZus{}ns}\PY{p}{,} \PY{n}{s\PYZus{}ns}\PY{p}{)}\PY{p}{)}
\end{Verbatim}
\end{tcolorbox}

    In this loop we basically rerun the full Markowitz procedure for each
5-year window. For every window, we:

\begin{itemize}
\tightlist
\item
  recompute the annualised mean vector and covariance matrix using only
  the returns in that period;
\item
  build an \textbf{unconstrained} tangency portfolio using the
  closed-form \(\Sigma^{-1}\mu\) formula, and record its return,
  volatility, and Sharpe ratio;
\item
  build a \textbf{no short-sale} tangency portfolio by scanning a grid
  of target returns, solving the quadratic program with \(w \ge 0\), and
  keeping the feasible portfolio with the highest Sharpe ratio.
\end{itemize}

We store both the weights and the performance statistics for each window
so we can later turn them into DataFrames and see how these ``optimal''
portfolios move over time.

    \hypertarget{rolling-tangency-portfolios-instability-and-concentration}{%
\subsubsection{3.2 Rolling tangency portfolios: instability and
concentration}\label{rolling-tangency-portfolios-instability-and-concentration}}

We now organise the rolling-window results into simple tables. The code
below turns the stored weight vectors and performance statistics into
DataFrames, with one row per 5-year window and one column per index (or
per moment). This makes it easier to see how the tangency portfolios and
their Sharpe ratios evolve over time, and to check how concentrated and
unstable they are across windows.

    \begin{tcolorbox}[breakable, size=fbox, boxrule=1pt, pad at break*=1mm,colback=cellbackground, colframe=cellborder]
\begin{Verbatim}[commandchars=\\\{\}]
\PY{n}{tan\PYZus{}uncon\PYZus{}weights\PYZus{}df} \PY{o}{=} \PY{n}{pd}\PY{o}{.}\PY{n}{DataFrame}\PY{p}{(}\PY{n}{tan\PYZus{}uncon\PYZus{}weights}\PY{p}{,} \PY{n}{index}\PY{o}{=}\PY{n}{window\PYZus{}labels}\PY{p}{,} \PY{n}{columns}\PY{o}{=}\PY{n}{assets}\PY{p}{)}
\PY{n}{tan\PYZus{}noshort\PYZus{}weights\PYZus{}df} \PY{o}{=} \PY{n}{pd}\PY{o}{.}\PY{n}{DataFrame}\PY{p}{(}\PY{n}{tan\PYZus{}noshort\PYZus{}weights}\PY{p}{,} \PY{n}{index}\PY{o}{=}\PY{n}{window\PYZus{}labels}\PY{p}{,} \PY{n}{columns}\PY{o}{=}\PY{n}{assets}\PY{p}{)}

\PY{n}{tan\PYZus{}uncon\PYZus{}stats\PYZus{}df} \PY{o}{=} \PY{n}{pd}\PY{o}{.}\PY{n}{DataFrame}\PY{p}{(}\PY{n}{tan\PYZus{}uncon\PYZus{}stats}\PY{p}{,} 
                                  \PY{n}{index}\PY{o}{=}\PY{n}{window\PYZus{}labels}\PY{p}{,} 
                                  \PY{n}{columns}\PY{o}{=}\PY{p}{[}\PY{l+s+s2}{\PYZdq{}}\PY{l+s+s2}{ann\PYZus{}return}\PY{l+s+s2}{\PYZdq{}}\PY{p}{,} \PY{l+s+s2}{\PYZdq{}}\PY{l+s+s2}{ann\PYZus{}vol}\PY{l+s+s2}{\PYZdq{}}\PY{p}{,} \PY{l+s+s2}{\PYZdq{}}\PY{l+s+s2}{sharpe}\PY{l+s+s2}{\PYZdq{}}\PY{p}{]}\PY{p}{)}
\PY{n}{tan\PYZus{}noshort\PYZus{}stats\PYZus{}df} \PY{o}{=} \PY{n}{pd}\PY{o}{.}\PY{n}{DataFrame}\PY{p}{(}\PY{n}{tan\PYZus{}noshort\PYZus{}stats}\PY{p}{,} 
                                    \PY{n}{index}\PY{o}{=}\PY{n}{window\PYZus{}labels}\PY{p}{,} 
                                    \PY{n}{columns}\PY{o}{=}\PY{p}{[}\PY{l+s+s2}{\PYZdq{}}\PY{l+s+s2}{ann\PYZus{}return}\PY{l+s+s2}{\PYZdq{}}\PY{p}{,} \PY{l+s+s2}{\PYZdq{}}\PY{l+s+s2}{ann\PYZus{}vol}\PY{l+s+s2}{\PYZdq{}}\PY{p}{,} \PY{l+s+s2}{\PYZdq{}}\PY{l+s+s2}{sharpe}\PY{l+s+s2}{\PYZdq{}}\PY{p}{]}\PY{p}{)}

\PY{n}{tan\PYZus{}uncon\PYZus{}weights\PYZus{}df}\PY{p}{,} \PY{n}{tan\PYZus{}noshort\PYZus{}weights\PYZus{}df}
\PY{n}{tan\PYZus{}uncon\PYZus{}weights\PYZus{}df}\PY{o}{.}\PY{n}{round}\PY{p}{(}\PY{l+m+mi}{3}\PY{p}{)}\PY{p}{,} \PY{n}{tan\PYZus{}noshort\PYZus{}weights\PYZus{}df}\PY{o}{.}\PY{n}{round}\PY{p}{(}\PY{l+m+mi}{3}\PY{p}{)}
\end{Verbatim}
\end{tcolorbox}

            \begin{tcolorbox}[breakable, size=fbox, boxrule=.5pt, pad at break*=1mm, opacityfill=0]
\begin{Verbatim}[commandchars=\\\{\}]
(             TSX    CAC    DAX  Eurostoxx50  NIKKEI225  FTSE100  SP500  \textbackslash{}
 2010-2014 -0.448 -1.038  2.196       -1.346      0.210    0.258  1.707
 2011-2015 -2.601  0.224  3.181       -3.174      0.771    0.218  3.707
 2012-2016 -0.740  2.337  2.649       -4.569      0.305   -0.524  1.768
 2013-2017 -0.556  2.431  1.987       -4.182      0.253   -0.319  1.488
 2014-2018 -2.489  9.735  4.944      -14.201      0.648   -1.118  3.180
 2015-2019 -1.101  6.911  2.983       -9.414      0.546   -0.632  1.455
 2016-2020 -0.440  6.160  4.921      -10.439      0.833   -1.396  1.139
 2017-2021 -0.550  5.453  2.211       -6.922      0.633   -1.006  1.599

            IBOVESPA
 2010-2014    -0.538
 2011-2015    -1.326
 2012-2016    -0.226
 2013-2017    -0.102
 2014-2018     0.300
 2015-2019     0.253
 2016-2020     0.222
 2017-2021    -0.418  ,
            TSX  CAC  DAX  Eurostoxx50  NIKKEI225  FTSE100  SP500  IBOVESPA
 2010-2014  0.0 -0.0  0.0          0.0      0.123      0.0  0.877    -0.000
 2011-2015  0.0 -0.0  0.0          0.0      0.116      0.0  0.884    -0.000
 2012-2016 -0.0  0.0  0.0          0.0      0.114      0.0  0.886     0.000
 2013-2017  0.0  0.0  0.0         -0.0      0.120      0.0  0.880     0.000
 2014-2018  0.0  0.0  0.0         -0.0      0.066      0.0  0.934     0.000
 2015-2019  0.0  0.0 -0.0          0.0      0.302      0.0  0.693     0.004
 2016-2020 -0.0  0.0 -0.0          0.0      0.455     -0.0  0.545     0.000
 2017-2021 -0.0  0.0  0.0          0.0      0.312      0.0  0.688     0.000)
\end{Verbatim}
\end{tcolorbox}
        
    The two tables show how \textbf{tangency portfolio weights} behave in
each 5-year window:

\begin{itemize}
\item
  \textbf{Unconstrained tangency (first table):}

  Weights are huge (far from ±1) and often negative. For example, in
  2014--2018 the portfolio is very long CAC and DAX, very short
  Eurostoxx50, and strongly long SP500. This means the portfolio is
  \textbf{heavily levered and very unstable over time}: small changes in
  inputs lead to completely different long--short bets from one window
  to the next.
\item
  \textbf{No-short tangency (second table):}

  Almost all the weight goes systematically to \textbf{SP500}, with a
  smaller share in \textbf{NIKKEI225}. The other indices are essentially
  zero in every window (the ``--0.0'' values are just rounding). These
  constrained portfolios are \textbf{much more stable but very
  concentrated} in a couple of indices.
\end{itemize}

So, rolling the MV analysis through time clearly shows the trade-off:\\
unconstrained MV gives extreme, unstable long--short portfolios, while
adding a no-short constraint produces more realistic, but highly
concentrated, allocations.

    Before interpreting the no-short portfolios, we remove tiny numerical
noise (weights of order \(10^{-16}\)) by setting very small coefficients
to zero and rebuilding a cleaned DataFrame of long-only weights. This
way, zeros really mean ``no position'' and the remaining weights are
easier to interpret.

    \begin{tcolorbox}[breakable, size=fbox, boxrule=1pt, pad at break*=1mm,colback=cellbackground, colframe=cellborder]
\begin{Verbatim}[commandchars=\\\{\}]
\PY{c+c1}{\PYZsh{} Clean no\PYZhy{}short weights: we just set tiny numerical values to zero}
\PY{n}{tan\PYZus{}noshort\PYZus{}weights\PYZus{}clean} \PY{o}{=} \PY{n}{np}\PY{o}{.}\PY{n}{where}\PY{p}{(}\PY{n}{np}\PY{o}{.}\PY{n}{abs}\PY{p}{(}\PY{n}{tan\PYZus{}noshort\PYZus{}weights\PYZus{}df}\PY{o}{.}\PY{n}{values}\PY{p}{)} \PY{o}{\PYZlt{}} \PY{l+m+mf}{1e\PYZhy{}10}\PY{p}{,}
                                     \PY{l+m+mi}{0}\PY{p}{,}
                                     \PY{n}{tan\PYZus{}noshort\PYZus{}weights\PYZus{}df}\PY{o}{.}\PY{n}{values}\PY{p}{)}

\PY{n}{tan\PYZus{}noshort\PYZus{}weights\PYZus{}clean\PYZus{}df} \PY{o}{=} \PY{n}{pd}\PY{o}{.}\PY{n}{DataFrame}\PY{p}{(}\PY{n}{tan\PYZus{}noshort\PYZus{}weights\PYZus{}clean}\PY{p}{,}
                                            \PY{n}{index}\PY{o}{=}\PY{n}{tan\PYZus{}noshort\PYZus{}weights\PYZus{}df}\PY{o}{.}\PY{n}{index}\PY{p}{,}
                                            \PY{n}{columns}\PY{o}{=}\PY{n}{tan\PYZus{}noshort\PYZus{}weights\PYZus{}df}\PY{o}{.}\PY{n}{columns}\PY{p}{)}

\PY{n}{tan\PYZus{}noshort\PYZus{}weights\PYZus{}clean\PYZus{}df}
\end{Verbatim}
\end{tcolorbox}

            \begin{tcolorbox}[breakable, size=fbox, boxrule=.5pt, pad at break*=1mm, opacityfill=0]
\begin{Verbatim}[commandchars=\\\{\}]
           TSX  CAC  DAX  Eurostoxx50  NIKKEI225  FTSE100     SP500  IBOVESPA
2010-2014  0.0  0.0  0.0          0.0   0.123354      0.0  0.876646  0.000000
2011-2015  0.0  0.0  0.0          0.0   0.115577      0.0  0.884423  0.000000
2012-2016  0.0  0.0  0.0          0.0   0.114084      0.0  0.885916  0.000000
2013-2017  0.0  0.0  0.0          0.0   0.120050      0.0  0.879950  0.000000
2014-2018  0.0  0.0  0.0          0.0   0.065745      0.0  0.934255  0.000000
2015-2019  0.0  0.0  0.0          0.0   0.302376      0.0  0.693289  0.004335
2016-2020  0.0  0.0  0.0          0.0   0.454591      0.0  0.545220  0.000189
2017-2021  0.0  0.0  0.0          0.0   0.311639      0.0  0.688361  0.000000
\end{Verbatim}
\end{tcolorbox}
        
    The cleaned no-short tangency weights over the rolling 5-year windows
show a very clear pattern:

\begin{itemize}
\tightlist
\item
  In \textbf{every window}, the portfolio is almost entirely split
  between \textbf{SP500} and \textbf{NIKKEI225}. The weight on SP500 is
  always between about \(0.55\) and \(0.93\), and the rest goes to
  NIKKEI225.
\item
  All the other indices (TSX, CAC, DAX, Eurostoxx50, FTSE100)
  consistently have \textbf{weight \(0\)}.
\item
  IBOVESPA only gets a \textbf{tiny} positive weight in 2015--2019 and
  2016--2020 (around \(0.4\%\) and \(0.02\%\)), which is basically
  negligible.
\end{itemize}

So the long-only tangency portfolio is very stable over time and heavily
concentrated in SP500 (with NIKKEI225 as a smaller satellite position),
while the other indices are never attractive enough in terms of
risk--return to enter the optimal MV portfolio under the no-short
constraint.

    To visualise how unstable the mean--variance solution is over time, we
now track the weight of a single asset - the S\&P 500 - across all
5-year windows. The plot below shows its weight in the unconstrained
tangency portfolio and in the no-short tangency portfolio, so we can see
how leverage and short-selling change the behaviour of the optimiser
through time.

    \begin{tcolorbox}[breakable, size=fbox, boxrule=1pt, pad at break*=1mm,colback=cellbackground, colframe=cellborder]
\begin{Verbatim}[commandchars=\\\{\}]
\PY{n}{plt}\PY{o}{.}\PY{n}{figure}\PY{p}{(}\PY{n}{figsize}\PY{o}{=}\PY{p}{(}\PY{l+m+mi}{8}\PY{p}{,} \PY{l+m+mi}{5}\PY{p}{)}\PY{p}{)}

\PY{c+c1}{\PYZsh{} SP500 weight in unconstrained tangency portfolio}
\PY{n}{plt}\PY{o}{.}\PY{n}{plot}\PY{p}{(}
    \PY{n}{tan\PYZus{}uncon\PYZus{}weights\PYZus{}df}\PY{o}{.}\PY{n}{index}\PY{p}{,}
    \PY{n}{tan\PYZus{}uncon\PYZus{}weights\PYZus{}df}\PY{p}{[}\PY{l+s+s2}{\PYZdq{}}\PY{l+s+s2}{SP500}\PY{l+s+s2}{\PYZdq{}}\PY{p}{]}\PY{p}{,}
    \PY{n}{marker}\PY{o}{=}\PY{l+s+s2}{\PYZdq{}}\PY{l+s+s2}{o}\PY{l+s+s2}{\PYZdq{}}\PY{p}{,}
    \PY{n}{linestyle}\PY{o}{=}\PY{l+s+s2}{\PYZdq{}}\PY{l+s+s2}{\PYZhy{}}\PY{l+s+s2}{\PYZdq{}}\PY{p}{,}
    \PY{n}{label}\PY{o}{=}\PY{l+s+s2}{\PYZdq{}}\PY{l+s+s2}{Tangency unconstrained}\PY{l+s+s2}{\PYZdq{}}
\PY{p}{)}

\PY{c+c1}{\PYZsh{} SP500 weight in no\PYZhy{}short tangency portfolio}
\PY{n}{plt}\PY{o}{.}\PY{n}{plot}\PY{p}{(}
    \PY{n}{tan\PYZus{}noshort\PYZus{}weights\PYZus{}clean\PYZus{}df}\PY{o}{.}\PY{n}{index}\PY{p}{,}
    \PY{n}{tan\PYZus{}noshort\PYZus{}weights\PYZus{}clean\PYZus{}df}\PY{p}{[}\PY{l+s+s2}{\PYZdq{}}\PY{l+s+s2}{SP500}\PY{l+s+s2}{\PYZdq{}}\PY{p}{]}\PY{p}{,}
    \PY{n}{marker}\PY{o}{=}\PY{l+s+s2}{\PYZdq{}}\PY{l+s+s2}{o}\PY{l+s+s2}{\PYZdq{}}\PY{p}{,}
    \PY{n}{linestyle}\PY{o}{=}\PY{l+s+s2}{\PYZdq{}}\PY{l+s+s2}{\PYZhy{}}\PY{l+s+s2}{\PYZdq{}}\PY{p}{,}
    \PY{n}{label}\PY{o}{=}\PY{l+s+s2}{\PYZdq{}}\PY{l+s+s2}{Tangency no\PYZhy{}short}\PY{l+s+s2}{\PYZdq{}}
\PY{p}{)}

\PY{n}{plt}\PY{o}{.}\PY{n}{title}\PY{p}{(}\PY{l+s+s2}{\PYZdq{}}\PY{l+s+s2}{SP500 weight in tangency portfolio (5\PYZhy{}year rolling windows)}\PY{l+s+s2}{\PYZdq{}}\PY{p}{)}
\PY{n}{plt}\PY{o}{.}\PY{n}{xlabel}\PY{p}{(}\PY{l+s+s2}{\PYZdq{}}\PY{l+s+s2}{Window}\PY{l+s+s2}{\PYZdq{}}\PY{p}{)}
\PY{n}{plt}\PY{o}{.}\PY{n}{ylabel}\PY{p}{(}\PY{l+s+s2}{\PYZdq{}}\PY{l+s+s2}{Portfolio weight}\PY{l+s+s2}{\PYZdq{}}\PY{p}{)}
\PY{n}{plt}\PY{o}{.}\PY{n}{xticks}\PY{p}{(}\PY{n}{rotation}\PY{o}{=}\PY{l+m+mi}{45}\PY{p}{)}
\PY{n}{plt}\PY{o}{.}\PY{n}{legend}\PY{p}{(}\PY{p}{)}
\PY{n}{plt}\PY{o}{.}\PY{n}{grid}\PY{p}{(}\PY{n}{alpha}\PY{o}{=}\PY{l+m+mf}{0.3}\PY{p}{)}
\PY{n}{plt}\PY{o}{.}\PY{n}{tight\PYZus{}layout}\PY{p}{(}\PY{p}{)}
\PY{n}{plt}\PY{o}{.}\PY{n}{show}\PY{p}{(}\PY{p}{)}
\end{Verbatim}
\end{tcolorbox}

    \begin{center}
    \adjustimage{max size={0.9\linewidth}{0.9\paperheight}}{Work-7_files/Work-7_100_0.png}
    \end{center}
    { \hspace*{\fill} \\}
    
    The figure shows that the SP500 weight in the \textbf{unconstrained}
tangency portfolio is very unstable and often well above (1). This means
the optimiser is using leverage and short positions in the other indices
to finance an oversized long position in the SP500.

By contrast, in the \textbf{no-short} tangency portfolio the SP500
weight stays between roughly (0.55) and (0.95) and moves much more
smoothly across windows. The rolling exercise therefore confirms that
the standard Markowitz solution is highly sample-sensitive and levered,
whereas the long-only solution remains strongly tilted to the SP500 but
in a more stable and realistic way.

    \hypertarget{comparison-with-market-cap-based-weights}{%
\subsubsection{3.3 Comparison with market cap-based
weights}\label{comparison-with-market-cap-based-weights}}

Finally, we compare our mean--variance tangency portfolios with a more
realistic benchmark: a portfolio weighted by the market capitalization
of each index. The Excel file provides these market-cap shares in the
sheet \texttt{market\_cap}. We read these weights, align them with our
eight indices, and compare them to the tangency portfolios in the last
5-year window (2017-2021).

    \begin{tcolorbox}[breakable, size=fbox, boxrule=1pt, pad at break*=1mm,colback=cellbackground, colframe=cellborder]
\begin{Verbatim}[commandchars=\\\{\}]
\PY{c+c1}{\PYZsh{} Read the market\PYZus{}cap sheet}
\PY{n}{mkt\PYZus{}raw} \PY{o}{=} \PY{n}{pd}\PY{o}{.}\PY{n}{read\PYZus{}excel}\PY{p}{(}\PY{n}{FILE}\PY{p}{,} \PY{n}{sheet\PYZus{}name}\PY{o}{=}\PY{l+s+s2}{\PYZdq{}}\PY{l+s+s2}{market\PYZus{}cap}\PY{l+s+s2}{\PYZdq{}}\PY{p}{,} \PY{n}{header}\PY{o}{=}\PY{l+m+mi}{1}\PY{p}{)}
\PY{n}{mkt\PYZus{}raw}\PY{o}{.}\PY{n}{columns} \PY{o}{=} \PY{p}{[}\PY{n+nb}{str}\PY{p}{(}\PY{n}{c}\PY{p}{)}\PY{o}{.}\PY{n}{strip}\PY{p}{(}\PY{p}{)} \PY{k}{for} \PY{n}{c} \PY{o+ow}{in} \PY{n}{mkt\PYZus{}raw}\PY{o}{.}\PY{n}{columns}\PY{p}{]}

\PY{n}{mkt\PYZus{}raw}
\end{Verbatim}
\end{tcolorbox}

            \begin{tcolorbox}[breakable, size=fbox, boxrule=.5pt, pad at break*=1mm, opacityfill=0]
\begin{Verbatim}[commandchars=\\\{\}]
   Unnamed: 0  SPTSX INDEX  CAC INDEX  DAX INDEX  SX5E INDEX  NKY INDEX  \textbackslash{}
0         NaN     0.041834    0.04318   0.036197    0.090955   0.092523

   UKX INDEX  SPX INDEX  IBOV INDEX
0   0.081589   0.602344    0.011378
\end{Verbatim}
\end{tcolorbox}
        
    The \texttt{market\_cap} sheet gives us one row of benchmark weights by
index. It shows a very concentrated universe: the S\&P 500 alone
represents about \textbf{60\%} of total market cap, while the other
developed markets are each in the single digits, and Brazil is just a
bit above \textbf{1\%}.

Next, we rename and reorder these entries so that this market-cap
portfolio can be directly compared to our mean-variance tangency
portfolios.

    \begin{tcolorbox}[breakable, size=fbox, boxrule=1pt, pad at break*=1mm,colback=cellbackground, colframe=cellborder]
\begin{Verbatim}[commandchars=\\\{\}]
\PY{c+c1}{\PYZsh{} Take the first (and only) row of weights}
\PY{n}{mkt\PYZus{}row} \PY{o}{=} \PY{n}{mkt\PYZus{}raw}\PY{o}{.}\PY{n}{iloc}\PY{p}{[}\PY{l+m+mi}{0}\PY{p}{]}\PY{o}{.}\PY{n}{drop}\PY{p}{(}\PY{n}{labels}\PY{o}{=}\PY{p}{[}\PY{l+s+s2}{\PYZdq{}}\PY{l+s+s2}{Unnamed: 0}\PY{l+s+s2}{\PYZdq{}}\PY{p}{]}\PY{p}{)}

\PY{c+c1}{\PYZsh{} We map Excel column names to our index names}
\PY{n}{mapping} \PY{o}{=} \PY{p}{\PYZob{}}
    \PY{l+s+s2}{\PYZdq{}}\PY{l+s+s2}{SPTSX INDEX}\PY{l+s+s2}{\PYZdq{}}\PY{p}{:} \PY{l+s+s2}{\PYZdq{}}\PY{l+s+s2}{TSX}\PY{l+s+s2}{\PYZdq{}}\PY{p}{,}
    \PY{l+s+s2}{\PYZdq{}}\PY{l+s+s2}{CAC INDEX}\PY{l+s+s2}{\PYZdq{}}\PY{p}{:} \PY{l+s+s2}{\PYZdq{}}\PY{l+s+s2}{CAC}\PY{l+s+s2}{\PYZdq{}}\PY{p}{,}
    \PY{l+s+s2}{\PYZdq{}}\PY{l+s+s2}{DAX INDEX}\PY{l+s+s2}{\PYZdq{}}\PY{p}{:} \PY{l+s+s2}{\PYZdq{}}\PY{l+s+s2}{DAX}\PY{l+s+s2}{\PYZdq{}}\PY{p}{,}
    \PY{l+s+s2}{\PYZdq{}}\PY{l+s+s2}{SX5E INDEX}\PY{l+s+s2}{\PYZdq{}}\PY{p}{:} \PY{l+s+s2}{\PYZdq{}}\PY{l+s+s2}{Eurostoxx50}\PY{l+s+s2}{\PYZdq{}}\PY{p}{,}
    \PY{l+s+s2}{\PYZdq{}}\PY{l+s+s2}{NKY INDEX}\PY{l+s+s2}{\PYZdq{}}\PY{p}{:} \PY{l+s+s2}{\PYZdq{}}\PY{l+s+s2}{NIKKEI225}\PY{l+s+s2}{\PYZdq{}}\PY{p}{,}
    \PY{l+s+s2}{\PYZdq{}}\PY{l+s+s2}{UKX INDEX}\PY{l+s+s2}{\PYZdq{}}\PY{p}{:} \PY{l+s+s2}{\PYZdq{}}\PY{l+s+s2}{FTSE100}\PY{l+s+s2}{\PYZdq{}}\PY{p}{,}
    \PY{l+s+s2}{\PYZdq{}}\PY{l+s+s2}{SPX INDEX}\PY{l+s+s2}{\PYZdq{}}\PY{p}{:} \PY{l+s+s2}{\PYZdq{}}\PY{l+s+s2}{SP500}\PY{l+s+s2}{\PYZdq{}}\PY{p}{,}
    \PY{l+s+s2}{\PYZdq{}}\PY{l+s+s2}{IBOV INDEX}\PY{l+s+s2}{\PYZdq{}}\PY{p}{:} \PY{l+s+s2}{\PYZdq{}}\PY{l+s+s2}{IBOVESPA}\PY{l+s+s2}{\PYZdq{}}\PY{p}{,}
\PY{p}{\PYZcb{}}

\PY{n}{mkt\PYZus{}weights} \PY{o}{=} \PY{n}{mkt\PYZus{}row}\PY{o}{.}\PY{n}{rename}\PY{p}{(}\PY{n}{index}\PY{o}{=}\PY{n}{mapping}\PY{p}{)}

\PY{c+c1}{\PYZsh{} We put it in order to be in the same order as \PYZdq{}asset\PYZdq{}}
\PY{n}{mkt\PYZus{}weights} \PY{o}{=} \PY{n}{mkt\PYZus{}weights}\PY{o}{.}\PY{n}{reindex}\PY{p}{(}\PY{n}{assets}\PY{p}{)}

\PY{c+c1}{\PYZsh{} We normalize just in case...}
\PY{n}{mkt\PYZus{}weights} \PY{o}{=} \PY{n}{mkt\PYZus{}weights} \PY{o}{/} \PY{n}{mkt\PYZus{}weights}\PY{o}{.}\PY{n}{sum}\PY{p}{(}\PY{p}{)}

\PY{n}{mkt\PYZus{}weights}
\end{Verbatim}
\end{tcolorbox}

            \begin{tcolorbox}[breakable, size=fbox, boxrule=.5pt, pad at break*=1mm, opacityfill=0]
\begin{Verbatim}[commandchars=\\\{\}]
TSX            0.041834
CAC            0.043180
DAX            0.036197
Eurostoxx50    0.090955
NIKKEI225      0.092523
FTSE100        0.081589
SP500          0.602344
IBOVESPA       0.011378
Name: 0, dtype: float64
\end{Verbatim}
\end{tcolorbox}
        
    Re-expressing the Excel row as a \texttt{pandas} Series, we now have the
\textbf{market-cap weights} in exactly the same format as our other
portfolios. The picture is very skewed: the S\&P 500 still dominates
with about \textbf{60\%} of total weight, the main developed markets
(TSX, CAC, DAX, Eurostoxx50, Nikkei, FTSE100) each get only 4--9\%, and
\textbf{IBOVESPA} stays tiny at just above \textbf{1\%}.

This Series is what we will now line up against the tangency portfolio
weights.

    To make the comparison concrete, we now focus on the last 5-year window
(2017-2021). We extract the tangency weights for that window, both with
and without short-selling, and place them side by side with the
market-cap weights in a single DataFrame. This will let us see, index by
index, how far the mean--variance portfolios are from a simple
market-cap benchmark.

    \begin{tcolorbox}[breakable, size=fbox, boxrule=1pt, pad at break*=1mm,colback=cellbackground, colframe=cellborder]
\begin{Verbatim}[commandchars=\\\{\}]
\PY{c+c1}{\PYZsh{} Tangency weights in the last 5\PYZhy{}year window (2017–2021)}
\PY{n}{w\PYZus{}mv\PYZus{}uncon\PYZus{}last}   \PY{o}{=} \PY{n}{tan\PYZus{}uncon\PYZus{}weights\PYZus{}df}\PY{o}{.}\PY{n}{loc}\PY{p}{[}\PY{l+s+s2}{\PYZdq{}}\PY{l+s+s2}{2017\PYZhy{}2021}\PY{l+s+s2}{\PYZdq{}}\PY{p}{]}
\PY{n}{w\PYZus{}mv\PYZus{}noshort\PYZus{}last} \PY{o}{=} \PY{n}{tan\PYZus{}noshort\PYZus{}weights\PYZus{}clean\PYZus{}df}\PY{o}{.}\PY{n}{loc}\PY{p}{[}\PY{l+s+s2}{\PYZdq{}}\PY{l+s+s2}{2017\PYZhy{}2021}\PY{l+s+s2}{\PYZdq{}}\PY{p}{]}

\PY{c+c1}{\PYZsh{} Comparison with market\PYZhy{}cap weights}
\PY{n}{compare\PYZus{}weights} \PY{o}{=} \PY{n}{pd}\PY{o}{.}\PY{n}{DataFrame}\PY{p}{(}\PY{p}{\PYZob{}}
    \PY{l+s+s2}{\PYZdq{}}\PY{l+s+s2}{MarketCap}\PY{l+s+s2}{\PYZdq{}}\PY{p}{:} \PY{n}{mkt\PYZus{}weights}\PY{p}{,}  \PY{c+c1}{\PYZsh{} serie built from market\PYZhy{}cap}
    \PY{l+s+s2}{\PYZdq{}}\PY{l+s+s2}{MV\PYZus{}Tangency\PYZus{}Uncon\PYZus{}2017\PYZus{}2021}\PY{l+s+s2}{\PYZdq{}}\PY{p}{:} \PY{n}{w\PYZus{}mv\PYZus{}uncon\PYZus{}last}\PY{p}{,}
    \PY{l+s+s2}{\PYZdq{}}\PY{l+s+s2}{MV\PYZus{}Tangency\PYZus{}NoShort\PYZus{}2017\PYZus{}2021}\PY{l+s+s2}{\PYZdq{}}\PY{p}{:} \PY{n}{w\PYZus{}mv\PYZus{}noshort\PYZus{}last}\PY{p}{,}
\PY{p}{\PYZcb{}}\PY{p}{,} \PY{n}{index}\PY{o}{=}\PY{n}{assets}\PY{p}{)}

\PY{n}{compare\PYZus{}weights}
\end{Verbatim}
\end{tcolorbox}

            \begin{tcolorbox}[breakable, size=fbox, boxrule=.5pt, pad at break*=1mm, opacityfill=0]
\begin{Verbatim}[commandchars=\\\{\}]
             MarketCap  MV\_Tangency\_Uncon\_2017\_2021  \textbackslash{}
TSX           0.041834                    -0.549505
CAC           0.043180                     5.453498
DAX           0.036197                     2.210759
Eurostoxx50   0.090955                    -6.922410
NIKKEI225     0.092523                     0.632892
FTSE100       0.081589                    -1.006031
SP500         0.602344                     1.599128
IBOVESPA      0.011378                    -0.418332

             MV\_Tangency\_NoShort\_2017\_2021
TSX                               0.000000
CAC                               0.000000
DAX                               0.000000
Eurostoxx50                       0.000000
NIKKEI225                         0.311639
FTSE100                           0.000000
SP500                             0.688361
IBOVESPA                          0.000000
\end{Verbatim}
\end{tcolorbox}
        
    The table compares three sets of weights for the 2017-2021 window:

\begin{itemize}
\item
  The \textbf{market-cap portfolio} is diversified across all indices
  but heavily tilted to the US: about 60\% in the S\&P 500, with most
  other indices between roughly 3--9\% and a tiny allocation to Brazil.
\item
  The \textbf{unconstrained tangency portfolio} is extremely aggressive.
  It takes very large leveraged longs in \textbf{CAC}, \textbf{DAX} and
  \textbf{SP500}, financed by big short positions in
  \textbf{Eurostoxx50}, \textbf{FTSE100}, \textbf{TSX} and
  \textbf{IBOVESPA}. This is the classic Markowitz pathology when we
  allow unlimited short-selling and rely on noisy sample means.
\item
  The \textbf{no short-sale tangency portfolio} is much more realistic:
  it puts about 69\% in \textbf{SP500} and 31\% in \textbf{NIKKEI225},
  with zero weight on the other indices. Compared to the market-cap
  benchmark it is still very concentrated (mainly US + Japan), but it
  avoids the huge short positions and leverage of the unconstrained
  solution.
\end{itemize}

    To visualise these differences more clearly, we now plot the three sets
of weights (market-cap, unconstrained tangency, and no-short tangency)
for the 2017-2021 window on a single bar chart. This makes it easier to
see, for each index, how far the mean-variance portfolios move away from
the simple market-cap allocation.

    \begin{tcolorbox}[breakable, size=fbox, boxrule=1pt, pad at break*=1mm,colback=cellbackground, colframe=cellborder]
\begin{Verbatim}[commandchars=\\\{\}]
\PY{n}{plt}\PY{o}{.}\PY{n}{figure}\PY{p}{(}\PY{n}{figsize}\PY{o}{=}\PY{p}{(}\PY{l+m+mi}{10}\PY{p}{,}\PY{l+m+mi}{6}\PY{p}{)}\PY{p}{)}

\PY{n}{x} \PY{o}{=} \PY{n}{np}\PY{o}{.}\PY{n}{arange}\PY{p}{(}\PY{n+nb}{len}\PY{p}{(}\PY{n}{assets}\PY{p}{)}\PY{p}{)}
\PY{n}{width} \PY{o}{=} \PY{l+m+mf}{0.25}

\PY{n}{plt}\PY{o}{.}\PY{n}{bar}\PY{p}{(}\PY{n}{x} \PY{o}{\PYZhy{}} \PY{n}{width}\PY{p}{,} \PY{n}{compare\PYZus{}weights}\PY{p}{[}\PY{l+s+s2}{\PYZdq{}}\PY{l+s+s2}{MarketCap}\PY{l+s+s2}{\PYZdq{}}\PY{p}{]}\PY{p}{,} 
        \PY{n}{width}\PY{o}{=}\PY{n}{width}\PY{p}{,} \PY{n}{label}\PY{o}{=}\PY{l+s+s2}{\PYZdq{}}\PY{l+s+s2}{Market\PYZhy{}cap}\PY{l+s+s2}{\PYZdq{}}\PY{p}{)}
\PY{n}{plt}\PY{o}{.}\PY{n}{bar}\PY{p}{(}\PY{n}{x}\PY{p}{,} \PY{n}{compare\PYZus{}weights}\PY{p}{[}\PY{l+s+s2}{\PYZdq{}}\PY{l+s+s2}{MV\PYZus{}Tangency\PYZus{}Uncon\PYZus{}2017\PYZus{}2021}\PY{l+s+s2}{\PYZdq{}}\PY{p}{]}\PY{p}{,} 
        \PY{n}{width}\PY{o}{=}\PY{n}{width}\PY{p}{,} \PY{n}{label}\PY{o}{=}\PY{l+s+s2}{\PYZdq{}}\PY{l+s+s2}{MV tangency (unconstrained)}\PY{l+s+s2}{\PYZdq{}}\PY{p}{)}
\PY{n}{plt}\PY{o}{.}\PY{n}{bar}\PY{p}{(}\PY{n}{x} \PY{o}{+} \PY{n}{width}\PY{p}{,} \PY{n}{compare\PYZus{}weights}\PY{p}{[}\PY{l+s+s2}{\PYZdq{}}\PY{l+s+s2}{MV\PYZus{}Tangency\PYZus{}NoShort\PYZus{}2017\PYZus{}2021}\PY{l+s+s2}{\PYZdq{}}\PY{p}{]}\PY{p}{,} 
        \PY{n}{width}\PY{o}{=}\PY{n}{width}\PY{p}{,} \PY{n}{label}\PY{o}{=}\PY{l+s+s2}{\PYZdq{}}\PY{l+s+s2}{MV tangency (no\PYZhy{}short)}\PY{l+s+s2}{\PYZdq{}}\PY{p}{)}

\PY{n}{plt}\PY{o}{.}\PY{n}{xticks}\PY{p}{(}\PY{n}{x}\PY{p}{,} \PY{n}{assets}\PY{p}{,} \PY{n}{rotation}\PY{o}{=}\PY{l+m+mi}{45}\PY{p}{)}
\PY{n}{plt}\PY{o}{.}\PY{n}{ylabel}\PY{p}{(}\PY{l+s+s2}{\PYZdq{}}\PY{l+s+s2}{Portfolio weight}\PY{l+s+s2}{\PYZdq{}}\PY{p}{)}
\PY{n}{plt}\PY{o}{.}\PY{n}{title}\PY{p}{(}\PY{l+s+s2}{\PYZdq{}}\PY{l+s+s2}{Market\PYZhy{}cap vs MV tangency portfolios (2017–2021)}\PY{l+s+s2}{\PYZdq{}}\PY{p}{)}
\PY{n}{plt}\PY{o}{.}\PY{n}{legend}\PY{p}{(}\PY{p}{)}
\PY{n}{plt}\PY{o}{.}\PY{n}{tight\PYZus{}layout}\PY{p}{(}\PY{p}{)}
\PY{n}{plt}\PY{o}{.}\PY{n}{show}\PY{p}{(}\PY{p}{)}
\end{Verbatim}
\end{tcolorbox}

    \begin{center}
    \adjustimage{max size={0.9\linewidth}{0.9\paperheight}}{Work-7_files/Work-7_111_0.png}
    \end{center}
    { \hspace*{\fill} \\}
    
    The bar chart makes the comparison very clear. The \textbf{market-cap}
portfolio is well diversified across all indices, with a dominant but
not extreme weight on the S\&P 500. By contrast, the
\textbf{unconstrained tangency} portfolio loads huge long and short
positions on a few indices (large long in CAC and DAX, very large short
in Eurostoxx50, etc.), which is completely unrealistic from a practical
point of view.

The \textbf{no-short tangency} portfolio sits much closer to the
market-cap benchmark: it is still tilted towards the S\&P 500 and
NIKKEI225, but without the extreme leverage and short positions. This
nicely illustrates how adding a simple non-negativity constraint can
turn the theoretical Markowitz solution into something that looks much
more like a realistic equity allocation.

    \hypertarget{sharpe-ratios-over-time}{%
\subsubsection{3.4 Sharpe ratios over
time}\label{sharpe-ratios-over-time}}

So far we have focused on the portfolio \textbf{weights}. To see how
meaningful these portfolios actually are, it is also useful to track
their \textbf{Sharpe ratios} over time. For each 5-year window we
already computed the tangency portfolio under:

\begin{itemize}
\tightlist
\item
  the \textbf{unconstrained} Markowitz model (short-selling allowed),
  and\\
\item
  the \textbf{no short-sale} version (weights constrained to be
  non-negative).
\end{itemize}

In this section we plot the Sharpe ratio of both tangency portfolios
across all rolling windows. This shows whether the very extreme
unconstrained portfolios actually deliver a consistently higher
risk-adjusted performance, or whether the more realistic no-short
portfolios achieve similar Sharpe ratios with much less leverage and
concentration.

    \begin{tcolorbox}[breakable, size=fbox, boxrule=1pt, pad at break*=1mm,colback=cellbackground, colframe=cellborder]
\begin{Verbatim}[commandchars=\\\{\}]
\PY{c+c1}{\PYZsh{} Collect Sharpe ratios for each 5\PYZhy{}year window}
\PY{n}{sharpe\PYZus{}df} \PY{o}{=} \PY{n}{pd}\PY{o}{.}\PY{n}{DataFrame}\PY{p}{(}\PY{p}{\PYZob{}}
    \PY{l+s+s2}{\PYZdq{}}\PY{l+s+s2}{Unconstrained}\PY{l+s+s2}{\PYZdq{}}\PY{p}{:} \PY{n}{tan\PYZus{}uncon\PYZus{}stats\PYZus{}df}\PY{p}{[}\PY{l+s+s2}{\PYZdq{}}\PY{l+s+s2}{sharpe}\PY{l+s+s2}{\PYZdq{}}\PY{p}{]}\PY{p}{,}
    \PY{l+s+s2}{\PYZdq{}}\PY{l+s+s2}{NoShort}\PY{l+s+s2}{\PYZdq{}}\PY{p}{:} \PY{n}{tan\PYZus{}noshort\PYZus{}stats\PYZus{}df}\PY{p}{[}\PY{l+s+s2}{\PYZdq{}}\PY{l+s+s2}{sharpe}\PY{l+s+s2}{\PYZdq{}}\PY{p}{]}\PY{p}{,}
\PY{p}{\PYZcb{}}\PY{p}{,} \PY{n}{index}\PY{o}{=}\PY{n}{window\PYZus{}labels}\PY{p}{)}

\PY{n}{sharpe\PYZus{}df}
\end{Verbatim}
\end{tcolorbox}

            \begin{tcolorbox}[breakable, size=fbox, boxrule=.5pt, pad at break*=1mm, opacityfill=0]
\begin{Verbatim}[commandchars=\\\{\}]
           Unconstrained   NoShort
2010-2014       1.995667  0.834794
2011-2015       2.061726  0.696750
2012-2016       2.053157  0.988011
2013-2017       2.086428  1.199801
2014-2018       1.484669  0.499447
2015-2019       1.731560  0.789048
2016-2020       1.519060  0.779839
2017-2021       1.534876  0.894899
\end{Verbatim}
\end{tcolorbox}
        
    \begin{tcolorbox}[breakable, size=fbox, boxrule=1pt, pad at break*=1mm,colback=cellbackground, colframe=cellborder]
\begin{Verbatim}[commandchars=\\\{\}]
\PY{n}{plt}\PY{o}{.}\PY{n}{figure}\PY{p}{(}\PY{n}{figsize}\PY{o}{=}\PY{p}{(}\PY{l+m+mi}{8}\PY{p}{,} \PY{l+m+mi}{5}\PY{p}{)}\PY{p}{)}

\PY{n}{plt}\PY{o}{.}\PY{n}{plot}\PY{p}{(}\PY{n}{sharpe\PYZus{}df}\PY{o}{.}\PY{n}{index}\PY{p}{,}
         \PY{n}{sharpe\PYZus{}df}\PY{p}{[}\PY{l+s+s2}{\PYZdq{}}\PY{l+s+s2}{Unconstrained}\PY{l+s+s2}{\PYZdq{}}\PY{p}{]}\PY{p}{,}
         \PY{n}{marker}\PY{o}{=}\PY{l+s+s2}{\PYZdq{}}\PY{l+s+s2}{o}\PY{l+s+s2}{\PYZdq{}}\PY{p}{,}
         \PY{n}{label}\PY{o}{=}\PY{l+s+s2}{\PYZdq{}}\PY{l+s+s2}{Tangency unconstrained}\PY{l+s+s2}{\PYZdq{}}\PY{p}{)}

\PY{n}{plt}\PY{o}{.}\PY{n}{plot}\PY{p}{(}\PY{n}{sharpe\PYZus{}df}\PY{o}{.}\PY{n}{index}\PY{p}{,}
         \PY{n}{sharpe\PYZus{}df}\PY{p}{[}\PY{l+s+s2}{\PYZdq{}}\PY{l+s+s2}{NoShort}\PY{l+s+s2}{\PYZdq{}}\PY{p}{]}\PY{p}{,}
         \PY{n}{marker}\PY{o}{=}\PY{l+s+s2}{\PYZdq{}}\PY{l+s+s2}{o}\PY{l+s+s2}{\PYZdq{}}\PY{p}{,}
         \PY{n}{label}\PY{o}{=}\PY{l+s+s2}{\PYZdq{}}\PY{l+s+s2}{Tangency no\PYZhy{}short}\PY{l+s+s2}{\PYZdq{}}\PY{p}{)}

\PY{n}{plt}\PY{o}{.}\PY{n}{title}\PY{p}{(}\PY{l+s+s2}{\PYZdq{}}\PY{l+s+s2}{Sharpe ratios of tangency portfolios (5\PYZhy{}year rolling windows)}\PY{l+s+s2}{\PYZdq{}}\PY{p}{)}
\PY{n}{plt}\PY{o}{.}\PY{n}{xlabel}\PY{p}{(}\PY{l+s+s2}{\PYZdq{}}\PY{l+s+s2}{Window}\PY{l+s+s2}{\PYZdq{}}\PY{p}{)}
\PY{n}{plt}\PY{o}{.}\PY{n}{ylabel}\PY{p}{(}\PY{l+s+s2}{\PYZdq{}}\PY{l+s+s2}{Sharpe ratio}\PY{l+s+s2}{\PYZdq{}}\PY{p}{)}
\PY{n}{plt}\PY{o}{.}\PY{n}{xticks}\PY{p}{(}\PY{n}{rotation}\PY{o}{=}\PY{l+m+mi}{45}\PY{p}{)}
\PY{n}{plt}\PY{o}{.}\PY{n}{legend}\PY{p}{(}\PY{p}{)}
\PY{n}{plt}\PY{o}{.}\PY{n}{tight\PYZus{}layout}\PY{p}{(}\PY{p}{)}
\PY{n}{plt}\PY{o}{.}\PY{n}{show}\PY{p}{(}\PY{p}{)}
\end{Verbatim}
\end{tcolorbox}

    \begin{center}
    \adjustimage{max size={0.9\linewidth}{0.9\paperheight}}{Work-7_files/Work-7_115_0.png}
    \end{center}
    { \hspace*{\fill} \\}
    
    The figure tracks the Sharpe ratio of the tangency portfolio in each
5-year window, with and without short-selling.

Across all windows, the \textbf{unconstrained} portfolio delivers a
higher Sharpe ratio (around 1.5-2.1), while the \textbf{no-short}
version stays in the 0.7-1.2 range. Both specifications follow a similar
pattern over time: performance is strongest in the early windows, drops
sharply around 2014-2018, and then recovers.

So allowing short-selling does improve risk-adjusted performance, but
the gain is not explosive, especially compared to how extreme the
unconstrained weights are. In practice, the long-only portfolios achieve
reasonably high Sharpe ratios without requiring huge leverage.

    To get a sense of how \emph{stable} these portfolios are over time, we
compute the portfolio \textbf{turnover} between consecutive 5-year
windows.

For each strategy (unconstrained and no-short), the function
\texttt{turnover} measures \(\frac{1}{2}\sum_i |w_{i,t} - w_{i,t-1}|\)
between windows, so high values mean the optimizer is radically
reshuffling positions from one period to the next.

We then compare \texttt{turn\_uncon} and \texttt{turn\_noshort} to see
how much the no-short constraint helps to stabilise the weights.

    \begin{tcolorbox}[breakable, size=fbox, boxrule=1pt, pad at break*=1mm,colback=cellbackground, colframe=cellborder]
\begin{Verbatim}[commandchars=\\\{\}]
\PY{k}{def}\PY{+w}{ }\PY{n+nf}{turnover}\PY{p}{(}\PY{n}{weights\PYZus{}df}\PY{p}{)}\PY{p}{:}
    \PY{c+c1}{\PYZsh{} sum of |w\PYZus{}t \PYZhy{} w\PYZus{}\PYZob{}t\PYZhy{}1\PYZcb{}| / 2 for each window transition}
    \PY{n}{arr} \PY{o}{=} \PY{n}{weights\PYZus{}df}\PY{o}{.}\PY{n}{values}
    \PY{n}{to} \PY{o}{=} \PY{n}{np}\PY{o}{.}\PY{n}{abs}\PY{p}{(}\PY{n}{arr}\PY{p}{[}\PY{l+m+mi}{1}\PY{p}{:}\PY{p}{]} \PY{o}{\PYZhy{}} \PY{n}{arr}\PY{p}{[}\PY{p}{:}\PY{o}{\PYZhy{}}\PY{l+m+mi}{1}\PY{p}{]}\PY{p}{)}\PY{o}{.}\PY{n}{sum}\PY{p}{(}\PY{n}{axis}\PY{o}{=}\PY{l+m+mi}{1}\PY{p}{)} \PY{o}{/} \PY{l+m+mi}{2}
    \PY{k}{return} \PY{n}{pd}\PY{o}{.}\PY{n}{Series}\PY{p}{(}\PY{n}{to}\PY{p}{,} \PY{n}{index}\PY{o}{=}\PY{n}{weights\PYZus{}df}\PY{o}{.}\PY{n}{index}\PY{p}{[}\PY{l+m+mi}{1}\PY{p}{:}\PY{p}{]}\PY{p}{)}

\PY{n}{turn\PYZus{}uncon}  \PY{o}{=} \PY{n}{turnover}\PY{p}{(}\PY{n}{tan\PYZus{}uncon\PYZus{}weights\PYZus{}df}\PY{p}{)}
\PY{n}{turn\PYZus{}noshort} \PY{o}{=} \PY{n}{turnover}\PY{p}{(}\PY{n}{tan\PYZus{}noshort\PYZus{}weights\PYZus{}clean\PYZus{}df}\PY{p}{)}

\PY{n}{turn\PYZus{}uncon}\PY{p}{,} \PY{n}{turn\PYZus{}noshort}
\end{Verbatim}
\end{tcolorbox}

            \begin{tcolorbox}[breakable, size=fbox, boxrule=.5pt, pad at break*=1mm, opacityfill=0]
\begin{Verbatim}[commandchars=\\\{\}]
(2011-2015     4.808523
 2012-2016     5.073856
 2013-2017     0.994292
 2014-2018    12.750338
 2015-2019     6.660629
 2016-2020     2.885531
 2017-2021     4.366548
 dtype: float64,
 2011-2015    0.007778
 2012-2016    0.001493
 2013-2017    0.005967
 2014-2018    0.054306
 2015-2019    0.240966
 2016-2020    0.152214
 2017-2021    0.143141
 dtype: float64)
\end{Verbatim}
\end{tcolorbox}
        
    The turnover numbers confirm how unstable the unconstrained Markowitz
solution is.

For the \textbf{unconstrained tangency portfolio}, turnover is often
well above 1 (and even above 10 in 2014-2018), which means the portfolio
is completely reshuffled from one window to the next.

By contrast, the \textbf{no-short tangency portfolio} has very low
turnover: values stay close to zero for the early windows and only rise
moderately (around 0.15--0.24) in the last years. This shows that
imposing \(w_i \ge 0\) not only removes extreme long--short bets, but
also makes the allocation much more stable over time.

    \hypertarget{x-bonus-riskreturn-comparison-with-the-market-cap-portfolio}{%
\subsubsection{3.x Bonus: Risk--return comparison with the market-cap
portfolio}\label{x-bonus-riskreturn-comparison-with-the-market-cap-portfolio}}

As a final check, we compare the realised risk--return profile of three
portfolios over the last window (2017-2021): the market-cap portfolio,
the unconstrained tangency portfolio, and the no-short tangency
portfolio. This shows whether the more ``aggressive'' Markowitz
allocations actually deliver a better in-sample Sharpe ratio than a
simple market-cap benchmark.

    \begin{tcolorbox}[breakable, size=fbox, boxrule=1pt, pad at break*=1mm,colback=cellbackground, colframe=cellborder]
\begin{Verbatim}[commandchars=\\\{\}]
\PY{c+c1}{\PYZsh{} We rcompute annualised mean and covariance for the last window (2017–2021) }
\PY{n}{start\PYZus{}date} \PY{o}{=} \PY{l+s+s2}{\PYZdq{}}\PY{l+s+s2}{2017\PYZhy{}01\PYZhy{}01}\PY{l+s+s2}{\PYZdq{}}
\PY{n}{end\PYZus{}date}   \PY{o}{=} \PY{l+s+s2}{\PYZdq{}}\PY{l+s+s2}{2021\PYZhy{}12\PYZhy{}31}\PY{l+s+s2}{\PYZdq{}}
\PY{n}{window\PYZus{}1721} \PY{o}{=} \PY{n}{returns}\PY{o}{.}\PY{n}{loc}\PY{p}{[}\PY{n}{start\PYZus{}date}\PY{p}{:}\PY{n}{end\PYZus{}date}\PY{p}{]}

\PY{n}{mu\PYZus{}daily\PYZus{}1721}    \PY{o}{=} \PY{n}{window\PYZus{}1721}\PY{o}{.}\PY{n}{mean}\PY{p}{(}\PY{p}{)}\PY{o}{.}\PY{n}{values}
\PY{n}{Sigma\PYZus{}daily\PYZus{}1721} \PY{o}{=} \PY{n}{window\PYZus{}1721}\PY{o}{.}\PY{n}{cov}\PY{p}{(}\PY{p}{)}\PY{o}{.}\PY{n}{values}

\PY{n}{mu\PYZus{}1721}    \PY{o}{=} \PY{n}{mu\PYZus{}daily\PYZus{}1721} \PY{o}{*} \PY{n}{ann\PYZus{}factor}
\PY{n}{Sigma\PYZus{}1721} \PY{o}{=} \PY{n}{Sigma\PYZus{}daily\PYZus{}1721} \PY{o}{*} \PY{n}{ann\PYZus{}factor}

\PY{c+c1}{\PYZsh{} Build the 3 portfolios\PYZsq{} weights for that window}
\PY{c+c1}{\PYZsh{} mkt\PYZus{}weights is the Series we constructed from the \PYZsq{}market\PYZus{}cap\PYZsq{} sheet}
\PY{n}{w\PYZus{}mkt\PYZus{}last}     \PY{o}{=} \PY{n}{mkt\PYZus{}weights}\PY{o}{.}\PY{n}{values}                       \PY{c+c1}{\PYZsh{} market\PYZhy{}cap portfolio}
\PY{n}{w\PYZus{}uncon\PYZus{}last}   \PY{o}{=} \PY{n}{tan\PYZus{}uncon\PYZus{}weights\PYZus{}df}\PY{o}{.}\PY{n}{loc}\PY{p}{[}\PY{l+s+s2}{\PYZdq{}}\PY{l+s+s2}{2017\PYZhy{}2021}\PY{l+s+s2}{\PYZdq{}}\PY{p}{]}\PY{o}{.}\PY{n}{values}      \PY{c+c1}{\PYZsh{} unconstrained tangency}
\PY{n}{w\PYZus{}noshort\PYZus{}last} \PY{o}{=} \PY{n}{tan\PYZus{}noshort\PYZus{}weights\PYZus{}clean\PYZus{}df}\PY{o}{.}\PY{n}{loc}\PY{p}{[}\PY{l+s+s2}{\PYZdq{}}\PY{l+s+s2}{2017\PYZhy{}2021}\PY{l+s+s2}{\PYZdq{}}\PY{p}{]}\PY{o}{.}\PY{n}{values}  \PY{c+c1}{\PYZsh{} no\PYZhy{}short tangency}

\PY{c+c1}{\PYZsh{} Compute (return, vol, Sharpe) for each}
\PY{n}{ret\PYZus{}mkt}\PY{p}{,} \PY{n}{vol\PYZus{}mkt}\PY{p}{,} \PY{n}{sh\PYZus{}mkt} \PY{o}{=} \PY{n}{portfolio\PYZus{}stats}\PY{p}{(}\PY{n}{w\PYZus{}mkt\PYZus{}last}\PY{p}{,}     \PY{n}{mu\PYZus{}1721}\PY{p}{,} \PY{n}{Sigma\PYZus{}1721}\PY{p}{,} \PY{n}{rf}\PY{o}{=}\PY{n}{rf}\PY{p}{)}
\PY{n}{ret\PYZus{}uc}\PY{p}{,}  \PY{n}{vol\PYZus{}uc}\PY{p}{,}  \PY{n}{sh\PYZus{}uc}  \PY{o}{=} \PY{n}{portfolio\PYZus{}stats}\PY{p}{(}\PY{n}{w\PYZus{}uncon\PYZus{}last}\PY{p}{,}   \PY{n}{mu\PYZus{}1721}\PY{p}{,} \PY{n}{Sigma\PYZus{}1721}\PY{p}{,} \PY{n}{rf}\PY{o}{=}\PY{n}{rf}\PY{p}{)}
\PY{n}{ret\PYZus{}ns}\PY{p}{,}  \PY{n}{vol\PYZus{}ns}\PY{p}{,}  \PY{n}{sh\PYZus{}ns}  \PY{o}{=} \PY{n}{portfolio\PYZus{}stats}\PY{p}{(}\PY{n}{w\PYZus{}noshort\PYZus{}last}\PY{p}{,} \PY{n}{mu\PYZus{}1721}\PY{p}{,} \PY{n}{Sigma\PYZus{}1721}\PY{p}{,} \PY{n}{rf}\PY{o}{=}\PY{n}{rf}\PY{p}{)}

\PY{n}{sharpe\PYZus{}compare} \PY{o}{=} \PY{n}{pd}\PY{o}{.}\PY{n}{DataFrame}\PY{p}{(}
    \PY{p}{\PYZob{}}
        \PY{l+s+s2}{\PYZdq{}}\PY{l+s+s2}{ann\PYZus{}return}\PY{l+s+s2}{\PYZdq{}}\PY{p}{:} \PY{p}{[}\PY{n}{ret\PYZus{}mkt}\PY{p}{,} \PY{n}{ret\PYZus{}uc}\PY{p}{,} \PY{n}{ret\PYZus{}ns}\PY{p}{]}\PY{p}{,}
        \PY{l+s+s2}{\PYZdq{}}\PY{l+s+s2}{ann\PYZus{}vol}\PY{l+s+s2}{\PYZdq{}}\PY{p}{:}    \PY{p}{[}\PY{n}{vol\PYZus{}mkt}\PY{p}{,} \PY{n}{vol\PYZus{}uc}\PY{p}{,} \PY{n}{vol\PYZus{}ns}\PY{p}{]}\PY{p}{,}
        \PY{l+s+s2}{\PYZdq{}}\PY{l+s+s2}{sharpe}\PY{l+s+s2}{\PYZdq{}}\PY{p}{:}     \PY{p}{[}\PY{n}{sh\PYZus{}mkt}\PY{p}{,} \PY{n}{sh\PYZus{}uc}\PY{p}{,} \PY{n}{sh\PYZus{}ns}\PY{p}{]}\PY{p}{,}
    \PY{p}{\PYZcb{}}\PY{p}{,}
    \PY{n}{index}\PY{o}{=}\PY{p}{[}\PY{l+s+s2}{\PYZdq{}}\PY{l+s+s2}{MarketCap}\PY{l+s+s2}{\PYZdq{}}\PY{p}{,} \PY{l+s+s2}{\PYZdq{}}\PY{l+s+s2}{MV\PYZus{}Tangency\PYZus{}Uncon}\PY{l+s+s2}{\PYZdq{}}\PY{p}{,} \PY{l+s+s2}{\PYZdq{}}\PY{l+s+s2}{MV\PYZus{}Tangency\PYZus{}NoShort}\PY{l+s+s2}{\PYZdq{}}\PY{p}{]}\PY{p}{,}
\PY{p}{)}

\PY{n}{sharpe\PYZus{}compare}
\end{Verbatim}
\end{tcolorbox}

            \begin{tcolorbox}[breakable, size=fbox, boxrule=.5pt, pad at break*=1mm, opacityfill=0]
\begin{Verbatim}[commandchars=\\\{\}]
                     ann\_return   ann\_vol    sharpe
MarketCap              0.134005  0.175058  0.765489
MV\_Tangency\_Uncon      0.491154  0.319996  1.534876
MV\_Tangency\_NoShort    0.148207  0.165613  0.894899
\end{Verbatim}
\end{tcolorbox}
        
    The table below compares the performance of three portfolios over
2017-2021:

\begin{itemize}
\item
  \textbf{MarketCap}: annual return ≈ \(13.4\%\), volatility ≈
  \(17.5\%\), Sharpe ≈ \(0.77\).
\item
  \textbf{Unconstrained tangency (MV\_Tangency\_Uncon)}: very high
  annual return (≈ \(49.1\%\)) but also much higher risk (volatility ≈
  \(32.0\%\)), with a Sharpe ratio around \(1.53\).
\item
  \textbf{No-short tangency (MV\_Tangency\_NoShort)}: annual return ≈
  \(14.8\%\), volatility ≈ \(16.6\%\), Sharpe ≈ \(0.89\).
\end{itemize}

This confirms the pattern seen in the weights:

\begin{itemize}
\tightlist
\item
  The \textbf{unconstrained Markowitz portfolio} offers the highest
  Sharpe ratio, but only by taking on a lot more risk and extreme
  long--short positions.
\item
  The \textbf{no-short tangency portfolio} slightly improves on the
  market-cap Sharpe ratio and keeps risk at a similar level, with a much
  more realistic, long-only allocation.
\item
  The \textbf{market-cap portfolio} sits in between: lower efficiency
  than the theoretical tangency portfolio, but very reasonable risk and
  no leverage or short-selling.
\end{itemize}

    \hypertarget{conclusion}{%
\subsection{Conclusion}\label{conclusion}}

In Question 3, the rolling 5-year windows make the message very clear:
the \textbf{unconstrained} tangency portfolio is extremely unstable.
Small changes in estimated returns and covariances lead to big shifts in
weights and high turnover from one window to the next. With the
\textbf{no short-sale constraint}, the portfolios are much more stable
over time and look much closer to something an actual investor could
hold.

Taken together with Questions 1 and 2, the whole assignment shows the
same story:\\
- The data are not perfectly Gaussian or i.i.d.,\\
- Unconstrained Markowitz can produce very high theoretical Sharpe
ratios, but with crazy, levered, long--short positions,\\
- Adding realistic constraints (like \(w_i \ge 0\)) and comparing to a
simple market-cap portfolio brings us back to something much more
sensible.

So in the end, mean--variance optimisation is a useful theoretical tool,
but in practice it really needs to be combined with constraints and
simple benchmarks if we want portfolios that are both \emph{efficient on
paper} and \emph{investable in reality}.


    % Add a bibliography block to the postdoc
    
    
    
\end{document}
